% Options for packages loaded elsewhere
\PassOptionsToPackage{unicode}{hyperref}
\PassOptionsToPackage{hyphens}{url}
\PassOptionsToPackage{dvipsnames,svgnames,x11names}{xcolor}
%
\documentclass[
  10pt,
]{article}
\usepackage{amsmath,amssymb}
\usepackage{iftex}
\ifPDFTeX
  \usepackage[T1]{fontenc}
  \usepackage[utf8]{inputenc}
  \usepackage{textcomp} % provide euro and other symbols
\else % if luatex or xetex
  \usepackage{unicode-math} % this also loads fontspec
  \defaultfontfeatures{Scale=MatchLowercase}
  \defaultfontfeatures[\rmfamily]{Ligatures=TeX,Scale=1}
\fi
\ifPDFTeX\else
  % xetex/luatex font selection
\fi
% Use upquote if available, for straight quotes in verbatim environments
\IfFileExists{upquote.sty}{\usepackage{upquote}}{}
\IfFileExists{microtype.sty}{% use microtype if available
  \usepackage[]{microtype}
  \UseMicrotypeSet[protrusion]{basicmath} % disable protrusion for tt fonts
}{}
\makeatletter
\@ifundefined{KOMAClassName}{% if non-KOMA class
  \IfFileExists{parskip.sty}{%
    \usepackage{parskip}
  }{% else
    \setlength{\parindent}{0pt}
    \setlength{\parskip}{6pt plus 2pt minus 1pt}}
}{% if KOMA class
  \KOMAoptions{parskip=half}}
\makeatother
\usepackage{xcolor}
\usepackage{color}
\usepackage{fancyvrb}
\newcommand{\VerbBar}{|}
\newcommand{\VERB}{\Verb[commandchars=\\\{\}]}
\DefineVerbatimEnvironment{Highlighting}{Verbatim}{commandchars=\\\{\}}
% Add ',fontsize=\small' for more characters per line
\newenvironment{Shaded}{}{}
\newcommand{\AlertTok}[1]{\textcolor[rgb]{1.00,0.00,0.00}{\textbf{#1}}}
\newcommand{\AnnotationTok}[1]{\textcolor[rgb]{0.38,0.63,0.69}{\textbf{\textit{#1}}}}
\newcommand{\AttributeTok}[1]{\textcolor[rgb]{0.49,0.56,0.16}{#1}}
\newcommand{\BaseNTok}[1]{\textcolor[rgb]{0.25,0.63,0.44}{#1}}
\newcommand{\BuiltInTok}[1]{\textcolor[rgb]{0.00,0.50,0.00}{#1}}
\newcommand{\CharTok}[1]{\textcolor[rgb]{0.25,0.44,0.63}{#1}}
\newcommand{\CommentTok}[1]{\textcolor[rgb]{0.38,0.63,0.69}{\textit{#1}}}
\newcommand{\CommentVarTok}[1]{\textcolor[rgb]{0.38,0.63,0.69}{\textbf{\textit{#1}}}}
\newcommand{\ConstantTok}[1]{\textcolor[rgb]{0.53,0.00,0.00}{#1}}
\newcommand{\ControlFlowTok}[1]{\textcolor[rgb]{0.00,0.44,0.13}{\textbf{#1}}}
\newcommand{\DataTypeTok}[1]{\textcolor[rgb]{0.56,0.13,0.00}{#1}}
\newcommand{\DecValTok}[1]{\textcolor[rgb]{0.25,0.63,0.44}{#1}}
\newcommand{\DocumentationTok}[1]{\textcolor[rgb]{0.73,0.13,0.13}{\textit{#1}}}
\newcommand{\ErrorTok}[1]{\textcolor[rgb]{1.00,0.00,0.00}{\textbf{#1}}}
\newcommand{\ExtensionTok}[1]{#1}
\newcommand{\FloatTok}[1]{\textcolor[rgb]{0.25,0.63,0.44}{#1}}
\newcommand{\FunctionTok}[1]{\textcolor[rgb]{0.02,0.16,0.49}{#1}}
\newcommand{\ImportTok}[1]{\textcolor[rgb]{0.00,0.50,0.00}{\textbf{#1}}}
\newcommand{\InformationTok}[1]{\textcolor[rgb]{0.38,0.63,0.69}{\textbf{\textit{#1}}}}
\newcommand{\KeywordTok}[1]{\textcolor[rgb]{0.00,0.44,0.13}{\textbf{#1}}}
\newcommand{\NormalTok}[1]{#1}
\newcommand{\OperatorTok}[1]{\textcolor[rgb]{0.40,0.40,0.40}{#1}}
\newcommand{\OtherTok}[1]{\textcolor[rgb]{0.00,0.44,0.13}{#1}}
\newcommand{\PreprocessorTok}[1]{\textcolor[rgb]{0.74,0.48,0.00}{#1}}
\newcommand{\RegionMarkerTok}[1]{#1}
\newcommand{\SpecialCharTok}[1]{\textcolor[rgb]{0.25,0.44,0.63}{#1}}
\newcommand{\SpecialStringTok}[1]{\textcolor[rgb]{0.73,0.40,0.53}{#1}}
\newcommand{\StringTok}[1]{\textcolor[rgb]{0.25,0.44,0.63}{#1}}
\newcommand{\VariableTok}[1]{\textcolor[rgb]{0.10,0.09,0.49}{#1}}
\newcommand{\VerbatimStringTok}[1]{\textcolor[rgb]{0.25,0.44,0.63}{#1}}
\newcommand{\WarningTok}[1]{\textcolor[rgb]{0.38,0.63,0.69}{\textbf{\textit{#1}}}}
\usepackage{longtable,booktabs,array}
\usepackage{calc} % for calculating minipage widths
% Correct order of tables after \paragraph or \subparagraph
\usepackage{etoolbox}
\makeatletter
\patchcmd\longtable{\par}{\if@noskipsec\mbox{}\fi\par}{}{}
\makeatother
% Allow footnotes in longtable head/foot
\IfFileExists{footnotehyper.sty}{\usepackage{footnotehyper}}{\usepackage{footnote}}
\makesavenoteenv{longtable}
\usepackage{graphicx}
\makeatletter
\def\maxwidth{\ifdim\Gin@nat@width>\linewidth\linewidth\else\Gin@nat@width\fi}
\def\maxheight{\ifdim\Gin@nat@height>\textheight\textheight\else\Gin@nat@height\fi}
\makeatother
% Scale images if necessary, so that they will not overflow the page
% margins by default, and it is still possible to overwrite the defaults
% using explicit options in \includegraphics[width, height, ...]{}
\setkeys{Gin}{width=\maxwidth,height=\maxheight,keepaspectratio}
% Set default figure placement to htbp
\makeatletter
\def\fps@figure{htbp}
\makeatother
\setlength{\emergencystretch}{3em} % prevent overfull lines
\providecommand{\tightlist}{%
  \setlength{\itemsep}{0pt}\setlength{\parskip}{0pt}}
\setcounter{secnumdepth}{-\maxdimen} % remove section numbering
\ifLuaTeX
\usepackage[bidi=basic]{babel}
\else
\usepackage[bidi=default]{babel}
\fi
\babelprovide[main,import]{american}
% get rid of language-specific shorthands (see #6817):
\let\LanguageShortHands\languageshorthands
\def\languageshorthands#1{}

\usepackage{fontspec}
\setmainfont{TeX Gyre Termes}[Scale=1.05]
\setsansfont{TeX Gyre Heros}
\setmonofont{DejaVu Sans Mono}[Scale=0.82]
\usepackage[letterpaper,margin=0.85in,headheight=14pt,headsep=0.2in,footskip=0.38in]{geometry}
\usepackage{booktabs}
\usepackage{colortbl}
\definecolor{cvblue}{HTML}{2E74B5}
\definecolor{cvdark}{HTML}{1F4D78}
\definecolor{cvdeep}{HTML}{1F4E79}
\definecolor{cvbox}{HTML}{8EA9C1}
\definecolor{cvfill}{HTML}{EEF3F8}
\definecolor{cvgrid}{HTML}{A6A6A6}
\definecolor{cvhead}{HTML}{DCE6F1}
\definecolor{cvgrey}{HTML}{595959}
\usepackage{caption}
\captionsetup{font=small,labelformat=empty,skip=4pt,justification=justified,singlelinecheck=false}
\usepackage{microtype}
\usepackage{xurl}
\usepackage{titlesec}
\titleformat{\section}{\fontsize{14}{17}\selectfont\bfseries\color{cvblue}}{}{0pt}{}
\titleformat{\subsection}{\fontsize{12.5}{15}\selectfont\bfseries\color{cvblue}}{}{0pt}{}
\titleformat{\subsubsection}{\normalsize\bfseries\color{cvdark}}{}{0pt}{}
\titlespacing*{\section}{0pt}{12pt plus 2pt minus 2pt}{5pt}
\titlespacing*{\subsection}{0pt}{9pt plus 2pt minus 2pt}{4pt}
\titlespacing*{\subsubsection}{0pt}{8pt plus 2pt minus 2pt}{4pt}
\usepackage{float}
\linespread{1.03}
\setlength{\parindent}{0pt}
\setlength{\parskip}{4pt plus 1pt}
\setlength{\emergencystretch}{3em}
\setlength{\textfloatsep}{12pt plus 2pt minus 2pt}
\setlength{\intextsep}{12pt plus 2pt minus 2pt}
\renewcommand{\arraystretch}{1.12}
\arrayrulecolor{cvgrid}
\usepackage{etoolbox}
\def\tablefont{\fontsize{8.5}{10.4}\selectfont}
\AtBeginEnvironment{longtable}{\tablefont}
\usepackage{fancyhdr}
\usepackage{lastpage}
\fancypagestyle{cave}{\fancyhf{}%
  \fancyhead[R]{\fontsize{9}{11}\selectfont\itshape\color{cvgrey}Did the Action Already Happen? Reliable Recovery for AI Agents Under Ambiguous Tool Outcomes}%
  \fancyfoot[C]{\fontsize{9}{11}\selectfont\color{cvgrey}Page \thepage\ of \pageref*{LastPage}}%
  \renewcommand{\headrule}{{\color{cvbox}\hrule height 0.4pt}}\renewcommand{\footrulewidth}{0pt}}
\pagestyle{cave}
\AtBeginDocument{\thispagestyle{cave}}
\fancypagestyle{plain}{\pagestyle{cave}}
\usepackage[most]{tcolorbox}
\renewenvironment{abstract}{\begin{tcolorbox}[enhanced,colback=cvfill,colframe=cvbox,boxrule=0.5pt,arc=0pt,
  left=6pt,right=6pt,top=5pt,bottom=5pt,fontupper=\fontsize{10.5}{12}\selectfont]%
  {\bfseries\color{cvdeep}ABSTRACT}\par\vspace{2pt}}{\end{tcolorbox}}
\makeatletter
\@ifundefined{Highlighting}{}{\RecustomVerbatimEnvironment{Highlighting}{Verbatim}{commandchars=\\\{\},fontsize=\small}}
\renewcommand{\@maketitle}{\newpage\null\vspace{-0.6em}%
  \begin{center}{\fontsize{16}{20}\selectfont\bfseries\@title\par}\vskip 12pt
  {\@author\par}\vskip 2pt{\fontsize{11}{14}\selectfont\@date\par}\end{center}%
  \vskip 2pt{\color{cvdeep}\hrule height 0.8pt}\vskip 10pt}
\makeatother
\ifLuaTeX
  \usepackage{selnolig}  % disable illegal ligatures
\fi
\IfFileExists{bookmark.sty}{\usepackage{bookmark}}{\usepackage{hyperref}}
\IfFileExists{xurl.sty}{\usepackage{xurl}}{} % add URL line breaks if available
\urlstyle{same}
\hypersetup{
  pdftitle={Did the Action Already Happen? Reliable Recovery for AI Agents Under Ambiguous Tool Outcomes},
  pdfauthor={Ankitha Vuyyala; Independent Researcher, Louisville, Kentucky, USA; ankithavuyyala08@gmail.com; ORCID: https://orcid.org/0009-0004-1916-8070},
  pdflang={en-US},
  colorlinks=true,
  linkcolor={Maroon},
  filecolor={Maroon},
  citecolor={Blue},
  urlcolor={Blue},
  pdfcreator={LaTeX via pandoc}}

\title{Did the Action Already Happen? Reliable Recovery for AI Agents
Under Ambiguous Tool Outcomes}
\author{{\fontsize{13}{16}\selectfont\bfseries Ankitha Vuyyala}\\[3pt]{\fontsize{11}{14}\selectfont Independent Researcher, Louisville, Kentucky, USA}\\{\fontsize{11}{14}\selectfont ankithavuyyala08@gmail.com}\\{\fontsize{11}{14}\selectfont ORCID: https://orcid.org/0009-0004-1916-8070}}
\date{October 2026}

\begin{document}
\maketitle
\begin{abstract}
When a tool call made by an AI agent times out or loses its reply, the
agent cannot tell whether the effect happened, did not happen, or is
still on its way. The standard remedy is an idempotency key, but in a
documentation census of 52 widely used write APIs, 38 offer no
server-side key on the operation an agent would call, and several keys
that exist expire within hours. I argue that recovery needs
\emph{attributability}, the ability to count the standing effects of
one's own operation, and derive for each API contract and read model
whether a client can achieve prevention of duplicates, convergence of
the final state, or only honest abstention; any policy that must stop at
a horizon has to bet between two indistinguishable worlds. Of the 38
key-less APIs, 31 persist a client-controlled field that can be looked
up. I build Settle, a recovery layer that journals each intent, stamps
an operation id into such a field, resolves unknown outcomes by that id
up to a horizon, and compensates only effects it can prove are surplus.
Under internally frozen protocols I evaluate it on six unmodified
services in containers, with injected network faults and process
crashes, against ten baselines including idempotency keys, a Temporal
workflow engine and a client-side guard, with a scripted agent and four
language models. Across the attributable and reversible grid, Settle
reduces inconsistent final states from 34.7\% for the strongest
matched-latency baseline of the main study to 3.2\%; without the two
fault modes that structurally require a final sweep or row-level checks,
the rates are 29.4\% and 3.4\%. Durable keys are exact only where the
service honours them. A follow-up study of 4,272 episodes bounds the
claims: a tag-based wait-then-count policy matches Settle on every
ordinary fault (it is inconsistent in 10.8\% of these attributable and
reversible worlds against Settle's 3.2\%, with 42 of its 65 failures on
transport redelivery), so convergence pays off mainly for duplicates the
transport creates; an author filter combined with waiting rescues
content checks against look-alikes from other accounts; and a frontier
model without the layer fails in the same cells as smaller ones (52.2\%
of hard worlds, against 6.0\% with Settle).
\end{abstract}

\textbf{Keywords:} LLM agents; tool use; exactly-once; idempotency;
fault tolerance; compensation; durable execution; API design

\hypertarget{introduction}{%
\section*{1. Introduction}\label{introduction}}
\addcontentsline{toc}{section}{1. Introduction}

An agent that acts on the world through tools inherits the oldest
problem of distributed computing. It sent a request and no answer came
back. The request may have been lost on the way; it may have been
carried out with only the reply lost; or it may still be in flight and
land a few seconds after the agent has given up. The three cases look
identical to the agent (Fig. 1). If the agent retries, it creates a
second issue, a second e-mail or a second payment whenever the first
attempt went through. If it does not, the task stays incomplete whenever
the first attempt did not.

\begin{figure}
\centering
\includegraphics[width=0.68\textwidth,height=\textheight]{fig_worlds.png}
\caption{\textbf{Fig. 1.} Three worlds that give the agent the same
observation. A request whose reply was lost, one that never arrived, and
one that will commit after the agent has given up all end in the same
timeout. Settle decides only after the in-flight bound \(\Delta\) and
the read lag \(\lambda\) have passed, by counting the effects that carry
the operation's tag.}
\end{figure}

The problem is as old as the two-generals argument {[}1{]}, {[}2{]}, and
it has a standard engineering answer: give each logical operation a key
and let the server execute each key at most once {[}3{]}--{[}5{]}.
Recent work shows that language-model agents run into exactly this
problem and that the standard answer does not carry over by itself.
Models frequently re-issue a write after an ambiguous failure without
first checking whether it took effect {[}6{]}, {[}7{]}. Li {[}8{]}
proves that no policy which only observes outcomes can be exactly-once
while a request may still be in flight, and finds that idempotency keys
remove most duplicates, but only on tools that honour them. Trofimov and
Novikov {[}9{]} argue that without a primitive that yields one
authoritative outcome, no recovery protocol is both safe under unknown
outcomes and safe to compensate. Agent runtimes that add transactions
{[}10{]}, {[}11{]}, durable execution {[}12{]}, {[}13{]} or rollback
{[}14{]}, {[}15{]} reach exactly-once behaviour only when the downstream
service honours a key or when failures are detectable, and neither holds
for a lost reply.

How large is the gap that this leaves? I examined the documentation of
52 widely used write APIs across payments, developer tools, messaging
and data services (Section 5). Only 14 of them offer a server-side key
on the operation an agent would call, and most of those are payment
APIs. Of these 14 keys, 4 are documented to expire within six hours or
when the object's state changes. None of the twelve issue trackers and
project tools and none of the thirteen messaging and mail services I
examined offers one. These are services where many agent writes land and
where a duplicate is visible to people. For these services, recovery
without a key is not a corner case but the common case.

\textbf{The key observation.} Preventing a duplicate requires a key, but
\emph{repairing} one requires only that the client can tell which
effects belong to which of its operations. Many key-less APIs persist a
field that the client controls: an issue label or body, a mail header, a
metadata map, a reference column, a custom property. If the client
writes an operation id into that field, it can later count the effects
of its own operation exactly. A search by content cannot do that,
because other users may create records with the same content. With an
exact count, the client can wait out the in-flight window, re-issue only
if no effect exists, and remove any surplus copy it can prove is its
own. In the census, 31 of the 38 key-less APIs offer such a field
together with a way to look it up, and 21 of them offer an exact lookup.
Payment engineers know this ``client reference plus authoritative
lookup'' pattern informally. To my knowledge it has not been formalised
for agents, compared against keys and workflow engines, or tested in the
places where it must fail.

This paper makes four contributions.

\begin{enumerate}
\def\labelenumi{\arabic{enumi}.}
\tightlist
\item
  \textbf{A theory of what recovery can achieve} (Section 3, Section 4).
  I define \emph{attributability} as the identifiability of the number
  of standing effects of an operation from what the client can observe.
  The notion unifies native keys, client tags, writer exclusivity and
  returned ids. I prove which guarantee each contract and read model
  admits. Without a known in-flight bound, prevention requires a durable
  key, and a key protects only within its retention. Convergence of the
  final state is achievable for attributable, reversible effects without
  any key; for irreversible effects it collapses to prevention.
  Content-based verification is unsound once look-alikes exist,
  aggregates are in general attributable only under exclusive writing,
  and the tag must survive a crash between dispatch and record. Finally,
  a policy that must stop at a horizon faces two indistinguishable
  worlds, so its residual risk is bounded below by the less likely of
  them. Settle's rule attains that bound when, given what it has seen, a
  late landing is less likely than a lost request. The result is an
  achievability map that predicts, cell by cell, where each recovery
  method must fail.
\item
  \textbf{A census of real API contracts} (Section 5). I code 52 APIs
  from their official documentation and re-check a random sample in a
  second pass. The census places real services in the achievability map
  and measures the premise of the paper rather than asserting it.
\item
  \textbf{Settle, a recovery layer} that implements the constructive
  side of the theory (Section 6). It writes each intent to a durable
  journal before dispatch, stamps an operation id into a persisted
  field, resolves unknown outcomes by that id after the read path's lag,
  and keeps reconciling up to a horizon. It compensates only effects it
  can prove are surplus, including effects the agent itself asks to
  delete, and reports \emph{unknown} or an unrecoverable duplicate where
  nothing can be guaranteed.
\item
  \textbf{An internally pre-registered evaluation on real services with
  real faults} (Section 7, Section 8). Six services run in containers:
  PostgreSQL, Gitea, Mailpit, Redis, SeaweedFS and TigerBeetle. Faults
  are injected at the network and process level and validated against
  each service's own records: lost replies, requests that never arrive,
  held connections, late commits with fixed and heavy-tailed delays,
  partial batches, transport redelivery and agent-process crashes. I
  compare sixteen conditions, namely ten baselines, Settle, four
  ablations of Settle and an outcome oracle, under a scripted agent and
  three OpenAI models. Grading uses service-side records only, and
  hypotheses, analysis code and source hashes were frozen internally
  before the first test episode. A second, separately frozen follow-up
  study of 4,272 episodes adds the strongest tag-based alternative I
  could build, Settle without its implementation gaps, delayed
  redelivery, look-alikes from another account, plan-free intent binding
  and a frontier model (Section 9).
\end{enumerate}

The evidence locates the boundary that the theory predicts. The scripted
grid replays 1,188 worlds under fifteen conditions each (17,820
episodes), and the language-model grid adds 2,680 episodes. In ambiguous
worlds whose write is attributable and reversible, Settle leaves 3.2\%
of episodes inconsistent, against 34.7\% for the strongest baseline, a
guard that waits exactly as long as Settle does. Without the two fault
modes that every baseline fails by construction, redelivery and partial
batches, the rates are 3.4\% and 29.4\%. Every remaining Settle failure
is accounted for by the theory: either Settle's assumed in-flight bound
was violated, whose consequence the theory bounds and which trades
against latency, or the write lies in a cell where the theory admits no
client-side guarantee. Durable keys, here those of a workflow engine,
are exact where the service honours the key, but on key-less writes they
are no better than plain retries (61.0\% for the workflow engine against
61.0\% for plain retries). Verification by content, the basis of current
client-side guards, fails in more than half of the worlds that contain
same-account look-alike records. Language-model agents either re-issue
blindly or check first. Checking helps, but it fails under look-alikes,
after crash replay and under long delays: with my worst-case look-alikes
and under long delays the observation itself is insufficient, and after
a crash the framework re-issues before the model can check. With Settle
underneath, all three models fall to 6.0\% inconsistency in these hard
worlds. Where nothing can be guaranteed, as with shared aggregates and
redelivered e-mail, every method short of an outcome oracle fails, and
Settle reports these cases instead of claiming success. The follow-up
study draws the boundaries more tightly, including where they work
against Settle. A tag-based wait-then-count policy is exactly as good as
Settle on ordinary faults, so Settle's convergence pays off mainly for
duplicates created by the transport. A content check with an author
filter, combined with waiting out the bound, is as good as a tag when
look-alikes come from another account. And the frontier model
\texttt{gpt-5.6-sol} without a recovery layer leaves 52.2\% of hard
worlds inconsistent, failing in the same cells as the smaller models,
against 6.0\% with Settle.

Section 2 positions the work. Section 3 and Section 4 give the model and
the propositions, Section 5 presents the census, and Section 6 describes
Settle. Section 7 and Section 8 present the experimental design and the
results, and Section 9 the internally pre-registered follow-up study.
Section 10 draws the consequences for API and agent-framework design and
states the limitations.

\hypertarget{background-and-related-work}{%
\section*{2. Background and related
work}\label{background-and-related-work}}
\addcontentsline{toc}{section}{2. Background and related work}

\textbf{Exactly-once effects in distributed systems.} Two parties cannot
agree on whether a message was acted upon when messages can be lost
{[}1{]}, {[}2{]}. Practical systems therefore build
\emph{effectively-once} behaviour from at-least-once delivery plus
deduplication. The ingredients are idempotent operations and request
identifiers {[}3{]}, durable completion records that let a server
recognise a retried request {[}4{]}, transactional logs in stream
processors {[}16{]}, and exactly-once semantics for stateful serverless
workflows {[}17{]}. Commercial APIs expose the same idea as an
idempotency key {[}5{]}, {[}18{]}. Durable-execution engines persist
workflow progress but execute each activity at least once, and advise
that activities be idempotent {[}12{]}, {[}13{]}. When an operation
cannot be undone atomically, sagas and compensating transactions restore
consistency by running a semantic inverse {[}19{]}, {[}20{]}. All of
these mechanisms assume either that the server deduplicates or that the
client knows which effects exist. My setting removes both assumptions
and asks what remains. The answer, attribution through a reference the
client writes, is close to what payment integrations do with merchant
references. My contribution is to state when it suffices and when it
cannot, and to measure how it compares with keys and durable execution
when an autonomous agent drives the calls.

\textbf{Ambiguous tool outcomes in language-model agents.} Sun {[}6{]}
builds observation-equivalent pairs of worlds, in which a write either
did or did not execute, and shows that an agent without further
information can be right in at most half of them. IdempotencyBench
{[}7{]} measures idempotency violations under retries with a hidden
effect ledger. Li {[}8{]} gives the most complete study to date: twelve
fault modes on simulated services, nine models and three production
harnesses. That study proves that no verification-only policy is
exactly-once under late commits without an in-flight bound. It also
finds that idempotency keys remove most duplicates, and that a
client-side guard helps weaker models more than frontier models.
Trofimov and Novikov {[}9{]} model effect histories and name the
anomalies that arise at the agent-tool boundary: duplicated, missing,
orphaned compensation and residue. They argue that without an
authoritative-outcome primitive no protocol is both safe under unknown
outcomes and safe to compensate. Their census of tool annotations in the
Model Context Protocol shows that such annotations cannot express these
properties, and the specification itself calls its idempotency and
destructiveness annotations untrusted hints {[}21{]}. Wang {[}22{]}
shows that interfaces which expose effect semantics reduce duplicate
effects.

My work accepts Li's impossibility result for \emph{prevention} and
studies \emph{convergence} instead. It also differs from this line of
work in four respects:

\begin{itemize}
\tightlist
\item
  I \emph{construct} the missing outcome primitive from fields that
  ordinary APIs already persist, and measure in a census how often such
  fields exist; Trofimov and Novikov's census covers tool annotations,
  not the services behind the tools.
\item
  I evaluate on real services rather than simulated ones.
\item
  I add agent-process crashes, a workflow engine, compensation-based
  baselines, look-alike records from concurrent users and aggregate-only
  reads.
\item
  I prove a lower bound for policies that must stop at a horizon, which
  characterises the residual failures I measure.
\end{itemize}

\textbf{Recovery mechanisms for agents.}

\begin{itemize}
\tightlist
\item
  \emph{Verification.} Verify-before-retry wrappers check a
  postcondition before repeating a call {[}23{]}. Li's guard adds
  automatic keys, a re-check after the documented lag and a block on
  repeats it cannot verify {[}8{]}.
\item
  \emph{Transactions.} Atomix gates commits on progress, compensates
  failed branches and deduplicates ambiguous retries with keys {[}10{]}.
  Cordon offers semantic transactions with an effect outbox {[}11{]}.
  Agentic Transaction proposes ACID-style guarantees for data agents
  {[}24{]}.
\item
  \emph{Compensation.} Robust Agent Compensation and SagaLLM bring
  log-based compensation and sagas to agents {[}14{]}, {[}15{]}.
\item
  \emph{Checkpoint and restore.} AgentRewind {[}25{]} and DART {[}26{]}
  restore agent state. Rolling back agent state does not undo effects in
  the world, and a restored agent may re-synthesise slightly different
  requests that a server treats as new {[}27{]}. CapLease measures such
  equivalent re-proposals and binds them to a durable authorisation
  state {[}28{]}.
\end{itemize}

Except for Li's guard and verify-before-retry, these mechanisms assume
keyed downstream services or detectable failures. I include baselines
from the verification family (verify-before-retry and Li's guard), the
compensation family (log-based rollback) and durable execution (a
Temporal workflow engine), together with idempotency keys, all with
matched retry budgets; transactional runtimes and checkpoint-and-restore
are not represented. I also add a variant of the guard that waits as
long as Settle, so that waiting alone cannot explain a difference.

\textbf{Fault injection and reliability benchmarks for agents.} Several
benchmarks stress agents with tool failures, timeouts, rate limits or
corrupted outputs {[}29{]}--{[}32{]}. Related work studies guards under
state races {[}33{]}, revisable streaming execution {[}34{]}, and agents
that claim success when the state disagrees {[}35{]}. These works
measure task success under faults. I measure the correctness of the
\emph{final external state} after an ambiguous write, which task success
does not capture, and I also score whether the agent's report is honest
about that state.

\hypertarget{problem-setting}{%
\section*{3. Problem setting}\label{problem-setting}}
\addcontentsline{toc}{section}{3. Problem setting}

\textbf{Operations, attempts and effects.} A \emph{logical operation}
\(\ell\) is one intended change to an external service: create one
issue, send one e-mail, make one transfer. The client makes
\emph{attempts} \(a_1, a_2, \dots\) of \(\ell\), dispatched at times
\(t_1 < t_2 < \dots\). An attempt either \emph{executes} at an unknown
time \(t_i + \delta_i\) with \(\delta_i \ge 0\), or never executes; the
transport may also \emph{redeliver} an attempt, so that it executes a
second time although the client issued it once. An executed attempt
creates one \emph{effect} \(e\), which is \emph{standing} until it is
neutralised. For every attempt the client receives an acknowledgement,
an error, a timeout or nothing, and it may issue reads through the read
paths that the service offers. The client can crash, and only what it
wrote to durable storage survives a crash. Other actors share the
service and create effects of their own.

A \emph{world} \(w\) fixes, for every attempt, whether, when and how
often it executes, together with the other actors' effects. Worlds are
drawn from a prior that describes the deployment, such as how often
requests are lost, how delays are distributed and how active other users
are. The client's \emph{history} \(h\) is the sequence of its own
actions, the responses and the read results, and \(\Pr[\cdot \mid h]\)
denotes the posterior over worlds. A \emph{policy} maps histories to
actions: re-issue, read, wait, neutralise an effect, or report. Two
worlds are \emph{observation-equivalent up to time \(t\)} if every
policy receives the same history in both until \(t\).

\textbf{Contract properties of a write.} Three properties belong to the
service's interface, not to the agent. The first is a \emph{native key}
(class K): attempts that carry the same key create at most one effect,
and a repeat returns the original outcome or an ``exists'' refusal. Real
keys come with a \emph{retention} \(r\), after which the service forgets
the key and treats a repeat as new. The second is the \emph{read model},
which says what reads reveal about standing effects. Under an
\emph{exact} read model, reads return all standing effects with their
fields. A \emph{lagged} read model is exact, but only for effects older
than a lag \(\lambda\). Under a \emph{content} read model, effects are
found by their content, and other actors may create effects with
identical content, which I call \emph{look-alikes}. An \emph{aggregate}
read model returns only a count or a sum. The third property is
\emph{reversibility}. A \emph{targeted} compensation \(c(e)\)
neutralises exactly the effect \(e\), such as deleting issue \#\(n\) or
row \(r\). A \emph{fungible} compensation neutralises one unit of an
aggregate, such as a decrement. Some effects cannot be neutralised at
all; a delivered e-mail is the standard example.

\textbf{In-flight bound.} An attempt that has not executed within
\(\Delta\) of its dispatch never executes. \(\Delta\) may be unknown to
the client, or the delay distribution \(F\) may have a heavy tail, in
which case any finite assumed bound is sometimes wrong.

\textbf{Attributability.} The central notion is the following.

\begin{quote}
\textbf{Definition 1 (Attributability).} Under a read model with lag
\(\lambda\), operation \(\ell\) is \emph{attributable} if, for every
time \(t\), the number \(N_\ell(t-\lambda)\) of standing effects of
\(\ell\) that executed before \(t-\lambda\) is a function of the
client's history at \(t\). Equivalently, no two worlds that are
observation-equivalent up to \(t\) disagree on \(N_\ell(t-\lambda)\).
\end{quote}

Attributability concerns effects that already exist. It does not make
the \emph{eventual} count identifiable, because an attempt may still be
in flight. That additionally requires the in-flight bound: once every
attempt is older than \(\Delta\), the eventual count equals the
attributable count. The theory below separates the two on purpose,
because attribution is a property of the API, whereas the in-flight
bound is a property of the network and the service's queues.

Five mechanisms make an operation attributable (Table 1). A native key
bounds \(N_\ell\) by one and answers queries for it. A \emph{client
tag}, an operation id \(\tau\) written into a field that the service
persists and returns, lets a read count exactly the effects that carry
\(\tau\), provided tags are unique per operation and no one else writes
\(\tau\) into an effect. \emph{Exclusivity}, being the only writer to an
aggregate, makes the change in the aggregate equal to \(N_\ell\). If no
one else writes with the client's account and the service records
authorship, the account itself identifies the client's effects, as the
follow-up study measures (Section 9). Finally, an id returned by the
service identifies the effect, but it is lost exactly when the reply is
lost, which is the case that matters.

\begin{longtable}[]{|
  >{\raggedright\arraybackslash}p{(\columnwidth - 2\tabcolsep - 2.1pt - 4\tabcolsep) * \real{0.0840}}|
  >{\raggedright\arraybackslash}p{(\columnwidth - 2\tabcolsep - 2.1pt - 4\tabcolsep) * \real{0.4580}}|
  >{\raggedright\arraybackslash}p{(\columnwidth - 2\tabcolsep - 2.1pt - 4\tabcolsep) * \real{0.4580}}|}
\caption{\textbf{Table 1.} Five ways an operation becomes attributable.
The returned id, the most common API design, gives no attributability in
the ambiguous case.}\tabularnewline
\noalign{\global\arrayrulewidth=1.1pt}\arrayrulecolor{cvdeep}\hline\noalign{\global\arrayrulewidth=0.4pt}\arrayrulecolor{cvgrid}\rowcolor{cvhead}
\begin{minipage}[b]{\linewidth}\raggedright\bfseries
Source
\end{minipage} & \begin{minipage}[b]{\linewidth}\raggedright\bfseries
Why \(N_\ell\) becomes identifiable
\end{minipage} & \begin{minipage}[b]{\linewidth}\raggedright\bfseries
Example in my stack
\end{minipage} \\
\hline
\endfirsthead
\noalign{\global\arrayrulewidth=1.1pt}\arrayrulecolor{cvdeep}\hline\noalign{\global\arrayrulewidth=0.4pt}\arrayrulecolor{cvgrid}\rowcolor{cvhead}
\begin{minipage}[b]{\linewidth}\raggedright\bfseries
Source
\end{minipage} & \begin{minipage}[b]{\linewidth}\raggedright\bfseries
Why \(N_\ell\) becomes identifiable
\end{minipage} & \begin{minipage}[b]{\linewidth}\raggedright\bfseries
Example in my stack
\end{minipage} \\
\hline
\endhead
\endlastfoot
Native key & the service enforces \(N_\ell \le 1\) and answers
``exists'' & TigerBeetle transfer id; S3 conditional create \\ \hline
Client tag & every effect carries \(\tau\); a read returns all effects
with \(\tau\) & Gitea issue body, mail Message-ID, row column, stream
field \\ \hline
Exclusivity & only the client writes, so the aggregate's change equals
\(N_\ell\) & a counter with no other writer \\ \hline
Own account & only the client writes with its account, and the service
records the author & Gitea issues and comments when other users use
other accounts \\ \hline
Returned id & the reply names the new effect & lost together with the
reply \\ \hline
\end{longtable}

\textbf{Guarantees.} I distinguish three guarantees. \emph{Prevention}
means that \(\ell\) executes exactly once over the episode; it is the
target of keys and of prior agent benchmarks {[}8{]}. \emph{Convergence
at horizon \(H\)} means that from time \(H\) on exactly one effect of
\(\ell\) stands, including after every in-flight attempt has landed, and
that no effect of another operation or actor has been neutralised; a
transient duplicate is allowed if it is removed by \(H\). \emph{Safe
abstention} means that, where neither guarantee is achievable, the
client reports \(\ell\) as unknown and never claims completion; it is a
property of the report, not of the state. Prevention implies convergence
once the single execution has landed, provided no other actor's effect
was touched; convergence does not imply prevention. My thesis is that
convergence is the right target wherever prevention is provably
impossible and effects are attributable and reversible.

\textbf{Primary metric.} The primary outcome of every episode is
therefore \emph{residual inconsistency at quiescence}. After the client
finishes and every in-flight request has landed, an episode is
inconsistent if a required effect is missing, if a surplus effect
stands, or if another actor's effect was removed. I also report
duplicate executions including compensated ones, which is the prevention
metric of prior work; irreversible duplicates; missing effects;
collateral damage; and \emph{overclaim}, an inconsistent final state
with nothing reported as uncertain.

\hypertarget{what-recovery-can-achieve}{%
\section*{4. What recovery can
achieve}\label{what-recovery-can-achieve}}
\addcontentsline{toc}{section}{4. What recovery can achieve}

I now state which guarantee is achievable under each contract and read
model. Proofs are in Appendix A. Several of the arguments are short
indistinguishability arguments. Their value lies less in any single
proposition than in the map they add up to (Table 2), which says, for
every combination of contract, read model and fault, which guarantee
some client-side policy can achieve. Propositions 1--6 and the
predictions of the map that the hypotheses test, including where Settle
itself must fail, were frozen with the analysis code before any test
episode ran. The retention clause of Proposition 1, the corrected
statements of Propositions 2, 3, 5 and 6, and Proposition 7 were worked
out while writing, after the runs, and I use them to interpret the
results rather than as tested predictions. Relative to the frozen map,
Table 2 changes the prevention column in two rows (prevention by waiting
is possible for reversible and exclusively written aggregate writes when
\(\Delta\) is known) and the wording of three others; the convergence
predictions that the hypotheses test are unchanged. The frozen version
is in the released repository.

\textbf{Proposition 1 (without a known in-flight bound, prevention needs
a durable key, within its retention).} Suppose the service deduplicates
a key for a retention period \(r\) after it first records it. Re-issuing
with the same key then prevents duplicates if the key is a function of
\emph{durable} state and every execution of every attempt falls within
\(r\) of the first execution. If the number of attempts is capped, this
gives at most one execution, and exactly one only if some attempt
executes. With a known in-flight bound, a sufficient condition is
\(t_{\text{last}} - t_1 + \Delta < r\). If the key is regenerated after
a client crash, as an in-memory random key would be, two attempts with
different keys can both execute, and prevention fails. If an attempt
executes later than \(r\) after the first execution, the service no
longer recognises the key, and the operation is protected only by
whatever attribution the contract offers apart from the key. Retention
that ends with a state change rather than after a time, such as a
deduplication key that lasts while an alert is open, behaves in the same
way, with the state change in place of \(r\). Without a key, no policy
that acts only on observations prevents duplicates under late commits
when \(\Delta\) is unknown {[}8{]}.

The retention condition matters in practice. A key that expires after
one hour does not protect an agent that resumes after a crash the next
morning, or one that waits for a human approval before retrying. The
census (Section 5) finds documented retentions from one hour to at least
thirty days, alongside conditional creates whose protection lasts as
long as the object exists.

\textbf{Proposition 2 (convergence for attributable, reversible
writes).} Assume the following:

\begin{enumerate}
\def\labelenumi{(\alph{enumi})}
\tightlist
\item
  the tag \(\tau\) of \(\ell\) can be recovered after a crash, for
  example from an intent record written before the first dispatch;
\item
  every executed attempt persists \(\tau\) unchanged;
\item
  the read path is complete and sound for \(\tau\) after lag
  \(\lambda\);
\item
  the in-flight bound \(\Delta\) holds;
\item
  a targeted compensation exists, and it is idempotent or itself keyed
  or tagged, so that it can be retried through the faulty channel;
\item
  no other actor neutralises effects that carry \(\tau\);
\item
  tags are unique per operation, and no other actor writes \(\tau\) into
  an effect.
\end{enumerate}

Consider the following policy, which proceeds in rounds. In each round
it waits until \(t_{\text{last}} + \Delta + \lambda\) and counts \(N\),
the effects that carry \(\tau\). If \(N = 0\), it re-issues with the
same \(\tau\) and starts a new round. If \(N > 1\), it compensates
\(N-1\) effects, chosen by id. The state is then consistent from the end
of the final round on. The number of rounds is finite with probability
one if each attempt executes with positive probability, and a policy
that caps the number of rounds, as Settle does, reports \emph{unknown}
when the cap is reached. If (d) is violated and delays follow a
distribution \(F\), the round that ends at horizon \(H\) leaves an
inconsistent state only if some attempt executes after \(H - \lambda\).
If attempts execute independently of each other, that happens with
probability at most \(\sum_i p_i\,\big(1 - F(H - \lambda - t_i)\big)\),
summed over the attempts \(i\) dispatched at \(t_i\), where \(p_i\) is
the probability that attempt \(i\) executes at all; this is a statement
about the round as a whole, not conditional on what the round observed.
The policy needs no assumption about the agent or about the order of
faults.

\textbf{Proposition 3 (irreversible writes: convergence collapses to
prevention).} If an effect cannot be neutralised, convergence requires
that \(\ell\) never executed twice, so convergence equals prevention,
and Proposition 1 applies. With a known \(\Delta\), the policy ``wait
\(\Delta + \lambda\), count by \(\tau\), re-issue only if \(N = 0\)''
prevents every duplicate that the \emph{client} could cause. The tag is
what makes the test \(N = 0\) exact; Proposition 4 shows that a test by
content is not exact. No client policy, however, prevents a duplicate
that the \emph{transport} creates by redelivering a request the client
issued once. For irreversible writes without a key, a redelivered
request is therefore an unrecoverable duplicate, and the best a client
can do is to detect it by its tag and report it. The same
wait-then-count policy would also prevent duplicates of reversible
writes. Settle does not use it for them, because convergence lets it
replace the wait by compensation.

\textbf{Proposition 4 (content verification is unsound under
look-alikes).} Consider a world \(w_1\) in which the client's attempt
executed and no look-alike exists, and a world \(w_2\) in which the
client's attempt did not execute and another actor created an effect
with identical content. Under a content read model the two worlds are
observation-equivalent, while \(N_\ell\) differs. Every content-only
policy is therefore wrong in one of them. A policy that skips the write
leaves \(\ell\) missing in \(w_2\). A policy that re-issues duplicates
\(\ell\) in \(w_1\). And a policy that re-issues and then removes one of
the two matching records either removes the other actor's effect in
\(w_2\) or removes its own new effect, which leaves \(\ell\) missing in
\(w_2\). A client tag restores attributability, because look-alikes do
not carry \(\tau\). Fields that real read models often expose, such as
the author account and the creation time, separate the two worlds only
when the other actor uses a different account or writes at a
distinguishable time; they are a partial substitute for a tag, not a
guarantee. The follow-up study (Section 9) measures both cases.

\textbf{Proposition 5 (aggregates: exclusivity is the only general
route).} Under an aggregate read model, \(\ell\) is attributable if no
other actor writes to the aggregate during the ambiguity window. If
other actors may write the same unit, or units of unknown size, then a
world in which the client's attempt executed and theirs did not gives
the same aggregate as a world in which theirs executed and the client's
did not. Neither prevention nor convergence is then achievable without a
key, and fungible compensation does not help, because the number of
surplus units is not identifiable. The only correct behaviour is safe
abstention. (Concurrent writes of known and distinct sizes would keep
the count identifiable. My shared-aggregate worlds, like most real
counters, do not have that property.)

\textbf{Proposition 6 (the tag must survive a crash).} Suppose the
client can crash after dispatching an attempt and before recording its
outcome. After the restart, the client must issue \(\ell\) if it was
never dispatched and must not issue it blindly if it was. If the
persisted state does not show whether \(\ell\) may have been dispatched,
and the tag of a possible dispatch cannot be recovered, the two cases
are observation-equivalent to the client, which then either duplicates
or omits \(\ell\). Two kinds of durable state suffice to decide
correctly: a write-ahead intent record that carries the tag, or a tag
derived deterministically from durable step identity, provided every
unfinished step is treated after a restart as possibly dispatched and
resolved by waiting and counting. A durable step checkpoint
\emph{without} a recoverable tag, which is what my drivers keep and what
the S−j ablation is left with, is not enough: it shows that the step may
have been dispatched, but leaves the client nothing to count by except
content, which Proposition 4 shows to be unsound. Workflow engines
persist the \emph{scheduling} of an activity durably but execute the
activity at least once {[}12{]}. Their durable identity is therefore
available, but unless it is written into the effect it does not make the
effect attributable: durability without attribution yields at-least-once
effects. Likewise, an agent framework that replays a pending tool call
after a restart turns every crash during a write into a blind re-issue.

\textbf{Proposition 7 (a policy that stops must bet).} Let \(\ell\) be
attributable, and let \(t_d\) be the time of the policy's last action on
\(\ell\), after which it neither observes nor acts on \(\ell\). For a
reversible write \(t_d\) is the final reconciliation; for an
irreversible write it is the decision whether to re-issue. Suppose an
attempt dispatched at \(t_0\) is still unresolved at \(t_d\). Let
\(w_0\) be the worlds in which it never executes, and \(w_L\) those in
which it executes after \(t_d - \lambda\), too late to be seen by the
last read. Up to \(t_d\) the two kinds of world are
observation-equivalent. At \(t_d\), the policy decides how many of its
\emph{other} effects of \(\ell\) remain standing: none, by not
re-issuing or by compensating its re-issue, or one, by re-issuing or
keeping the re-issue. The final count is that number plus one in
\(w_L\), so every choice is wrong in either \(w_0\) or \(w_L\). Hence,
for every policy \(\pi\) and every history \(h\) that satisfies the
premise,
\[\Pr[\pi \text{ inconsistent} \mid h] \;\ge\; \min\{\Pr[w_0 \mid h],\; \Pr[w_L \mid h]\}.\]
Settle always keeps one effect standing, so its error at such a history
is \(\Pr[w_L \mid h]\), plus \(\Pr[w_0 \mid h]\) times the probability
that its own re-issue fails to execute. It attains the bound exactly
when its re-issue executes and a late landing is no more likely, given
the history, than an attempt that never executes.

Proposition 7 makes two points about the residual failures I measure.
First, they are the price of a bet that every stopping policy must make,
not a defect of Settle's rule: a policy that bet the other way would
fail in the worlds where the request was lost instead. Settle bets on
the lost request because then the operation is complete and the error,
if any, is a surplus effect rather than a missing one. In a deployment
where late commits after the horizon are more likely than lost requests,
the opposite bet has the lower error. My late-tail worlds are designed
so that every request eventually commits, so there the opposite bet
would have been right. Second, at a fixed decision time no rule does
better than the bound, so the residual risk falls only if \(t_d\) moves
later or if the posterior shifts. For an irreversible write, a later
\(t_d\) means a longer wait before the re-issue, which the agent
experiences as latency. For a reversible write the reconciliation can
instead continue \emph{after} the agent has reported, because
compensation needs no agent. Section 8.6 measures the trade-off.

\begin{longtable}[]{|
  >{\raggedright\arraybackslash}p{(\columnwidth - 2\tabcolsep - 2.9pt - 8\tabcolsep) * \real{0.0913}}|
  >{\raggedright\arraybackslash}p{(\columnwidth - 2\tabcolsep - 2.9pt - 8\tabcolsep) * \real{0.1261}}|
  >{\raggedright\arraybackslash}p{(\columnwidth - 2\tabcolsep - 2.9pt - 8\tabcolsep) * \real{0.2609}}|
  >{\raggedright\arraybackslash}p{(\columnwidth - 2\tabcolsep - 2.9pt - 8\tabcolsep) * \real{0.2609}}|
  >{\raggedright\arraybackslash}p{(\columnwidth - 2\tabcolsep - 2.9pt - 8\tabcolsep) * \real{0.2609}}|}
\caption{\textbf{Table 2.} Achievability map. For each contract and read
model: whether \emph{some} client-side policy can prevent duplicates or
make the final state converge, and the policy that attains it. Class K
has a native key, class A is key-less but taggable, class U exposes only
an aggregate, and class C is a key-less write that can be found only by
content. The experiments fill each row with measured rates (Table
7).}\tabularnewline
\noalign{\global\arrayrulewidth=1.1pt}\arrayrulecolor{cvdeep}\hline\noalign{\global\arrayrulewidth=0.4pt}\arrayrulecolor{cvgrid}\rowcolor{cvhead}
\begin{minipage}[b]{\linewidth}\raggedright\bfseries
Contract
\end{minipage} & \begin{minipage}[b]{\linewidth}\raggedright\bfseries
Read model
\end{minipage} & \begin{minipage}[b]{\linewidth}\raggedright\bfseries
Prevention
\end{minipage} & \begin{minipage}[b]{\linewidth}\raggedright\bfseries
Convergence
\end{minipage} & \begin{minipage}[b]{\linewidth}\raggedright\bfseries
Attaining policy
\end{minipage} \\
\hline
\endfirsthead
\noalign{\global\arrayrulewidth=1.1pt}\arrayrulecolor{cvdeep}\hline\noalign{\global\arrayrulewidth=0.4pt}\arrayrulecolor{cvgrid}\rowcolor{cvhead}
\begin{minipage}[b]{\linewidth}\raggedright\bfseries
Contract
\end{minipage} & \begin{minipage}[b]{\linewidth}\raggedright\bfseries
Read model
\end{minipage} & \begin{minipage}[b]{\linewidth}\raggedright\bfseries
Prevention
\end{minipage} & \begin{minipage}[b]{\linewidth}\raggedright\bfseries
Convergence
\end{minipage} & \begin{minipage}[b]{\linewidth}\raggedright\bfseries
Attaining policy
\end{minipage} \\
\hline
\endhead
\endlastfoot
K (native key) & any & yes, with a durable key, while every execution is
within retention \(r\) & yes, within retention & key + recoverable
identity \\ \hline
A, reversible & exact or lagged, \(\Delta\) known & yes, by waiting
\(\Delta + \lambda\) before each re-issue (P3), except against
redelivery & yes, without the wait, and also against redelivery that
lands before the final count & recoverable tag + horizon + targeted
compensation \\ \hline
A, reversible & lagged, heavy-tailed \(\delta\) & no & per round, with
probability at least \(1 - \sum_i p_i(1 - F(H - \lambda - t_i))\) for
independent attempts; no stopping policy beats the bound of P7 & same;
residual risk falls as \(H\) grows \\ \hline
A, irreversible & exact or lagged & only if \(\Delta\) is known and
waited out, and never against redelivery & equals prevention & tag +
wait, then count; report redelivered duplicates \\ \hline
C, content only & content with look-alikes & no & no & safe abstention
(a tag, where one exists, makes the write class A) \\ \hline
U, exclusive writer & aggregate & yes, by waiting \(\Delta + \lambda\)
before each re-issue, except against redelivery & yes, if fungibly
reversible & delta + horizon \\ \hline
U, concurrent writers & aggregate & no & no & safe abstention \\ \hline
\end{longtable}

Table 2 summarises the propositions and yields three predictions that
the experiments test. Idempotency keys and workflow engines should match
Settle on class-K writes and fall short on class-A writes. Every method
that recognises its own effects by content should fail under
look-alikes. And no method should succeed in the last row; the best a
method can do there is to say so.

\hypertarget{how-common-is-each-contract-a-census-of-52-apis}{%
\section*{5. How common is each contract? A census of 52
APIs}\label{how-common-is-each-contract-a-census-of-52-apis}}
\addcontentsline{toc}{section}{5. How common is each contract? A census
of 52 APIs}

The theory makes the contract of a write, rather than the agent, decide
what recovery can achieve. Whether this matters depends on how real
services are distributed over the cells of Table 2. I therefore coded
the write contracts of 52 widely used APIs from their official
documentation. The census is exploratory; it was coded after the test
runs and was not pre-registered (Section 7).

\textbf{Sample and coding.} I chose APIs that agents are commonly
connected to, in four domains of twelve to fifteen APIs each: payments
and commerce, developer and project tools, messaging and e-mail, and
CRM, data and cloud services. For each API I examined the \emph{create}
operation an agent would typically call (for example, creating an issue,
sending a message, inserting a record or creating a payment) and
recorded six properties from the vendor's documentation. The first is
whether the operation accepts a client-supplied key that the
\emph{server} deduplicates, and for how long the key is retained. The
second is whether a client-controlled field is persisted with the
created object, such as a metadata map, an external id, a custom field,
a header or a label. The third is whether objects can be retrieved by
the exact value of that field. Lookups are coded as exact, as a search
that the documentation describes as fuzzy or eventually consistent, as
possible only by scanning a list of recent objects, or as impossible.
The fourth is whether the field is hidden from end users. The fifth is
whether the created effect can be undone through the API. The sixth is
the URL of every documentation page relied on.

Each API receives a class from these properties: \textbf{K} if it has a
server-side key; \textbf{A} if, without a key, it has a taggable field
with an exact lookup; \textbf{A−} if the field can be found only through
a lagged search or by scanning lists; and \textbf{C} if there is no
usable field or no lookup at all. A second pass re-checked a random
sample of 12 entries against the documentation, starting from the first
pass's codes and sources. It confirmed all 12 class assignments and
found 1 factual error in a detail that did not affect the class (an
additional search endpoint for one payment API). Twelve agreements out
of twelve bound the class error rate only loosely, at about 25\% with
95\% confidence, so individual entries should be read with their
sources. The full coded table, with documentation links, is part of the
released artifacts, and Appendix E lists it.

\begin{figure}
\centering
\includegraphics[width=0.66\textwidth,height=\textheight]{fig_census.png}
\caption{\textbf{Fig. 2.} Write contracts of 52 APIs by domain, coded
from vendor documentation. K: the create operation accepts a
server-deduplicated key. A: no key, but a client-controlled field with
an exact lookup. A−: a field that can be found only through a lagged
search or by scanning lists. C: no usable field or no lookup.}
\end{figure}

\textbf{Keys are rare outside payments.} Only 14 of the 52 APIs (27\%)
have a server-side key on the examined operation (Fig. 2). Eight of the
twelve payment APIs have one; the other six are in the CRM, data and
cloud group: Zendesk, Microsoft Graph calendar events, PagerDuty, and
the conditional creates of S3, DynamoDB and Elasticsearch. None of the
twelve developer and project tools and none of the thirteen messaging
and e-mail services offers a key. Eleven key-less APIs have something
\emph{close to} a key, with weaker or undocumented guarantees: upserts
on a client-chosen field (Salesforce, Airtable) without a documented
concurrency guarantee; uniqueness constraints that reject a repeat with
an error instead of returning the original outcome (Razorpay's order
receipt, Confluence page titles, HubSpot's unique properties);
client-chosen ids whose duplicate handling is undocumented (Linear) or
whose collisions the provider explicitly does not guarantee to detect
(Google Calendar); duplicate checks that last thirty seconds (Braintree)
or minutes (Discord's nonce, SQS FIFO queues); and deduplication within
one producer session (Kafka). An agent can build on several of these,
but none is documented to give the guarantee of Proposition 1, and I do
not count them as keys.

\textbf{Keys expire.} Retention is documented for most keys, and several
are short. Mollie keeps keys for one hour, Zendesk for two and PayPal
for six by default. PagerDuty's deduplication key lasts only while the
alert is open, so a retry after it has been resolved opens a new
incident. Square, Wise and Microsoft Graph do not document retention at
all. By Proposition 1, a key protects only re-issues within its
retention, and agents that resume after a crash, a deployment or a human
approval can exceed one hour. Worse, five of the eight keyed payment
APIs can find a payment by a client field only by scanning a list of
recent payments, one (Adyen) offers no lookup at all, and one (Stripe)
offers only an eventually consistent search. Once a key has expired,
these APIs leave the client no exact way to attribute.

\textbf{Most key-less APIs are taggable.} Of the 38 key-less APIs, 34
accept a client-controlled field with the write, and 31 (82\%) persist
one that can also be looked up. Of these, 21 (55\% of the key-less APIs)
offer an exact lookup by the field, which makes them class A in the
sense of Section 4; under a strict reading that treats Mailgun's events
API and the Sent Items copy of Microsoft Graph mail as lagged, the count
is 19. For class-A APIs Proposition 2 applies as soon as the effect is
reversible: 14 of the 21 class-A operations can be undone through the
API, and 2 more only partly (GitHub's REST API closes rather than
deletes an issue, and Notion moves a page to the trash). An exact lookup
on GitHub needs the tag as a label, one per operation; a tag in the
issue body can be found only through the lagged search index. Most of
the remaining class-A operations send e-mail, which cannot be recalled,
so Proposition 3 applies to them. Ten APIs are class A−. Jira's and
Freshdesk's searches, for example, are documented as possibly missing
recent writes (Freshdesk: for a few minutes), and Slack messages and
Google Sheets rows can be found only by scanning recent items. Here
attribution is possible, but only under a read lag \(\lambda\) that is
sometimes undocumented, or at the cost of a scan.

\textbf{A minority offers no attribution at all.} Class C, with no
usable tag field or no lookup, covers 7 APIs (18\% of the key-less
ones): Microsoft Teams channel messages, Twilio SMS, Telegram, WhatsApp,
Intercom conversations, Google Pub/Sub and standard SQS queues. Their
writes cannot be tagged, or a tag cannot be looked up. Under a strict
reading, which also moves Discord and Amazon SES to class C because
their lookups depend on undocumented or sampled data, the share is 24\%.
For class-C writes the theory leaves only safe abstention; waiting out a
known in-flight bound tells the client that no attempt is still pending,
but not whether one executed.

\textbf{Tags need not be visible.} A tag written into an issue body is
visible to the people who read the issue. Among the 31 taggable key-less
APIs, 9 offer a field the documentation describes as hidden from end
users, for example Asana's external data, Shopify's metafields, Jira's
issue properties, Postmark's and SendGrid's message metadata, and mail
headers. Another 10 offer only visible fields, such as GitHub labels or
a Notion property, and for 12 the documentation does not say.

\textbf{What the census implies.} The census supports the premise of
this paper on documentary evidence. Prevention by a server-side key is
available for about a quarter of the APIs I examined and concentrated in
payments. Convergence by attribution is available for most of the rest.
Only a minority of services, mostly messaging services, leaves nothing
but abstention. The census has limits. It covers documented behaviour,
not behaviour under fault, and the sample is purposive rather than
random. It codes one create operation per API, and several ``unclear''
entries could move with better documentation. The experiments that
follow test whether the documented contract classes behave as the theory
predicts, on real services and under real faults.

\hypertarget{settle-recovery-by-attribution}{%
\section*{6. Settle: recovery by
attribution}\label{settle-recovery-by-attribution}}
\addcontentsline{toc}{section}{6. Settle: recovery by attribution}

Settle is a layer between an agent and its tools that turns the
constructive side of Section 4 into a protocol. The agent, scripted or a
language model, calls tools as usual. Settle intercepts every write and,
in the language-model setting, every deletion the agent requests. It
returns one of three observations: \emph{ok} with the effect's id,
\emph{unknown} with a reason, or the service's own error when the
outcome is known. Settle has six parts.

\textbf{Operation identity.} Each logical intent receives one operation
id, which is kept across retries, restarts and rephrasings. For the
scripted agent, the intent is the plan step. For a language-model agent,
a write is bound to an intent by its \emph{intent fields}: the title of
an issue, the recipient and subject of a mail, the item of a record. A
rephrased re-issue of the same issue therefore maps to the same
operation. In my implementation this binding reuses the harness's
mapping of model writes onto the task's plan steps by these fields
(Appendix C). A deployed layer has no plan to map onto and would bind by
the intent fields alone. It would then also merge two deliberately
identical writes, such as the same reminder sent twice on purpose,
unless the agent marks them as distinct. The follow-up study runs this
plan-free binding (S-np) and finds identical outcomes (Section 9).

\textbf{Durable intent journal.} Before the first dispatch, Settle
appends \texttt{intent\ →\ \{op\ id,\ arguments,\ aggregate\ before\}}
to a journal and calls \texttt{fsync}. Each attempt moves the record
through the states \texttt{dispatching}, then \texttt{confirmed},
\texttt{in\_doubt} or \texttt{unknown}. After a crash, a record found
\texttt{dispatching} or \texttt{in\_doubt} is resolved before anything
is re-issued, which is what Proposition 6 requires.

\textbf{Attribution.} Settle uses the strongest attribution that the
contract offers. For a class-K write, the op id is the idempotency key:
TigerBeetle derives the transfer id from it, and the object store uses a
conditional create. For a class-A write, the op id is written into a
field that the service persists, such as an \texttt{op-id:} line in an
issue or comment body, the mail \texttt{Message-ID} header, a reference
column, a stream field or object metadata. The agent is not told about
the tag, although it can see it when it reads the record back. For a
class-U write, Settle records the aggregate before the write and uses
the change, which is valid only under exclusivity (Proposition 5). A
write that can be found only by content (class C in Table 2; in my
stack, the bulk e-mail) is not attributable. For such writes the
implementation falls back to a content match, as the baselines do, and
lists the operation as uncertain whenever it had to re-issue it. It does
\emph{not} abstain when the content match succeeds at the first check,
which is the case a look-alike exploits; Section 8.5 reports the
consequence. A design that abstains on every ambiguous class-C write
would report these failures, at the cost of reporting uncertainty in
every ambiguous class-C episode; the follow-up study measures it (S2).

\textbf{Resolution of unknown outcomes.} On an error or timeout, Settle
re-issues a keyed write with the same key. For a taggable write, it
waits out the read path's documented lag and counts the effects carrying
the tag. If there are one or more, the operation is done and the id of
the kept effect is returned. If there are none, the write is re-issued
under the same tag. For an irreversible write such as an e-mail, Settle
first waits until the assumed in-flight bound \(\Delta\) has passed
since the last attempt, so that a re-issue cannot overtake a late commit
(Proposition 3). For an aggregate shared with other writers it stops and
returns \emph{unknown}. After three attempts it returns \emph{unknown}
in every case.

\textbf{Reconciliation horizon.} When the agent finishes, Settle keeps
reconciling every operation that was ever in doubt until
\(t_{\text{last}} + \Delta + \lambda\). For each such operation it
counts the effects by tag, re-issues if none exists, and marks surplus
effects for compensation. The final sweep also covers operations that
were confirmed at once, which catches transport-level redelivery. If no
operation was ever in doubt, the sweep runs at once without waiting, so
it catches a redelivery only if the duplicate is already visible when
the agent finishes. My main redelivery fault delivers the duplicate
immediately; a delayed redelivery, as from a queue's visibility timeout,
needs the sweep to wait in every episode, which the follow-up study
tests (S2).

\textbf{Evidence-gated compensation.} A surplus effect is neutralised
only if it carries the operation's tag and is not the effect that the
operation keeps, which is the one whose id later steps have already
received. Compensations are themselves writes. Settle's compensations in
my stack are deletions, which are idempotent, and decrements of
exclusively written aggregates, which are not. In the experiments every
compensation travelled on a fault-free path, so assumption (e) of
Proposition 2 was not stressed. Three further consequences follow. An
effect that cannot be attributed is never compensated, which avoids the
\emph{orphaned compensation} anomaly {[}9{]}. An irreversible surplus is
reported as an unrecoverable duplicate instead of being hidden. And in
the language-model setting, the same rule guards the agent's own
deletion tools: a deletion of an effect that carries none of the task's
tags, or of the only effect of a completed operation, is refused with a
short explanation. I added this rule on the development split, after
observing that a careful agent which sees two identical issues ``cleans
up'' by deleting the other user's issue (Section 8). Clean-up by content
is exactly the unsound compensation of Proposition 4.

The pseudocode below summarises the write path for a taggable,
reversible tool; the released code contains the full implementation.

\begin{Shaded}
\begin{Highlighting}[]
\NormalTok{write(intent, args):}
\NormalTok{  rec = journal.get(intent)}
\NormalTok{  if rec is None:}
\NormalTok{    rec = journal.put(intent, op=new\_id(), state=intent)}
\NormalTok{  elif rec.confirmed:                     \# done before a crash or re{-}issue}
\NormalTok{    return ok(rec.value)}
\NormalTok{  elif rec.state in (dispatching, in\_doubt):    \# restart after a crash:}
\NormalTok{    return resolve(intent)       \# wait lag (+ Δ if irreversible), count}
\NormalTok{  journal.put(intent, state=dispatching)  \# write{-}ahead, before sending}
\NormalTok{  r = dispatch(args + tag(rec.op))}
\NormalTok{  if r.ok: return ok(journal.confirm(intent, r.id))}
\NormalTok{  journal.put(intent, state=in\_doubt)}
\NormalTok{  sleep(lag(tool))}
\NormalTok{  found = read\_by\_tag(rec.op)             \# attribution, not content}
\NormalTok{  if found: return ok(journal.confirm(intent, kept(found)))}
\NormalTok{  return write\_again(intent)        \# same tag; ≤ 3 attempts, then unknown}

\NormalTok{finish():}
\NormalTok{  if journal.ever\_in\_doubt():}
\NormalTok{    sleep\_until(max(last\_attempt) + Δ + λ)    \# reconciliation horizon}
\NormalTok{  for op in journal.ever\_in\_doubt() ∪ journal.confirmed():}
\NormalTok{    found = read\_by\_tag(op)}
\NormalTok{    if not found: re{-}issue(op)}
\NormalTok{    for e in found − \{kept(op)\}:}
\NormalTok{      compensate(e)                   \# only effects carrying op\textquotesingle{}s tag}
\NormalTok{  report(unknown = unattributable re{-}issued ops, shared aggregates,}
\NormalTok{         irreversible duplicates)}
\end{Highlighting}
\end{Shaded}

\textbf{Assumptions and parameters.} Settle does not assume an
idempotency key on the service, a well-behaved agent (the agent may
retry, rephrase, crash or try to delete things), or a known order of
faults. It does assume that the tag survives byte-exact in the persisted
field, assumption (b) of Proposition 2, which the environment validation
checks per service, and that no one else copies the tag into another
record, assumption (g). A tag in a visible body line can be copied, for
instance when a user quotes the record; requiring the tag \emph{and} the
client's own author account before counting or compensating an effect
would close most of this gap, and my implementation does not yet do so.
Its parameters are the assumed in-flight bound \(\Delta\) and the
documented read lags. \(\Delta\) was set on the development split to 20
s and frozen before testing. Where it is violated, Propositions 2 and 7
predict the residual risk; the experiments include a heavy-tailed delay
distribution to measure that risk, and Section 8.6 traces how it moves
with \(\Delta\).

\textbf{Ablations.} Four ablations each remove one component. S−j keeps
the intent journal, and with it the tag, in memory, so by Proposition 6
it tests whether the tag must survive a crash; S−t matches effects by
content instead of by tag; S−h verifies once and has no horizon; and S−c
flags surplus effects without compensating them. Each ablation should
fail in the cells where the theory says its component is needed, and
should be close to Settle elsewhere.

\hypertarget{experimental-design}{%
\section*{7. Experimental design}\label{experimental-design}}
\addcontentsline{toc}{section}{7. Experimental design}

The design is a paired factorial. Every \emph{world} is replayed under
every recovery condition of a driver. A world fixes the task instance,
the focal write that receives the fault, the fault mode and the read
model the agent sees, so all comparisons between conditions are
within-world. The study was pre-registered internally: hypotheses,
analysis code, task definitions, prompts and the SHA-256 hashes of all
source files were frozen before the first test episode, and the test
presets refuse to run if any hashed file changes (Appendix D). The
freeze is recorded in the released repository rather than with a
third-party registry.

\hypertarget{services-tools-and-contracts}{%
\subsection*{7.1 Services, tools and
contracts}\label{services-tools-and-contracts}}
\addcontentsline{toc}{subsection}{7.1 Services, tools and contracts}

All services run unmodified in containers on one workstation: PostgreSQL
16 for records and stock levels; Gitea 1.24 {[}36{]} for issues and
comments, with its real asynchronous search indexer; Mailpit {[}37{]}
for SMTP with a search API; Redis 7 streams and counters {[}38{]};
SeaweedFS {[}39{]} as an S3 object store; and TigerBeetle {[}18{]} for
ledger transfers. The agent reaches them through thin tool adapters
(Table 3), each a few lines over the service's own API. Where Settle
writes a tag, it uses a field of the kind that the census found in real
APIs: a body line as in GitHub issues, a mail header as in Gmail, a
reference column as in Salesforce's external ids, a stream field as in
Kafka headers. The stack covers every row of the achievability map.

\begin{longtable}[]{|
  >{\raggedright\arraybackslash}p{(\columnwidth - 2\tabcolsep - 3.3pt - 10\tabcolsep) * \real{0.1393}}|
  >{\raggedright\arraybackslash}p{(\columnwidth - 2\tabcolsep - 3.3pt - 10\tabcolsep) * \real{0.0984}}|
  >{\raggedright\arraybackslash}p{(\columnwidth - 2\tabcolsep - 3.3pt - 10\tabcolsep) * \real{0.1311}}|
  >{\raggedright\arraybackslash}p{(\columnwidth - 2\tabcolsep - 3.3pt - 10\tabcolsep) * \real{0.1639}}|
  >{\raggedright\arraybackslash}p{(\columnwidth - 2\tabcolsep - 3.3pt - 10\tabcolsep) * \real{0.1721}}|
  >{\raggedright\arraybackslash}p{(\columnwidth - 2\tabcolsep - 3.3pt - 10\tabcolsep) * \real{0.2951}}|}
\caption{\textbf{Table 3.} Tool contracts. Class K has a native key,
class A is key-less but taggable, class U exposes only an aggregate, and
class C can be found only by content.}\tabularnewline
\noalign{\global\arrayrulewidth=1.1pt}\arrayrulecolor{cvdeep}\hline\noalign{\global\arrayrulewidth=0.4pt}\arrayrulecolor{cvgrid}\rowcolor{cvhead}
\begin{minipage}[b]{\linewidth}\raggedright\bfseries
Tool
\end{minipage} & \begin{minipage}[b]{\linewidth}\raggedright\bfseries
Service
\end{minipage} & \begin{minipage}[b]{\linewidth}\raggedright\bfseries
Class
\end{minipage} & \begin{minipage}[b]{\linewidth}\raggedright\bfseries
Reversible
\end{minipage} & \begin{minipage}[b]{\linewidth}\raggedright\bfseries
Tag field or key
\end{minipage} & \begin{minipage}[b]{\linewidth}\raggedright\bfseries
Read paths
\end{minipage} \\
\hline
\endfirsthead
\noalign{\global\arrayrulewidth=1.1pt}\arrayrulecolor{cvdeep}\hline\noalign{\global\arrayrulewidth=0.4pt}\arrayrulecolor{cvgrid}\rowcolor{cvhead}
\begin{minipage}[b]{\linewidth}\raggedright\bfseries
Tool
\end{minipage} & \begin{minipage}[b]{\linewidth}\raggedright\bfseries
Service
\end{minipage} & \begin{minipage}[b]{\linewidth}\raggedright\bfseries
Class
\end{minipage} & \begin{minipage}[b]{\linewidth}\raggedright\bfseries
Reversible
\end{minipage} & \begin{minipage}[b]{\linewidth}\raggedright\bfseries
Tag field or key
\end{minipage} & \begin{minipage}[b]{\linewidth}\raggedright\bfseries
Read paths
\end{minipage} \\
\hline
\endhead
\endlastfoot
create issue & Gitea & A & targeted (delete) & \texttt{op-id:} line in
body & list + body (exact); search (lagged) \\ \hline
add comment & Gitea & A & targeted & tag line in body & comment list \\ \hline
send e-mail & Mailpit SMTP & A & none & client \texttt{Message-ID} &
search by header or content \\ \hline
send bulk e-mail & Mailpit SMTP & C (content only) & none & none &
search by content \\ \hline
insert record(s) & PostgreSQL & A & targeted & \texttt{reference} column
& select \\ \hline
append event & Redis stream & A & targeted (XDEL) & \texttt{op\_id}
field & range scan \\ \hline
adjust stock & PostgreSQL & U & fungible (inverse) & none & sum only \\ \hline
increment counter & Redis & U & fungible (decrement) & none & value
only \\ \hline
create transfer & TigerBeetle & K & targeted (reversal) & 128-bit
transfer id & lookup by id \\ \hline
put object & SeaweedFS S3 & K & targeted & key +
\texttt{If-None-Match:\ *} & head \\ \hline
\end{longtable}

\textbf{Tasks.} Eight task templates (Table 4) define short workflows of
two to four writes across services, one of which (T5) inserts a batch of
six records. Each comes with an exact specification of the final state:
every write must leave exactly one standing effect, and every aggregate
must change by exactly the required amount. The development split
consists of five templates on instance 0, and every method parameter was
set on it. The test split uses all eight templates on fresh instances
100--102. Three templates (T3, T6, T8) were not used during development;
for these, only a fault-free and error check of the harness was run, on
an instance outside the test set.

\begin{longtable}[]{|
  >{\raggedright\arraybackslash}p{(\columnwidth - 2\tabcolsep - 2.5pt - 6\tabcolsep) * \real{0.0656}}|
  >{\raggedright\arraybackslash}p{(\columnwidth - 2\tabcolsep - 2.5pt - 6\tabcolsep) * \real{0.1721}}|
  >{\raggedright\arraybackslash}p{(\columnwidth - 2\tabcolsep - 2.5pt - 6\tabcolsep) * \real{0.4918}}|
  >{\raggedright\arraybackslash}p{(\columnwidth - 2\tabcolsep - 2.5pt - 6\tabcolsep) * \real{0.2705}}|}
\caption{\textbf{Table 4.} Task templates. Test-only templates: T3, T6,
T8.}\tabularnewline
\noalign{\global\arrayrulewidth=1.1pt}\arrayrulecolor{cvdeep}\hline\noalign{\global\arrayrulewidth=0.4pt}\arrayrulecolor{cvgrid}\rowcolor{cvhead}
\begin{minipage}[b]{\linewidth}\raggedright\bfseries
Template
\end{minipage} & \begin{minipage}[b]{\linewidth}\raggedright\bfseries
Workflow
\end{minipage} & \begin{minipage}[b]{\linewidth}\raggedright\bfseries
Writes (class)
\end{minipage} & \begin{minipage}[b]{\linewidth}\raggedright\bfseries
Stresses
\end{minipage} \\
\hline
\endfirsthead
\noalign{\global\arrayrulewidth=1.1pt}\arrayrulecolor{cvdeep}\hline\noalign{\global\arrayrulewidth=0.4pt}\arrayrulecolor{cvgrid}\rowcolor{cvhead}
\begin{minipage}[b]{\linewidth}\raggedright\bfseries
Template
\end{minipage} & \begin{minipage}[b]{\linewidth}\raggedright\bfseries
Workflow
\end{minipage} & \begin{minipage}[b]{\linewidth}\raggedright\bfseries
Writes (class)
\end{minipage} & \begin{minipage}[b]{\linewidth}\raggedright\bfseries
Stresses
\end{minipage} \\
\hline
\endhead
\endlastfoot
T1 & open an incident & issue (A), event (A), page e-mail (A,
irreversible), record (A) & mixed reversibility \\ \hline
T2 & pay out a batch & 2 transfers (K), record (A), e-mail (A, irr.) &
keyed and key-less in one task \\ \hline
T3 & ship a release & object (K), issue (A), comment (A), bulk e-mail
(C, irr.) & content-only irreversible write \\ \hline
T4 & receive stock & 2 stock adjustments (U), counter (U), record (A) &
aggregate reads \\ \hline
T5 & migrate a batch & 6 records in one non-atomic batch (A), event (A)
& partial failure \\ \hline
T6 & send customer notices & 3 e-mails (A, irr.), issue (A) &
irreversible-heavy \\ \hline
T7 & triage a ticket & issue (A), 2 comments (A), record (A) & dependent
steps, look-alike prone \\ \hline
T8 & refund a customer & transfer (K), record (A), e-mail (A, irr.) &
keyed payment + notification \\ \hline
\end{longtable}

\hypertarget{faults-and-read-models}{%
\subsection*{7.2 Faults and read models}\label{faults-and-read-models}}
\addcontentsline{toc}{subsection}{7.2 Faults and read models}

\textbf{Network and process faults.} The faults are real. They are
injected on the client path with Toxiproxy {[}40{]}, through one proxy
per worker, armed by a hook immediately before the focal write's bytes
are sent. In \emph{lost\_ack}, the request is delivered and committed
and the reply is cut. In \emph{not\_exec\_err}, the connection is cut
before the request is delivered, and the agent sees the same error as in
lost\_ack. In \emph{held}, the reply is withheld past the tool deadline
after the request has committed; in \emph{not\_exec\_to}, the request is
held upstream and dropped, and the agent sees the same timeout. In
\emph{late}, the request reaches the service 12 s after dispatch, long
after the 4 s tool deadline has expired. \emph{late\_tail} is like late
but draws the delay per world from a log-uniform distribution on 5--120
s, so that in part of the worlds it exceeds Settle's assumed bound of 20
s. In \emph{partial}, a non-atomic batch is cut after its third row. In
\emph{redelivery}, the transport delivers the same request twice. In
\emph{crash\_after} and \emph{crash\_before}, the agent process exits
(\texttt{os.\_exit}) right after or right before the focal dispatch and
is restarted from its checkpoint.

Each pair of modes (lost\_ack and not\_exec\_err, held and
not\_exec\_to, crash\_after and crash\_before) gives the agent the same
observation with opposite hidden outcomes. Transport errors are
normalised to one exception type and message, so the pairs are
observation-equivalent by construction. The two network pairs make the
design symmetric in what it penalises, though not in size (held occurs
in 42 of the H1 worlds and not\_exec\_to in 96; Table 6): a method that
re-issues too eagerly fails in one mode of each pair, and a method that
re-issues too reluctantly fails in the other. The crash pair is not
symmetric in this sense. After a crash before dispatch, the driver's
checkpoint re-issues the step in every condition, so crash\_before ends
consistent everywhere, and the pair penalises only eager re-issuing. An
environment validation checks on every service, from the service's own
records, that each mode produces its intended hidden outcome; 493 of 510
trials did, and the exceptions are TigerBeetle's not\_exec modes, whose
client library retries internally (Appendix B).

\textbf{Read models.} The read model is varied per world. Under
\emph{exact} reads the agent has the default read paths. Under
\emph{lagged} reads it can find issues only through Gitea's asynchronous
search index. Under \emph{look-alike}, a background actor creates an
effect with content identical to the focal write at the moment of the
fault, just before the focal request is sent, through the same service
account. This is the worst case for any check by content, since no
author or creation-time field separates the two records; a look-alike
from another account, or from an earlier time, could be told apart by
those fields. Under \emph{shared aggregate}, a background actor adjusts
the same counter or stock level.

\hypertarget{conditions}{%
\subsection*{7.3 Conditions}\label{conditions}}
\addcontentsline{toc}{subsection}{7.3 Conditions}

Every condition sits between the driver and the tools, and only the
recovery logic differs.

\begin{itemize}
\tightlist
\item
  \textbf{C0 plain:} the agent's own behaviour; for the scripted driver,
  one immediate re-issue on error.
\item
  \textbf{C1 never retry:} stop and report at the first failed write.
\item
  \textbf{C2 SDK retry:} transparent client retries, three with
  back-off.
\item
  \textbf{C3 idempotency keys:} a key per call, generated in memory and
  passed wherever the service honours one, which is common client
  practice {[}5{]}.
\item
  \textbf{C4 workflow engine:} each write is a Temporal activity with a
  retry policy of three attempts, and keys are derived from the durable
  workflow and activity ids {[}12{]}.
\item
  \textbf{C5 verify-before-retry:} re-issue only if a content-based
  postcondition fails {[}23{]}.
\item
  \textbf{C6 guard:} my re-implementation of Li's client-side guard from
  its published description {[}8{]}, with automatic keys, a re-check
  after the documented lag, a block on repeats it cannot verify, and an
  uncertainty annotation in the observation.
\item
  \textbf{Gw guard + wait:} C6, but waiting the same in-flight bound
  \(\Delta\) as Settle. This matched-latency control shows whether
  waiting alone explains Settle's results.
\item
  \textbf{C7 reflection:} after a failed tool call the model is asked to
  reason about the failure before acting {[}41{]}; language-model
  drivers only.
\item
  \textbf{C8 rollback:} on failure, compensate the completed steps from
  a log, then restart the workflow {[}14{]}, {[}15{]}.
\item
  \textbf{S Settle}, and its ablations S−j, S−t, S−h and S−c.
\item
  \textbf{O outcome oracle:} knows the hidden outcome that the fault
  mode is designed to produce for the focal write, including when an
  in-flight request should land, and re-issues only if the write is
  designed never to execute. It does not compensate. It shows what
  perfect knowledge of the focal outcome achieves \emph{without}
  compensation; it is not a lower bound for methods that compensate, and
  Settle does better than it where compensation matters.
\end{itemize}

Retry budgets are small and comparable. The scripted driver re-issues a
failed step once, C2 retries up to three times, and the guards and
Settle make at most three attempts per write. Language-model conditions
share the same model, prompts, and limits of 30 turns and 40 tool calls.
Wall-clock time is reported, not capped.

\textbf{Limiting bias towards my own method.} Because I designed Settle,
the theory and the harness, I took several precautions against a
comparison that favours Settle by construction. The predictions of the
achievability map that the hypotheses test, including where Settle
itself must fail, were frozen with the analysis code before the test
runs. Each baseline family is represented by its strongest configuration
that I could build without attribution: durable rather than in-memory
keys (C4), a guard that waits as long as Settle (Gw), and an outcome
oracle. Variants that \emph{do} use the tag but lack another component
are the ablations, and I report them next to the baselines so that the
contribution of attribution can be separated from that of the journal,
the horizon and compensation. Three of the eight test templates were not
used during development. The network fault pairs penalise eager and
reluctant re-issuing equally. And grading is done by code from the
services' own records, independently of every recovery layer. These
precautions do not remove the main limitation: the baselines are denied
tags by design, so part of the comparison confirms Proposition 4 rather
than testing an engineering choice. The re-implemented baselines may
also differ from their authors' code in detail; I release mine and
describe the behaviour of each.

\hypertarget{drivers}{%
\subsection*{7.4 Drivers}\label{drivers}}
\addcontentsline{toc}{subsection}{7.4 Drivers}

\textbf{Scripted driver.} The scripted driver is a deterministic agent
without a model. It executes the plan in order, keeps a durable step
checkpoint as agent frameworks do, re-issues a failed step once, finds
earlier results by content when it needs them, and reports honestly what
it believes. It isolates the effect of the recovery layer and makes the
large grid possible.

\textbf{Language-model drivers.} Three OpenAI models run, each at its
default reasoning setting, through the Responses API {[}42{]}:
\texttt{gpt-5.6-luna}, which reasons by default; \texttt{gpt-5.4-nano},
which does not; and \texttt{gpt-4.1-mini}, which is not a reasoning
model. The language-model study does not aim to rank models. It asks two
questions that the scripted driver cannot answer: what a model does on
its own after an ambiguous write, and whether Settle's guarantee
survives an agent that rephrases writes, deletes records, waits, and
crashes in the middle of a conversation. Each model receives a system
instruction that describes it as an autonomous operations agent; the
task goal, with the exact required strings and the sentence ``Each
action must take effect exactly once.''; and the tools of its task,
namely writes, reads, its own deletion tools, a wait tool and finish
(Appendix C). The agent ends by calling finish with a status, a summary
and a list of operations whose outcome it is unsure of. The conversation
is checkpointed after every tool result. After a crash, the run resumes
from the checkpoint and re-executes any pending tool call, which is the
replay behaviour of common agent frameworks.

The language-model grid uses the first two writes of every template on
test instance 100; the modes none, lost\_ack, not\_exec\_err, late,
late\_tail and crash\_after; all four read models where they apply; and
conditions C0, C5, C6, C7 and S. \texttt{gpt-5.6-luna} runs two repeats
and the other models one, for 2,680 planned episodes. Every model
response is cached by a hash of the request, so the released cache
replays every decision without an API key.

\hypertarget{grading-statistics-and-pre-registration}{%
\subsection*{7.5 Grading, statistics and
pre-registration}\label{grading-statistics-and-pre-registration}}
\addcontentsline{toc}{subsection}{7.5 Grading, statistics and
pre-registration}

\textbf{Grading.} Grading is done by code from the services' own
records, read with separate credentials after quiescence: PostgreSQL
audit triggers; the Redis stream's monotonic \texttt{entries-added}
counter and an atomic counter-audit list; TigerBeetle's immutable
transfer history; Mailpit's capture store, which no tool can delete
from; and Gitea state together with a compensation log written by the
harness. Background-actor effects are identified from the actor's own
log. Neither agents nor recovery layers can import the grader.

\textbf{Statistics.} Rates carry cluster-bootstrap 95\% intervals that
resample template × instance clusters 10,000 times. Paired comparisons
use the exact McNemar test with Holm correction within each hypothesis,
and equivalence uses two one-sided tests through a 90\% interval.
Because the scripted driver is deterministic and the instances of a
template differ mainly in their strings, worlds are not independent
draws, and the intervals measure replication across my design grid
rather than uncertainty about a population of deployments. As a
sensitivity check I also resample the eight templates as clusters
(Section 8.2). The rates themselves are averages over a designed grid of
worlds, so I report them per fault mode and per contract cell as well as
pooled. The six hypotheses, their thresholds and every analysis
deviation are listed in Appendix D, and Section 8 reports every
hypothesis, confirmed or not. Two analyses are exploratory and were not
pre-registered: the API census (Section 5) and the horizon trade-off
(Section 8.6). I label them as such wherever they appear.

\hypertarget{results}{%
\section*{8. Results}\label{results}}
\addcontentsline{toc}{section}{8. Results}

\hypertarget{the-test-runs}{%
\subsection*{8.1 The test runs}\label{the-test-runs}}
\addcontentsline{toc}{subsection}{8.1 The test runs}

The test split comprises 17,820 scripted and 2,680 language-model
episodes. Every scheduled episode finished, so no record is excluded.
The language-model test runs cost US\$ 3.24, and the main study as a
whole, including its pilot and development checks, cost US\$ 3.40. With
the follow-up study (US\$ 6.09, Section 9), all model calls of the
project cost US\$ 9.49. Two amendments were made after the freeze
(Appendix D). The first fixed a lock race in the cost ledger during the
runs and could not change any outcome. The second fixed a defect found
in the audit before writing: transfer keys chosen by the models collided
across episodes that replay the same task, so the ledger refused the
transfer while the tool reported success. All affected language-model
episodes were re-run under the amended code. The verdicts of H1--H4 and
H6 did not change, H5's cost criterion for \texttt{gpt-5.4-nano} moved
from inconclusive to confirmed, and Section 8.9 also reports the numbers
from before the amendment. Every condition left a consistent final state
in every one of the 1,305 fault-free scripted episodes, so the harness
and the grader create no inconsistency of their own.

\hypertarget{settle-converges-where-convergence-is-possible-h1}{%
\subsection*{8.2 Settle converges where convergence is possible
(H1)}\label{settle-converges-where-convergence-is-possible-h1}}
\addcontentsline{toc}{subsection}{8.2 Settle converges where convergence
is possible (H1)}

\begin{figure}
\centering
\includegraphics[width=0.53\textwidth,height=\textheight]{fig_h1_dots.png}
\caption{\textbf{Fig. 3.} Inconsistent final states for class-A,
targeted-reversible writes in ambiguous worlds (the H1 population;
scripted driver, all read models). Points are rates; bars are 95\%
cluster-bootstrap intervals over template × instance clusters.}
\end{figure}

Fig. 3 shows the H1 population: ambiguous worlds whose focal write is
taggable and reversible. Settle leaves 3.2\% of these episodes
inconsistent (95\% CI 2.2\%--4.1\%). The best baseline is Gw guard +
wait, the guard that waits as long as Settle, at 34.7\%, and the paired
difference is −31.5 pp (pp: percentage points; 95\% CI −32.5 pp to −30.5
pp; Holm-adjusted exact McNemar \(p\) \textless{} 0.001). H1 is
confirmed.

Plain retries, client retries, in-memory keys and the workflow engine
all end at about the same level: 60.7\%, 60.7\%, 60.5\% and 60.7\%. None
of them can tell whether an unacknowledged write took effect, and a key
or a durable workflow history does not help when the service ignores the
key. Verifying by content helps, but only partly: verify-before-retry
reaches 47.8\% and the guard 46.9\%. Waiting out the in-flight bound
lowers the guard to 34.7\%, so waiting alone does not explain the
difference between the guard and Settle. Rollback is the worst condition
at 62.5\%: it compensates and restarts on every error, including errors
after which the write \emph{had} taken effect, so the restarted workflow
duplicates it. The oracle knows every hidden outcome but cannot remove a
duplicate, and ends at 7.0\%. All its failures in this population are
transport redeliveries, which no knowledge of outcomes prevents and only
compensation repairs. This is why Settle does better than the oracle in
this population; without redelivery and partial batches the oracle is
perfect (Table 6).

\begin{longtable}[]{|
  >{\raggedright\arraybackslash}p{(\columnwidth - 2\tabcolsep - 4.1pt - 14\tabcolsep) * \real{0.2095}}|
  >{\raggedright\arraybackslash}p{(\columnwidth - 2\tabcolsep - 4.1pt - 14\tabcolsep) * \real{0.0381}}|
  >{\raggedright\arraybackslash}p{(\columnwidth - 2\tabcolsep - 4.1pt - 14\tabcolsep) * \real{0.1333}}|
  >{\raggedright\arraybackslash}p{(\columnwidth - 2\tabcolsep - 4.1pt - 14\tabcolsep) * \real{0.1905}}|
  >{\raggedright\arraybackslash}p{(\columnwidth - 2\tabcolsep - 4.1pt - 14\tabcolsep) * \real{0.0857}}|
  >{\raggedright\arraybackslash}p{(\columnwidth - 2\tabcolsep - 4.1pt - 14\tabcolsep) * \real{0.1143}}|
  >{\raggedright\arraybackslash}p{(\columnwidth - 2\tabcolsep - 4.1pt - 14\tabcolsep) * \real{0.1048}}|
  >{\raggedright\arraybackslash}p{(\columnwidth - 2\tabcolsep - 4.1pt - 14\tabcolsep) * \real{0.1238}}|}
\caption{\textbf{Table 5.} All ambiguous worlds by condition, scripted
driver. ``Duplicate executed'' counts episodes in which any step's
effect was executed more than once, \emph{including} duplicates that
were later compensated, and ``missing'' episodes with at least one
missing step. This is the prevention metric of prior work {[}8{]}. It
shows why prevention and convergence differ: Settle allows transient
duplicates and removes them.}\tabularnewline
\noalign{\global\arrayrulewidth=1.1pt}\arrayrulecolor{cvdeep}\hline\noalign{\global\arrayrulewidth=0.4pt}\arrayrulecolor{cvgrid}\rowcolor{cvhead}
\begin{minipage}[b]{\linewidth}\raggedright\bfseries
Condition
\end{minipage} & \begin{minipage}[b]{\linewidth}\raggedright\bfseries
n
\end{minipage} & \begin{minipage}[b]{\linewidth}\raggedright\bfseries
Inconsistent \%
\end{minipage} & \begin{minipage}[b]{\linewidth}\raggedright\bfseries
Duplicate executed \%
\end{minipage} & \begin{minipage}[b]{\linewidth}\raggedright\bfseries
Missing \%
\end{minipage} & \begin{minipage}[b]{\linewidth}\raggedright\bfseries
Collateral \%
\end{minipage} & \begin{minipage}[b]{\linewidth}\raggedright\bfseries
Overclaim \%
\end{minipage} & \begin{minipage}[b]{\linewidth}\raggedright\bfseries
Mean time (s)
\end{minipage} \\
\hline
\endfirsthead
\noalign{\global\arrayrulewidth=1.1pt}\arrayrulecolor{cvdeep}\hline\noalign{\global\arrayrulewidth=0.4pt}\arrayrulecolor{cvgrid}\rowcolor{cvhead}
\begin{minipage}[b]{\linewidth}\raggedright\bfseries
Condition
\end{minipage} & \begin{minipage}[b]{\linewidth}\raggedright\bfseries
n
\end{minipage} & \begin{minipage}[b]{\linewidth}\raggedright\bfseries
Inconsistent \%
\end{minipage} & \begin{minipage}[b]{\linewidth}\raggedright\bfseries
Duplicate executed \%
\end{minipage} & \begin{minipage}[b]{\linewidth}\raggedright\bfseries
Missing \%
\end{minipage} & \begin{minipage}[b]{\linewidth}\raggedright\bfseries
Collateral \%
\end{minipage} & \begin{minipage}[b]{\linewidth}\raggedright\bfseries
Overclaim \%
\end{minipage} & \begin{minipage}[b]{\linewidth}\raggedright\bfseries
Mean time (s)
\end{minipage} \\
\hline
\endhead
\endlastfoot
C0 plain & 1101 & 60.8 & 60.8 & 0.0 & 0.0 & 60.8 & 3.3 \\ \hline
C1 never retry & 1101 & 46.5 & 14.2 & 32.3 & 0.0 & 14.2 & 3.1 \\ \hline
C2 SDK retry & 1101 & 55.9 & 55.9 & 0.0 & 0.0 & 55.9 & 3.7 \\ \hline
C3 keys & 1101 & 55.8 & 55.8 & 0.0 & 0.0 & 55.7 & 3.3 \\ \hline
C4 workflow engine & 1101 & 55.0 & 55.0 & 0.0 & 0.0 & 55.0 & 7.4 \\ \hline
C5 verify-before-retry & 1101 & 46.6 & 32.5 & 14.2 & 0.0 & 46.6 & 3.6 \\ \hline
C6 guard & 1101 & 42.8 & 27.0 & 15.8 & 0.0 & 39.5 & 3.5 \\ \hline
Gw guard + wait & 1101 & 33.7 & 17.9 & 15.8 & 0.0 & 30.4 & 14.1 \\ \hline
C8 rollback & 1101 & 63.9 & 83.7 & 0.0 & 3.3 & 63.9 & 4.0 \\ \hline
S Settle & 1101 & 7.4 & 22.0 & 2.2 & 0.0 & 3.4 & 18.5 \\ \hline
S−j no journal & 1101 & 15.0 & 29.6 & 2.2 & 0.0 & 10.7 & 16.3 \\ \hline
S−t no tags & 1101 & 22.7 & 18.3 & 17.4 & 0.0 & 18.6 & 14.8 \\ \hline
S−h no horizon & 1101 & 23.9 & 27.1 & 2.6 & 0.0 & 19.7 & 4.8 \\ \hline
S−c no compensation & 1101 & 24.2 & 22.0 & 2.2 & 0.0 & 3.4 & 18.5 \\ \hline
O oracle & 1101 & 6.8 & 6.5 & 0.3 & 0.0 & 6.8 & 6.7 \\ \hline
\end{longtable}

Table 5 extends the comparison to all ambiguous worlds. Settle's rate of
duplicate executions (22.0\%) is higher than its inconsistency rate
(7.4\%); the difference is the duplicates that Settle let happen and
then removed. Under the prevention metric of prior work Settle looks
mediocre. Under convergence it is the only condition close to the
oracle.

\textbf{Where the difference comes from.} A pooled rate over a designed
grid depends on the mix of fault modes, so Table 6 breaks the H1
population down by mode and adds the ablations. Three of them (S−h, S−c,
S−j) use tags but lack another component; S−t keeps every component but
matches by content. Two modes favour Settle by construction: every
condition without a final sweep fails on redelivery, and every condition
that does not check rows individually fails on a partial batch. Without
these two modes, Settle leaves 3.4\% of 558 worlds inconsistent against
29.4\% for the waiting guard, and the oracle is perfect. The ablations
separate the components. A tag with a journal and compensation but
without the horizon (S−h) reaches 25.1\%, about the level of the waiting
guard. A tag with a horizon but without compensation (S−c) reaches
24.7\%, and a tag that is lost in a crash (S−j) reaches 11.5\%.
Attribution is therefore necessary but not sufficient: a tag-aware
variant that lacks the horizon or compensation is not much better than
the waiting guard, as Proposition 2 leads one to expect. These ablations
remove a component from Settle; they are not the strongest tag-based
alternative. Proposition 3 implies that a tag-based
\emph{wait-then-count} policy, which waits out \(\Delta\) before every
re-issue and never compensates, prevents duplicates whenever \(\Delta\)
holds. The follow-up study (Section 9) runs exactly this policy and
finds it identical to Settle on every ordinary fault, so Settle's
advantage over it lies mainly in transport redelivery. The comparison
with the waiting guard does not depend on the clustering. Resampling the
eight templates instead of the 24 template × instance clusters gives a
paired difference of −31.5 pp (95\% CI −33.2 pp to −29.6 pp), or −26.0
pp (−27.3 pp to −24.3 pp) without redelivery and partial batches, and
Settle is better than the waiting guard within every one of the eight
templates (a sign test over templates gives \(p\) = 0.008, two-sided).

\begin{longtable}[]{|
  >{\raggedright\arraybackslash}p{(\columnwidth - 2\tabcolsep - 5.7pt - 22\tabcolsep) * \real{0.3614}}|
  >{\raggedright\arraybackslash}p{(\columnwidth - 2\tabcolsep - 5.7pt - 22\tabcolsep) * \real{0.0723}}|
  >{\raggedright\arraybackslash}p{(\columnwidth - 2\tabcolsep - 5.7pt - 22\tabcolsep) * \real{0.0602}}|
  >{\raggedright\arraybackslash}p{(\columnwidth - 2\tabcolsep - 5.7pt - 22\tabcolsep) * \real{0.0602}}|
  >{\raggedright\arraybackslash}p{(\columnwidth - 2\tabcolsep - 5.7pt - 22\tabcolsep) * \real{0.0602}}|
  >{\raggedright\arraybackslash}p{(\columnwidth - 2\tabcolsep - 5.7pt - 22\tabcolsep) * \real{0.0602}}|
  >{\raggedright\arraybackslash}p{(\columnwidth - 2\tabcolsep - 5.7pt - 22\tabcolsep) * \real{0.0482}}|
  >{\raggedright\arraybackslash}p{(\columnwidth - 2\tabcolsep - 5.7pt - 22\tabcolsep) * \real{0.0602}}|
  >{\raggedright\arraybackslash}p{(\columnwidth - 2\tabcolsep - 5.7pt - 22\tabcolsep) * \real{0.0602}}|
  >{\raggedright\arraybackslash}p{(\columnwidth - 2\tabcolsep - 5.7pt - 22\tabcolsep) * \real{0.0482}}|
  >{\raggedright\arraybackslash}p{(\columnwidth - 2\tabcolsep - 5.7pt - 22\tabcolsep) * \real{0.0482}}|
  >{\raggedright\arraybackslash}p{(\columnwidth - 2\tabcolsep - 5.7pt - 22\tabcolsep) * \real{0.0602}}|}
\caption{\textbf{Table 6.} The H1 population by fault mode: inconsistent
final states (\%), scripted driver, all read models. S−h, S−c and S−j
use tags but lack the horizon, compensation and a crash-surviving
journal respectively; S−t has every component but matches by content
instead of by tag. The last row excludes the two modes that every
condition without a final sweep or row-level checks fails by
construction.}\tabularnewline
\noalign{\global\arrayrulewidth=1.1pt}\arrayrulecolor{cvdeep}\hline\noalign{\global\arrayrulewidth=0.4pt}\arrayrulecolor{cvgrid}\rowcolor{cvhead}
\begin{minipage}[b]{\linewidth}\raggedright\bfseries
Fault mode
\end{minipage} & \begin{minipage}[b]{\linewidth}\raggedright\bfseries
worlds
\end{minipage} & \begin{minipage}[b]{\linewidth}\raggedright\bfseries
C0
\end{minipage} & \begin{minipage}[b]{\linewidth}\raggedright\bfseries
C4
\end{minipage} & \begin{minipage}[b]{\linewidth}\raggedright\bfseries
C6
\end{minipage} & \begin{minipage}[b]{\linewidth}\raggedright\bfseries
Gw
\end{minipage} & \begin{minipage}[b]{\linewidth}\raggedright\bfseries
S−h
\end{minipage} & \begin{minipage}[b]{\linewidth}\raggedright\bfseries
S−c
\end{minipage} & \begin{minipage}[b]{\linewidth}\raggedright\bfseries
S−j
\end{minipage} & \begin{minipage}[b]{\linewidth}\raggedright\bfseries
S−t
\end{minipage} & \begin{minipage}[b]{\linewidth}\raggedright\bfseries
S
\end{minipage} & \begin{minipage}[b]{\linewidth}\raggedright\bfseries
O
\end{minipage} \\
\hline
\endfirsthead
\noalign{\global\arrayrulewidth=1.1pt}\arrayrulecolor{cvdeep}\hline\noalign{\global\arrayrulewidth=0.4pt}\arrayrulecolor{cvgrid}\rowcolor{cvhead}
\begin{minipage}[b]{\linewidth}\raggedright\bfseries
Fault mode
\end{minipage} & \begin{minipage}[b]{\linewidth}\raggedright\bfseries
worlds
\end{minipage} & \begin{minipage}[b]{\linewidth}\raggedright\bfseries
C0
\end{minipage} & \begin{minipage}[b]{\linewidth}\raggedright\bfseries
C4
\end{minipage} & \begin{minipage}[b]{\linewidth}\raggedright\bfseries
C6
\end{minipage} & \begin{minipage}[b]{\linewidth}\raggedright\bfseries
Gw
\end{minipage} & \begin{minipage}[b]{\linewidth}\raggedright\bfseries
S−h
\end{minipage} & \begin{minipage}[b]{\linewidth}\raggedright\bfseries
S−c
\end{minipage} & \begin{minipage}[b]{\linewidth}\raggedright\bfseries
S−j
\end{minipage} & \begin{minipage}[b]{\linewidth}\raggedright\bfseries
S−t
\end{minipage} & \begin{minipage}[b]{\linewidth}\raggedright\bfseries
S
\end{minipage} & \begin{minipage}[b]{\linewidth}\raggedright\bfseries
O
\end{minipage} \\
\hline
\endhead
\endlastfoot
lost\_ack & 96 & 100.0 & 100.0 & 7.3 & 6.3 & 0.0 & 0.0 & 0.0 & 43.8 &
0.0 & 0.0 \\ \hline
not\_exec\_err & 96 & 0.0 & 0.0 & 43.8 & 43.8 & 1.0 & 0.0 & 0.0 & 43.8 &
0.0 & 0.0 \\ \hline
held & 42 & 100.0 & 100.0 & 0.0 & 0.0 & 0.0 & 0.0 & 0.0 & 0.0 & 0.0 &
0.0 \\ \hline
not\_exec\_to & 96 & 0.0 & 0.0 & 43.8 & 43.8 & 4.2 & 0.0 & 0.0 & 43.8 &
0.0 & 0.0 \\ \hline
late & 96 & 100.0 & 100.0 & 62.5 & 6.3 & 99.0 & 100.0 & 0.0 & 6.3 & 0.0
& 0.0 \\ \hline
late\_tail & 42 & 100.0 & 100.0 & 100.0 & 54.8 & 95.2 & 100.0 & 45.2 &
45.2 & 45.2 & 0.0 \\ \hline
crash\_after & 45 & 100.0 & 100.0 & 100.0 & 100.0 & 0.0 & 0.0 & 100.0 &
0.0 & 0.0 & 0.0 \\ \hline
crash\_before & 45 & 0.0 & 0.0 & 0.0 & 0.0 & 0.0 & 0.0 & 0.0 & 0.0 & 0.0
& 0.0 \\ \hline
redelivery & 42 & 100.0 & 100.0 & 100.0 & 100.0 & 0.0 & 100.0 & 0.0 &
0.0 & 0.0 & 100.0 \\ \hline
partial & 3 & 100.0 & 100.0 & 100.0 & 100.0 & 0.0 & 0.0 & 0.0 & 0.0 &
0.0 & 0.0 \\ \hline
all but redelivery and partial & 558 & 57.5 & 57.5 & 42.7 & 29.4 & 25.1
& 24.7 & 11.5 & 27.1 & 3.4 & 0.0 \\ \hline
\end{longtable}

\hypertarget{where-keys-and-workflow-engines-are-enough-and-where-they-are-not-h3}{%
\subsection*{8.3 Where keys and workflow engines are enough, and where
they are not
(H3)}\label{where-keys-and-workflow-engines-are-enough-and-where-they-are-not-h3}}
\addcontentsline{toc}{subsection}{8.3 Where keys and workflow engines
are enough, and where they are not (H3)}

Table 7 fills in the achievability map with measured rates. On class-K
writes, where the service honours a key, Settle and the workflow engine
with durable keys are both exact: neither leaves any of the 108 paired
worlds inconsistent, so the equivalence margin is met trivially. H3 is
confirmed. The oracle's three failures in this cell (2.8\% of the 108
keyed worlds) are late-tail worlds in which the delayed TigerBeetle
transfer, injected to commit after 50--84 s, never executed; conditions
that re-issued with the same id ended consistent there. This is a fault
that did not realise its intended outcome, and it affects only the
oracle, which by design never re-issues a write it expects to land. Keys
generated in memory (C3) fail in exactly the crash worlds (37.5\%
inconsistent) and nowhere else (0.0\%), as Proposition 1 predicts.

On key-less writes the picture reverses. Keys and the workflow engine
are no better than plain retries, because there is nothing to send the
key to. Over all key-less ambiguous worlds their rates are 60.9\% and
61.0\%, against 61.0\% for plain retries. The workflow engine's durable
history does not stop a crashed step from taking effect twice: in
key-less crash-after-dispatch worlds it is inconsistent in 100.0\% of
episodes, against 0.0\% for Settle, because resuming the step means
executing the activity again (Proposition 6). In every key-less cell
Settle has the lowest rate of the non-oracle conditions or ties for it,
and its remaining failures in the attributable, reversible cells are all
delays beyond \(\Delta\).

\begin{longtable}[]{|
  >{\raggedright\arraybackslash}p{(\columnwidth - 2\tabcolsep - 5.3pt - 20\tabcolsep) * \real{0.4085}}|
  >{\raggedright\arraybackslash}p{(\columnwidth - 2\tabcolsep - 5.3pt - 20\tabcolsep) * \real{0.0845}}|
  >{\raggedright\arraybackslash}p{(\columnwidth - 2\tabcolsep - 5.3pt - 20\tabcolsep) * \real{0.0563}}|
  >{\raggedright\arraybackslash}p{(\columnwidth - 2\tabcolsep - 5.3pt - 20\tabcolsep) * \real{0.0563}}|
  >{\raggedright\arraybackslash}p{(\columnwidth - 2\tabcolsep - 5.3pt - 20\tabcolsep) * \real{0.0563}}|
  >{\raggedright\arraybackslash}p{(\columnwidth - 2\tabcolsep - 5.3pt - 20\tabcolsep) * \real{0.0563}}|
  >{\raggedright\arraybackslash}p{(\columnwidth - 2\tabcolsep - 5.3pt - 20\tabcolsep) * \real{0.0563}}|
  >{\raggedright\arraybackslash}p{(\columnwidth - 2\tabcolsep - 5.3pt - 20\tabcolsep) * \real{0.0563}}|
  >{\raggedright\arraybackslash}p{(\columnwidth - 2\tabcolsep - 5.3pt - 20\tabcolsep) * \real{0.0563}}|
  >{\raggedright\arraybackslash}p{(\columnwidth - 2\tabcolsep - 5.3pt - 20\tabcolsep) * \real{0.0563}}|
  >{\raggedright\arraybackslash}p{(\columnwidth - 2\tabcolsep - 5.3pt - 20\tabcolsep) * \real{0.0563}}|}
\caption{\textbf{Table 7.} The achievability map, measured: inconsistent
final states (\%) by contract cell of the focal write and condition, in
ambiguous worlds (scripted driver). ``Worlds'' is the number of
ambiguous worlds per condition.}\tabularnewline
\noalign{\global\arrayrulewidth=1.1pt}\arrayrulecolor{cvdeep}\hline\noalign{\global\arrayrulewidth=0.4pt}\arrayrulecolor{cvgrid}\rowcolor{cvhead}
\begin{minipage}[b]{\linewidth}\raggedright\bfseries
Contract cell (focal write)
\end{minipage} & \begin{minipage}[b]{\linewidth}\raggedright\bfseries
worlds
\end{minipage} & \begin{minipage}[b]{\linewidth}\raggedright\bfseries
C0
\end{minipage} & \begin{minipage}[b]{\linewidth}\raggedright\bfseries
C3
\end{minipage} & \begin{minipage}[b]{\linewidth}\raggedright\bfseries
C4
\end{minipage} & \begin{minipage}[b]{\linewidth}\raggedright\bfseries
C5
\end{minipage} & \begin{minipage}[b]{\linewidth}\raggedright\bfseries
C6
\end{minipage} & \begin{minipage}[b]{\linewidth}\raggedright\bfseries
Gw
\end{minipage} & \begin{minipage}[b]{\linewidth}\raggedright\bfseries
C8
\end{minipage} & \begin{minipage}[b]{\linewidth}\raggedright\bfseries
S
\end{minipage} & \begin{minipage}[b]{\linewidth}\raggedright\bfseries
O
\end{minipage} \\
\hline
\endfirsthead
\noalign{\global\arrayrulewidth=1.1pt}\arrayrulecolor{cvdeep}\hline\noalign{\global\arrayrulewidth=0.4pt}\arrayrulecolor{cvgrid}\rowcolor{cvhead}
\begin{minipage}[b]{\linewidth}\raggedright\bfseries
Contract cell (focal write)
\end{minipage} & \begin{minipage}[b]{\linewidth}\raggedright\bfseries
worlds
\end{minipage} & \begin{minipage}[b]{\linewidth}\raggedright\bfseries
C0
\end{minipage} & \begin{minipage}[b]{\linewidth}\raggedright\bfseries
C3
\end{minipage} & \begin{minipage}[b]{\linewidth}\raggedright\bfseries
C4
\end{minipage} & \begin{minipage}[b]{\linewidth}\raggedright\bfseries
C5
\end{minipage} & \begin{minipage}[b]{\linewidth}\raggedright\bfseries
C6
\end{minipage} & \begin{minipage}[b]{\linewidth}\raggedright\bfseries
Gw
\end{minipage} & \begin{minipage}[b]{\linewidth}\raggedright\bfseries
C8
\end{minipage} & \begin{minipage}[b]{\linewidth}\raggedright\bfseries
S
\end{minipage} & \begin{minipage}[b]{\linewidth}\raggedright\bfseries
O
\end{minipage} \\
\hline
\endhead
\endlastfoot
K, keyed & 108 & 58.3 & 8.3 & 0.0 & 41.7 & 8.3 & 8.3 & 58.3 & 0.0 &
2.8 \\ \hline
A, reversible, exact read & 387 & 66.7 & 66.4 & 66.7 & 45.0 & 45.2 &
29.2 & 60.2 & 4.9 & 10.9 \\ \hline
A, reversible, lagged read & 48 & 50.0 & 50.0 & 50.0 & 37.5 & 25.0 & 0.0
& 62.5 & 0.0 & 0.0 \\ \hline
A, reversible, look-alikes & 168 & 50.0 & 50.0 & 50.0 & 57.1 & 57.1 &
57.1 & 67.9 & 0.0 & 0.0 \\ \hline
A, irreversible (mail) & 234 & 61.5 & 61.5 & 61.5 & 46.2 & 45.7 & 36.3 &
71.8 & 12.8 & 7.7 \\ \hline
C, content only, irreversible & 39 & 61.5 & 61.5 & 61.5 & 46.2 & 46.2 &
35.9 & 61.5 & 28.2 & 7.7 \\ \hline
U, exclusive aggregate & 81 & 66.7 & 66.7 & 66.7 & 44.4 & 44.4 & 44.4 &
66.7 & 3.7 & 11.1 \\ \hline
U, shared aggregate & 36 & 50.0 & 50.0 & 50.0 & 50.0 & 50.0 & 50.0 &
50.0 & 50.0 & 0.0 \\ \hline
\end{longtable}

\hypertarget{content-based-verification-under-look-alikes-h2}{%
\subsection*{8.4 Content-based verification under look-alikes
(H2)}\label{content-based-verification-under-look-alikes-h2}}
\addcontentsline{toc}{subsection}{8.4 Content-based verification under
look-alikes (H2)}

When another user of the same account creates a look-alike, every method
that recognises its own effects by content fails in more than half of
the 252 ambiguous look-alike worlds of all contract classes (Table 7
restricts them to class-A reversible writes, 168 worlds):
verify-before-retry in 54.8\%, the guard in 54.8\%, and Settle without
tags in 69.0\%. When the agent's write never executed, these methods
take the other user's record for their own and skip the write, so their
own effect is missing. Settle's rate is 2.4\%, and H2 is confirmed.
Rollback is the only baseline that deletes, and it causes collateral
damage in 14.3\% of these worlds. Every case is a world whose focal
write was a comment: after the comment failed, rollback compensated the
completed steps by deleting the agent's own issue, and with it the other
user's comment on that issue. Compensation that is correct for the
agent's own effect still destroys what others attached to it.

\hypertarget{what-no-method-can-do-and-what-settle-does-instead-h4}{%
\subsection*{8.5 What no method can do, and what Settle does instead
(H4)}\label{what-no-method-can-do-and-what-settle-does-instead-h4}}
\addcontentsline{toc}{subsection}{8.5 What no method can do, and what
Settle does instead (H4)}

When another writer adjusts the same counter or stock level, every
non-oracle condition is inconsistent in at least 50.0\% of ambiguous
worlds (H4a confirmed). Settle abstains there: it returns \emph{unknown}
instead of guessing. Across shared-aggregate and irreversible worlds,
Settle overclaims in 4.9\% of its episodes (15 of 309; H4b confirmed),
and it reports a problem in 44 of the 59 episodes that ended
inconsistent. Plain retries, by comparison, overclaim in 60.8\% of all
ambiguous worlds. In irreversible worlds every condition, the oracle
included, shows duplicates or missing effects (H4c confirmed). H4a and
H4c, like the lenient thresholds of H5 below, confirm properties that
the design of the fault grid makes very likely: half of each
shared-aggregate pair is a never-executed world, and redelivery of an
irreversible write cannot be undone. They check that the harness behaves
as the theory says rather than testing Settle.

\textbf{Every one of Settle's failures is accounted for by the theory.}
Of Settle's 1,101 ambiguous scripted episodes, 81 ended inconsistent
(Table 8). Each is either a violation of the in-flight bound that Settle
assumes, whose consequence Propositions 2 and 7 bound, or falls into a
cell where the theory admits no client-side guarantee. The first kind is
not a prediction in the strong sense; it is a failure of an assumption,
and the theory says how likely it is.

\begin{longtable}[]{|
  >{\raggedright\arraybackslash}p{(\columnwidth - 2\tabcolsep - 2.1pt - 4\tabcolsep) * \real{0.6522}}|
  >{\raggedright\arraybackslash}p{(\columnwidth - 2\tabcolsep - 2.1pt - 4\tabcolsep) * \real{0.1630}}|
  >{\raggedright\arraybackslash}p{(\columnwidth - 2\tabcolsep - 2.1pt - 4\tabcolsep) * \real{0.1848}}|}
\caption{\textbf{Table 8.} Every inconsistent Settle episode (scripted
driver), by the theoretical reason it could not converge, and how many
Settle itself reported as uncertain or as an unrecoverable
duplicate.}\tabularnewline
\noalign{\global\arrayrulewidth=1.1pt}\arrayrulecolor{cvdeep}\hline\noalign{\global\arrayrulewidth=0.4pt}\arrayrulecolor{cvgrid}\rowcolor{cvhead}
\begin{minipage}[b]{\linewidth}\raggedright\bfseries
Cause (proposition)
\end{minipage} & \begin{minipage}[b]{\linewidth}\raggedright\bfseries
Settle failures
\end{minipage} & \begin{minipage}[b]{\linewidth}\raggedright\bfseries
of which reported
\end{minipage} \\
\hline
\endfirsthead
\noalign{\global\arrayrulewidth=1.1pt}\arrayrulecolor{cvdeep}\hline\noalign{\global\arrayrulewidth=0.4pt}\arrayrulecolor{cvgrid}\rowcolor{cvhead}
\begin{minipage}[b]{\linewidth}\raggedright\bfseries
Cause (proposition)
\end{minipage} & \begin{minipage}[b]{\linewidth}\raggedright\bfseries
Settle failures
\end{minipage} & \begin{minipage}[b]{\linewidth}\raggedright\bfseries
of which reported
\end{minipage} \\
\hline
\endhead
\endlastfoot
assumed in-flight bound violated: delay longer than \(\Delta\) (risk
bounded by Propositions 2 and 7) & 36 & 5 \\ \hline
transport redelivery of an irreversible key-less write (no client policy
can prevent it; Proposition 3) & 21 & 21 \\ \hline
content-only write under a look-alike, not attributable (Proposition 4)
& 6 & 0 \\ \hline
shared aggregate, not attributable (Proposition 5) & 18 & 18 \\ \hline
anything else & 0 & -- \\ \hline
\end{longtable}

The late-tail worlds test the qualitative prediction of Proposition 2
that Settle can fail only when the in-flight bound is violated. Of 84
late-tail worlds, 48 had a delay above the assumed bound of 20 s. Settle
failed in 36 of those and in 0 of the worlds whose delay was within the
bound; the shortest delay among the failures was 21.2 s. Of the twelve
worlds above the bound in which Settle did not fail, seven had a keyed
focal write, which the bound does not affect, and in the other five the
late commit still landed before Settle's final sweep at the end of the
episode. Within the late-tail worlds Settle therefore fails only when
its assumed bound is wrong (36 of the 48 worlds above the bound, none of
the others). Of the 36 delay failures, 22 concern reversible and 14
irreversible writes. Settle's unreported failures are all of two kinds:
31 delay failures and the 6 content-only failures. The content-only
failures are a gap in my implementation rather than in the theory:
Settle resolves an ambiguous content-only write by content (Section 6),
and a look-alike satisfies the check. A Settle that never confirms a
content-only write by content reports them all, as the follow-up study
confirms (Section 9). When a late commit lands after the horizon, Settle
has already confirmed the operation, so only 5 of the 36 delay failures
were flagged. In contrast, every shared-aggregate abstention (18 of 18)
and every unrecoverable duplicate (21 of 21) was reported.

\hypertarget{how-the-remaining-risk-trades-against-latency}{%
\subsection*{8.6 How the remaining risk trades against
latency}\label{how-the-remaining-risk-trades-against-latency}}
\addcontentsline{toc}{subsection}{8.6 How the remaining risk trades
against latency}

Settle has one parameter that the theory does not fix, the assumed
in-flight bound \(\Delta\). A larger bound buys a lower residual risk at
the cost of latency, and Proposition 7 says that no smarter rule escapes
this trade-off at a fixed decision time. This analysis is exploratory.
To trace the trade-off without new runs, I use a property of Settle's
decision rule that the recorded episodes let me check. In a late-commit
world, Settle ends inconsistent if and only if the commit lands after
its last decision about the operation. For an irreversible write that
decision is taken \(\Delta\) after the last attempt, which in these
worlds is the first; for a reversible one it is the final sweep at the
end of the episode. At the frozen \(\Delta\) = 20 s, this rule
reproduces the measured outcome in 167 of the 168 recorded late and
late-tail worlds with exact reads, which is the population I replay. (It
does not apply to shared-aggregate worlds, where Settle abstains.) Each
recorded world therefore also tells me its outcome under a different
bound: the decision time shifts by the difference in \(\Delta\), while
the recorded delay of the commit stays the same. The replay assumes that
Settle's other actions do not change with \(\Delta\), which holds for
its decision rule but was verified only at the frozen value.

\begin{figure}
\centering
\includegraphics[width=0.68\textwidth,height=\textheight]{fig_horizon.png}
\caption{\textbf{Fig. 4.} Exploratory replay of the recorded Settle
episodes under other assumed in-flight bounds \(\Delta\) (scripted
driver). Left: inconsistent final states in the late-commit worlds with
exact reads, for the fixed 12 s delay and the heavy-tailed 5--120 s
delay. Right: median duration of Settle episodes in the lost-reply,
timeout and late-commit modes. At the frozen \(\Delta\) = 20 s the
replay reproduces the measured outcome in all but one of these 168
worlds.}
\end{figure}

Fig. 4 shows the result. With the fixed 12 s delay, a bound shorter than
the delay is a cliff: at \(\Delta\) = 10 s, 25\% of these worlds end
inconsistent, and from \(\Delta\) = 12.5 s on none does. With the
heavy-tailed delay the trade-off is gradual. The replayed inconsistency
falls from 55\% at 10 s through 42\% at 20 s (measured: 36 of 84, one
world more than the replay) and 20\% at 60 s to 0\% at 120 s, the upper
end of the distribution I chose. Over the same range the median duration
of an episode in the lost-reply, timeout and late-commit modes grows
from 15.0 s to 125.0 s, against 5.0 s for plain retries; only the
episodes that actually waited out the bound move with it. The curve is
the tail of the delay distribution seen through Settle's decision times,
exactly as Propositions 2 and 7 predict. Two consequences follow. A
deployment should set \(\Delta\) from the delay tail of its own
services. Settle's journal records every late commit that a sweep
catches, but those observations are censored at the horizon, so the part
of the tail that matters most is visible only if reconciliation
continues beyond it. And since the horizon of a reversible write can
continue after the agent has reported, the latency cost needs to be paid
in full only for irreversible writes (Section 10).

\hypertarget{each-component-is-needed-where-the-theory-says-ablations}{%
\subsection*{8.7 Each component is needed where the theory says
(ablations)}\label{each-component-is-needed-where-the-theory-says-ablations}}
\addcontentsline{toc}{subsection}{8.7 Each component is needed where the
theory says (ablations)}

\begin{longtable}[]{|
  >{\raggedright\arraybackslash}p{(\columnwidth - 2\tabcolsep - 3.7pt - 12\tabcolsep) * \real{0.1881}}|
  >{\raggedright\arraybackslash}p{(\columnwidth - 2\tabcolsep - 3.7pt - 12\tabcolsep) * \real{0.1683}}|
  >{\raggedright\arraybackslash}p{(\columnwidth - 2\tabcolsep - 3.7pt - 12\tabcolsep) * \real{0.1287}}|
  >{\raggedright\arraybackslash}p{(\columnwidth - 2\tabcolsep - 3.7pt - 12\tabcolsep) * \real{0.1386}}|
  >{\raggedright\arraybackslash}p{(\columnwidth - 2\tabcolsep - 3.7pt - 12\tabcolsep) * \real{0.1485}}|
  >{\raggedright\arraybackslash}p{(\columnwidth - 2\tabcolsep - 3.7pt - 12\tabcolsep) * \real{0.1386}}|
  >{\raggedright\arraybackslash}p{(\columnwidth - 2\tabcolsep - 3.7pt - 12\tabcolsep) * \real{0.0891}}|}
\caption{\textbf{Table 9.} Ablations. Δ is the ablation's inconsistency
minus Settle's, paired by world, in the cell where the removed component
is needed and in all other ambiguous worlds.}\tabularnewline
\noalign{\global\arrayrulewidth=1.1pt}\arrayrulecolor{cvdeep}\hline\noalign{\global\arrayrulewidth=0.4pt}\arrayrulecolor{cvgrid}\rowcolor{cvhead}
\begin{minipage}[b]{\linewidth}\raggedright\bfseries
Ablation
\end{minipage} & \begin{minipage}[b]{\linewidth}\raggedright\bfseries
Component removed
\end{minipage} & \begin{minipage}[b]{\linewidth}\raggedright\bfseries
Pairs in cell
\end{minipage} & \begin{minipage}[b]{\linewidth}\raggedright\bfseries
Δ in cell (pp)
\end{minipage} & \begin{minipage}[b]{\linewidth}\raggedright\bfseries
95\% CI
\end{minipage} & \begin{minipage}[b]{\linewidth}\raggedright\bfseries
Δ outside (pp)
\end{minipage} & \begin{minipage}[b]{\linewidth}\raggedright\bfseries
Verdict
\end{minipage} \\
\hline
\endfirsthead
\noalign{\global\arrayrulewidth=1.1pt}\arrayrulecolor{cvdeep}\hline\noalign{\global\arrayrulewidth=0.4pt}\arrayrulecolor{cvgrid}\rowcolor{cvhead}
\begin{minipage}[b]{\linewidth}\raggedright\bfseries
Ablation
\end{minipage} & \begin{minipage}[b]{\linewidth}\raggedright\bfseries
Component removed
\end{minipage} & \begin{minipage}[b]{\linewidth}\raggedright\bfseries
Pairs in cell
\end{minipage} & \begin{minipage}[b]{\linewidth}\raggedright\bfseries
Δ in cell (pp)
\end{minipage} & \begin{minipage}[b]{\linewidth}\raggedright\bfseries
95\% CI
\end{minipage} & \begin{minipage}[b]{\linewidth}\raggedright\bfseries
Δ outside (pp)
\end{minipage} & \begin{minipage}[b]{\linewidth}\raggedright\bfseries
Verdict
\end{minipage} \\
\hline
\endhead
\endlastfoot
S−j no journal & journal (P6) & 87 & +96.6 & {[}+92.6, +100.0{]} & +0.0
& confirmed \\ \hline
S−h no horizon & horizon (P2) & 252 & +70.2 & {[}+64.7, +75.5{]} & +0.6
& confirmed \\ \hline
S−t no tags & tags (P4) & 462 & +36.4 & {[}+32.3, +40.3{]} & +0.2 &
confirmed \\ \hline
S−c no compensation & compensation (P2) & 216 & +74.5 & {[}+65.4,
+82.4{]} & +2.7 & confirmed \\ \hline
\end{longtable}

Each ablation fails in its predicted cell and almost nowhere else (Table
9, Fig. 5). Without the journal, Settle fails after a crash. Without the
horizon, it fails on late commits. Without tags, it fails under
look-alikes, where it mistakes the other user's record for its own. And
without compensation, it leaves standing the duplicates it detected.
Outside its cell, no ablation is more than three points worse than
Settle, although the small excess of S−c (+2.7 points) is itself
significant.

\begin{figure}
\centering
\includegraphics[width=0.6\textwidth,height=\textheight]{fig_modes_heatmap.png}
\caption{\textbf{Fig. 5.} Inconsistent final states (\%) by fault mode
and condition, scripted driver, all read models. Each network pair
(lost\_ack/not\_exec\_err, held/not\_exec\_to) shows a method's
trade-off: re-issuing fails in the first mode, not re-issuing in the
second. The crash pair (crash\_after/crash\_before) penalises only
re-issuing, because after a crash before dispatch every condition
re-issues the step. Settle and the oracle are low on both sides of every
pair. Each ablation fails almost only where its removed component is
needed.}
\end{figure}

\hypertarget{how-language-model-agents-treat-an-ambiguous-outcome}{%
\subsection*{8.8 How language-model agents treat an ambiguous
outcome}\label{how-language-model-agents-treat-an-ambiguous-outcome}}
\addcontentsline{toc}{subsection}{8.8 How language-model agents treat an
ambiguous outcome}

Without a recovery layer, the three models split into two behaviours
after their first failed or timed-out write (Table 10, Fig. 6).
\texttt{gpt-4.1-mini} and \texttt{gpt-5.4-nano} re-issue the write
immediately in 80\% and 68\% of these episodes. \texttt{gpt-5.6-luna},
the only model that reasons by default, re-issues at once in only 17\%,
and its first action is a read in 69\%. The hidden truth changes what
the models do only modestly. The paired worlds give identical tool
observations, so the remaining differences reflect sampling and the
small differences between the episodes' histories, such as ids and
timings, rather than information about the outcome. The immediate
re-issue rate is 16\% in committed and 17\% in never-executed worlds for
\texttt{gpt-5.6-luna}, 77\% and 90\% for \texttt{gpt-4.1-mini}, and 70\%
and 62\% for \texttt{gpt-5.4-nano}.

Checking first helps, but only as far as Propositions 4 and 6 allow.
Without a layer, \texttt{gpt-5.6-luna} ends inconsistent in 32\% of
these episodes, against 70\% for \texttt{gpt-4.1-mini}, and its failures
concentrate exactly where checking cannot work. In look-alike worlds
(70\% inconsistent) the other user's identical record convinces it that
its own write took effect. In crash worlds (66\%) the framework replays
the pending call after the restart, before the model can intervene.
Under heavy-tailed delays (33\%) the record is not there yet when the
model looks. With exact reads and a short delay, the same model fails in
only 3\% of the worlds with a lost reply. In the look-alike and
late-commit cells the model's observation is insufficient: a more
capable model would read the same look-alike and the same empty result,
unless it used fields my look-alikes do not distinguish, such as author
and creation time, or waited longer than my late commits. Crash replay
is different: it is a framework behaviour that no model can influence.

Asking the model to reflect on the error (C7) gave no benefit in hard
worlds: it was two to six points worse than the plain agent for every
model (Table 11). In lost-reply worlds with exact reads it moved
\texttt{gpt-4.1-mini} from 80\% to 47\% and \texttt{gpt-5.4-nano} from
40\% to 73\%, but with fifteen episodes per cell these differences are
within sampling noise. Reports were often not honest about the
uncertainty. After an ambiguous write, \texttt{gpt-4.1-mini} listed an
operation as uncertain in 1\% of episodes and claimed completion in
96\%; \texttt{gpt-5.6-luna} listed one in 22\% and claimed completion in
80\%; \texttt{gpt-5.4-nano} listed one in 49\%.

\begin{longtable}[]{|
  >{\raggedright\arraybackslash}p{(\columnwidth - 2\tabcolsep - 4.9pt - 18\tabcolsep) * \real{0.0960}}|
  >{\raggedright\arraybackslash}p{(\columnwidth - 2\tabcolsep - 4.9pt - 18\tabcolsep) * \real{0.0320}}|
  >{\raggedright\arraybackslash}p{(\columnwidth - 2\tabcolsep - 4.9pt - 18\tabcolsep) * \real{0.0800}}|
  >{\raggedright\arraybackslash}p{(\columnwidth - 2\tabcolsep - 4.9pt - 18\tabcolsep) * \real{0.0560}}|
  >{\raggedright\arraybackslash}p{(\columnwidth - 2\tabcolsep - 4.9pt - 18\tabcolsep) * \real{0.0480}}|
  >{\raggedright\arraybackslash}p{(\columnwidth - 2\tabcolsep - 4.9pt - 18\tabcolsep) * \real{0.0560}}|
  >{\raggedright\arraybackslash}p{(\columnwidth - 2\tabcolsep - 4.9pt - 18\tabcolsep) * \real{0.2000}}|
  >{\raggedright\arraybackslash}p{(\columnwidth - 2\tabcolsep - 4.9pt - 18\tabcolsep) * \real{0.1680}}|
  >{\raggedright\arraybackslash}p{(\columnwidth - 2\tabcolsep - 4.9pt - 18\tabcolsep) * \real{0.1520}}|
  >{\raggedright\arraybackslash}p{(\columnwidth - 2\tabcolsep - 4.9pt - 18\tabcolsep) * \real{0.1120}}|}
\caption{\textbf{Table 10.} What a language-model agent does after the
first failed or timed-out write when no recovery layer is present
(condition C0): its first action, whether it checked state before
re-issuing, and what it reported. ``Inconsistent'' is the share of these
episodes whose final state was wrong.}\tabularnewline
\noalign{\global\arrayrulewidth=1.1pt}\arrayrulecolor{cvdeep}\hline\noalign{\global\arrayrulewidth=0.4pt}\arrayrulecolor{cvgrid}\rowcolor{cvhead}
\begin{minipage}[b]{\linewidth}\raggedright\bfseries
Model
\end{minipage} & \begin{minipage}[b]{\linewidth}\raggedright\bfseries
n
\end{minipage} & \begin{minipage}[b]{\linewidth}\raggedright\bfseries
Re-issue \%
\end{minipage} & \begin{minipage}[b]{\linewidth}\raggedright\bfseries
Check \%
\end{minipage} & \begin{minipage}[b]{\linewidth}\raggedright\bfseries
Wait \%
\end{minipage} & \begin{minipage}[b]{\linewidth}\raggedright\bfseries
Other \%
\end{minipage} & \begin{minipage}[b]{\linewidth}\raggedright\bfseries
Checked before re-issue \%
\end{minipage} & \begin{minipage}[b]{\linewidth}\raggedright\bfseries
Listed as uncertain \%
\end{minipage} & \begin{minipage}[b]{\linewidth}\raggedright\bfseries
Claimed completed \%
\end{minipage} & \begin{minipage}[b]{\linewidth}\raggedright\bfseries
Inconsistent \%
\end{minipage} \\
\hline
\endfirsthead
\noalign{\global\arrayrulewidth=1.1pt}\arrayrulecolor{cvdeep}\hline\noalign{\global\arrayrulewidth=0.4pt}\arrayrulecolor{cvgrid}\rowcolor{cvhead}
\begin{minipage}[b]{\linewidth}\raggedright\bfseries
Model
\end{minipage} & \begin{minipage}[b]{\linewidth}\raggedright\bfseries
n
\end{minipage} & \begin{minipage}[b]{\linewidth}\raggedright\bfseries
Re-issue \%
\end{minipage} & \begin{minipage}[b]{\linewidth}\raggedright\bfseries
Check \%
\end{minipage} & \begin{minipage}[b]{\linewidth}\raggedright\bfseries
Wait \%
\end{minipage} & \begin{minipage}[b]{\linewidth}\raggedright\bfseries
Other \%
\end{minipage} & \begin{minipage}[b]{\linewidth}\raggedright\bfseries
Checked before re-issue \%
\end{minipage} & \begin{minipage}[b]{\linewidth}\raggedright\bfseries
Listed as uncertain \%
\end{minipage} & \begin{minipage}[b]{\linewidth}\raggedright\bfseries
Claimed completed \%
\end{minipage} & \begin{minipage}[b]{\linewidth}\raggedright\bfseries
Inconsistent \%
\end{minipage} \\
\hline
\endhead
\endlastfoot
gpt-4.1-mini & 102 & 80 & 20 & 0 & 0 & 16.3 & 1.0 & 96.1 & 69.6 \\ \hline
gpt-5.6-luna & 204 & 17 & 69 & 7 & 8 & 72.1 & 22.1 & 80.4 & 32.4 \\ \hline
gpt-5.4-nano & 102 & 68 & 22 & 0 & 11 & 20.2 & 49.0 & 64.7 & 53.9 \\ \hline
\end{longtable}

\begin{figure}
\centering
\includegraphics[width=0.49\textwidth,height=\textheight]{fig_rq1_behaviour.png}
\caption{\textbf{Fig. 6.} First action of each model after an ambiguous
write result, without a recovery layer (C0).}
\end{figure}

\hypertarget{settle-under-language-model-agents-h6}{%
\subsection*{8.9 Settle under language-model agents
(H6)}\label{settle-under-language-model-agents-h6}}
\addcontentsline{toc}{subsection}{8.9 Settle under language-model agents
(H6)}

\begin{figure}
\centering
\includegraphics[width=0.54\textwidth,height=\textheight]{fig_llm_dots.png}
\caption{\textbf{Fig. 7.} Inconsistent final states per model and
condition in hard worlds (late\_tail, crash\_after, look-alike or lagged
reads) and easy worlds (exact reads; lost\_ack, not\_exec\_err, late).
95\% cluster-bootstrap intervals.}
\end{figure}

With Settle underneath, the differences between models disappear (Fig.
7, Table 11). In hard worlds, which are the late-tail, crash and
look-alike or lagged-read worlds, each of the three models ends
inconsistent in 6.0\% of episodes, while every baseline condition of
every model lies between 47.8\% and 83.6\%. Hard and easy worlds as
defined before the test runs exclude the shared-aggregate worlds; they
hold 8 of Settle's language-model failures, all of them abstentions that
Settle reported. Every comparison is significant (largest Holm-adjusted
\(p\) \textless{} 0.001), so H6 is confirmed for every model. The three
Settle rates are identical because the remaining failures are determined
by the world, not by the model. Of 472 ambiguous Settle episodes, 24
ended inconsistent: 16 with a delay beyond the assumed bound and 8 on
shared aggregates, the same predicted kinds as with the scripted driver,
and none for any other reason. In easy worlds no Settle episode of any
model ended inconsistent, whereas even the careful \texttt{gpt-5.6-luna}
loses something there on its own (12.2\%).

Identical rates across models are what one should expect when the layer
resolves every write before the model sees an error, so they are not by
themselves evidence of robustness. A recovery guarantee that depends on
how carefully a model reasons would not be a guarantee, and the scripted
results already show that Settle's rests on the contract. What the
language-model runs add is a check that the guarantee survives the
behaviours a model brings and a scripted agent does not: rephrased
re-issues, waits and deletions. The main runs bound writes to intents
through the harness's step mapping (Section 6); the follow-up repeats
them with plan-free binding and gets identical outcomes, and adds a
frontier model (Section 9). No agent under Settle removed another user's
record, against 6 episodes with such collateral damage under the
baseline conditions.

\textbf{Effect of amendment 2.} Before amendment 2, the ledger-key
collision had inflated the baselines' rates for the two models that
choose their own transfer keys. For \texttt{gpt-5.6-luna} without a
layer, the hard-world rate moved from 65.7\% to 54.5\% and the
easy-world rate from 37.8\% to 12.2\%. For \texttt{gpt-5.4-nano} the
hard-world rate moved from 74.6\% to 62.7\% and the easy-world rate from
68.9\% to 44.4\%. H6 was confirmed for both models before and after the
amendment, and every number reported here is the corrected one.

\begin{longtable}[]{|
  >{\raggedright\arraybackslash}p{(\columnwidth - 2\tabcolsep - 3.7pt - 12\tabcolsep) * \real{0.0938}}|
  >{\raggedright\arraybackslash}p{(\columnwidth - 2\tabcolsep - 3.7pt - 12\tabcolsep) * \real{0.1719}}|
  >{\raggedright\arraybackslash}p{(\columnwidth - 2\tabcolsep - 3.7pt - 12\tabcolsep) * \real{0.2109}}|
  >{\raggedright\arraybackslash}p{(\columnwidth - 2\tabcolsep - 3.7pt - 12\tabcolsep) * \real{0.2109}}|
  >{\raggedright\arraybackslash}p{(\columnwidth - 2\tabcolsep - 3.7pt - 12\tabcolsep) * \real{0.1797}}|
  >{\raggedright\arraybackslash}p{(\columnwidth - 2\tabcolsep - 3.7pt - 12\tabcolsep) * \real{0.0391}}|
  >{\raggedright\arraybackslash}p{(\columnwidth - 2\tabcolsep - 3.7pt - 12\tabcolsep) * \real{0.0938}}|}
\caption{\textbf{Table 11.} Language-model agents on the test split:
inconsistent final states in hard and easy ambiguous worlds;
overclaiming, as a share of all episodes in ambiguous fault modes (all
read models, including shared aggregates), where an episode overclaims
if its final state is inconsistent although the agent claimed completion
and neither the agent nor the layer marked any operation as uncertain or
duplicated; turns and cost per episode.}\tabularnewline
\noalign{\global\arrayrulewidth=1.1pt}\arrayrulecolor{cvdeep}\hline\noalign{\global\arrayrulewidth=0.4pt}\arrayrulecolor{cvgrid}\rowcolor{cvhead}
\begin{minipage}[b]{\linewidth}\raggedright\bfseries
Model
\end{minipage} & \begin{minipage}[b]{\linewidth}\raggedright\bfseries
Condition
\end{minipage} & \begin{minipage}[b]{\linewidth}\raggedright\bfseries
Inconsistent \%, hard worlds
\end{minipage} & \begin{minipage}[b]{\linewidth}\raggedright\bfseries
Inconsistent \%, easy worlds
\end{minipage} & \begin{minipage}[b]{\linewidth}\raggedright\bfseries
Overclaim \% (ambiguous)
\end{minipage} & \begin{minipage}[b]{\linewidth}\raggedright\bfseries
Turns
\end{minipage} & \begin{minipage}[b]{\linewidth}\raggedright\bfseries
US\$ / episode
\end{minipage} \\
\hline
\endfirsthead
\noalign{\global\arrayrulewidth=1.1pt}\arrayrulecolor{cvdeep}\hline\noalign{\global\arrayrulewidth=0.4pt}\arrayrulecolor{cvgrid}\rowcolor{cvhead}
\begin{minipage}[b]{\linewidth}\raggedright\bfseries
Model
\end{minipage} & \begin{minipage}[b]{\linewidth}\raggedright\bfseries
Condition
\end{minipage} & \begin{minipage}[b]{\linewidth}\raggedright\bfseries
Inconsistent \%, hard worlds
\end{minipage} & \begin{minipage}[b]{\linewidth}\raggedright\bfseries
Inconsistent \%, easy worlds
\end{minipage} & \begin{minipage}[b]{\linewidth}\raggedright\bfseries
Overclaim \% (ambiguous)
\end{minipage} & \begin{minipage}[b]{\linewidth}\raggedright\bfseries
Turns
\end{minipage} & \begin{minipage}[b]{\linewidth}\raggedright\bfseries
US\$ / episode
\end{minipage} \\
\hline
\endhead
\endlastfoot
gpt-4.1-mini & C0 plain & 80.6 & 62.2 & 67.8 & 8.4 & 0.0018 \\ \hline
gpt-4.1-mini & C5 verify-before-retry & 70.1 & 37.8 & 55.9 & 7.1 &
0.0014 \\ \hline
gpt-4.1-mini & C6 guard & 56.7 & 24.4 & 32.2 & 8.0 & 0.0016 \\ \hline
gpt-4.1-mini & C7 reflection & 83.6 & 51.1 & 62.7 & 7.4 & 0.0015 \\ \hline
gpt-4.1-mini & S Settle & 6.0 & 0.0 & 3.4 & 7.2 & 0.0014 \\ \hline
gpt-5.6-luna & C0 plain & 54.5 & 12.2 & 13.6 & 9.6 & 0.0010 \\ \hline
gpt-5.6-luna & C5 verify-before-retry & 64.2 & 25.6 & 47.5 & 7.8 &
0.0007 \\ \hline
gpt-5.6-luna & C6 guard & 52.2 & 11.1 & 12.3 & 9.5 & 0.0009 \\ \hline
gpt-5.6-luna & C7 reflection & 56.7 & 12.2 & 15.7 & 9.4 & 0.0009 \\ \hline
gpt-5.6-luna & S Settle & 6.0 & 0.0 & 3.4 & 7.6 & 0.0008 \\ \hline
gpt-5.4-nano & C0 plain & 62.7 & 44.4 & 33.1 & 5.5 & 0.0008 \\ \hline
gpt-5.4-nano & C5 verify-before-retry & 62.7 & 24.4 & 42.4 & 4.8 &
0.0006 \\ \hline
gpt-5.4-nano & C6 guard & 47.8 & 24.4 & 11.0 & 5.4 & 0.0007 \\ \hline
gpt-5.4-nano & C7 reflection & 68.7 & 51.1 & 43.2 & 5.6 & 0.0007 \\ \hline
gpt-5.4-nano & S Settle & 6.0 & 0.0 & 3.4 & 4.8 & 0.0006 \\ \hline
\end{longtable}

\hypertarget{cost-h5}{%
\subsection*{8.10 Cost (H5)}\label{cost-h5}}
\addcontentsline{toc}{subsection}{8.10 Cost (H5)}

Settle's cost is latency, not model calls. In fault-free worlds it adds
1.01 s per scripted episode relative to the guard (95\% upper bound 1.29
s). In ambiguous worlds it adds 15.0 s on average, almost all of it the
reconciliation horizon, whose trade-off Section 8.6 traces. Settle
\emph{saves} model calls, because the agent rarely sees an error it has
to reason about. Its cost per episode relative to the guard is 0.84 for
\texttt{gpt-5.6-luna} (95\% upper bound 0.94), 0.87 for
\texttt{gpt-4.1-mini} (0.99) and 0.91 for \texttt{gpt-5.4-nano} (1.16),
against the pre-registered limit of 1.25; the cost criterion is
confirmed for all three, and the latency criterion is confirmed.

\hypertarget{an-internally-pre-registered-follow-up-study}{%
\section*{9. An internally pre-registered follow-up
study}\label{an-internally-pre-registered-follow-up-study}}
\addcontentsline{toc}{section}{9. An internally pre-registered follow-up
study}

The main study left six questions open that only new experiments could
settle. Was a simpler tag-based \emph{prevention} policy as good as
Settle? Did Settle's implementation gaps (Section 10.1) matter? Would a
duplicate that the network delivers late escape it? Do look-alikes from
another account defeat content checks as easily as same-account
look-alikes do? Does the language-model result depend on the harness's
plan mapping? And would a frontier model change it? I answered them in a
follow-up study with its own analysis plan, hypotheses F1--F7 and pass
criteria, frozen before the first follow-up episode
(\texttt{docs/07\_followup\_plan.md}; freeze digest prefix
\texttt{87ddadb5}). The main study's records and verdicts are unchanged.
The follow-up ran 4,272 episodes with no exclusions: 3,468 scripted
episodes, 536 episodes with the three original models and 268 with
\texttt{gpt-5.6-sol}. It cost US\$ 6.09 in model calls, under a hard cap
of US\$ 9.90. All seven hypotheses were confirmed (Table 12). One
analysis defect was corrected after the runs: the frozen script treated
a rate of exactly zero as missing, so it reported the two F4 criteria on
S2 as not confirmed although S2's measured rates (0 of 51 and 0 of 21)
met them. Every other verdict is the same under the frozen and the
corrected script; both are released (Appendix D). Several follow-up
hypotheses check properties that follow from the conditions'
definitions, such as that a policy without a final sweep misses
redelivery; F1(a), F5 and F7 are the informative ones.

\textbf{Prevention versus convergence (F1).} The new condition Sp is the
policy that the corrected Proposition 3 recommends for every write. It
keeps the durable journal and the tag, but after an ambiguous outcome it
waits out the in-flight bound, counts by tag and re-issues only if no
effect exists. It never compensates and has no final sweep. On the H1
population without redelivery and heavy-tailed delays, Sp and Settle
were \emph{identical}: neither left any of the 519 paired worlds
inconsistent. The two differ only where the theory says they must. On
transport redelivery Sp was inconsistent in 42 of 42 worlds and Settle
in 0, because preventing one's own re-issues does nothing about a
duplicate the transport creates. Under heavy-tailed delays Sp was worse
by +9.5 pp (54.8\% against 45.2\%), but this rests on 4 discordant
worlds (exact McNemar \(p\) = 0.125). In all of them the commit landed
between 21.0 and 25.6 s, just past the 20 s bound: Sp's last decision
comes \(\Delta\) after the attempt and Settle's at the end of the
episode, a few seconds later (Proposition 7). The difference is
therefore a matter of decision timing, which a slightly larger bound for
Sp would remove, not a robust advantage of convergence. In exchange, Sp
finished the episodes of the H1 population about 3 s sooner on average
(paired difference −3.0 s, 95\% CI −3.3 to −2.7 s). Over the whole H1
population Sp was inconsistent in 10.8\% of worlds against Settle's
3.2\%, and 42 of its 65 failures were redeliveries. The comparison
settles what Settle's extra machinery buys, and it is narrower than the
headline comparison suggests. When the failure is a lost reply, a
request that never arrived, a crash, a partial batch or a bounded delay,
a tag-based wait-then-count policy is exactly as good. Convergence
matters for duplicates that the transport creates.

\textbf{Settle as designed (F2, F3).} The variant S2 closes the first
two implementation gaps. It never confirms an ambiguous content-only
write by a content match, and its final sweep always waits out the
bound. On the 39 ambiguous content-only worlds, S2 left the same share
inconsistent as Settle (28.2\%). But it overclaimed in none of them,
against 6 for Settle, so every content-only failure is now reported.
S2's ambiguous content-only writes are reported as uncertain also when
they turned out fine, which is the cost of abstaining. On the H1
population S2 and Settle were again identical. The price of always
waiting is paid in every episode: S2 adds 19.8 s to a fault-free episode
(95\% CI 19.6 to 19.9 s), close to the assumed bound of 20 s. This
exceeds the main study's own latency criterion for fault-free episodes
(H5: at most 2 s more than the guard), which S2 would therefore fail.

\textbf{Delayed redelivery (F4).} In the new mode
\texttt{redelivery\_late}, the write is acknowledged normally, and the
transport delivers an identical copy 8 s later, after a Settle episode
without ambiguity has already finished. On the 51 worlds whose write is
reversible and attributable, every condition except S2 ended
inconsistent in every world, including Settle and the outcome oracle; S2
ended inconsistent in none. On the 21 worlds whose write is an e-mail,
every condition ended inconsistent, as Proposition 3 says it must: no
client can recall a delivered duplicate. But S2 reported every one of
them, while every other condition claimed success in all of them. On the
12 keyed worlds no condition failed, because the service drops the
duplicate key. The design choice is therefore real. A sweep that waits
only when something looked ambiguous saves about 20 s per clean episode,
and misses every duplicate that arrives after the agent has finished. S2
catches only duplicates that arrive within \(\Delta + \lambda\) of the
last attempt; a redelivery after its horizon is untested and would be
missed as well.

\textbf{Look-alikes from another account (F5).} In the new read variant
\texttt{lookalike\_other}, a second, ordinary Gitea user creates the
look-alike at the moment of the fault. Content checks without an author
filter fared as before: verify-before-retry and the waiting guard were
each inconsistent in 64.3\% of the 84 worlds. They failed even in some
lost-reply and late-commit worlds (6 of 21 each), in a way the main
study already showed for same-account look-alikes: after an issue's
reply was lost, they resolved the issue's number by its title, could
pick the other user's issue, and the task's comments were then missing
from the agent's own issue. With an author filter, the waiting guard
(Gwa) was inconsistent in 0 worlds, exactly like Settle.
Verify-before-retry with the filter (C5a) still failed every late-commit
world (100\%), because it does not wait. This result limits the scope of
Proposition 4 in practice, and it is shown here only for Gitea issues
and comments, four fault modes and the scripted driver. The tag is
\emph{necessary} when the duplicate-looking record can come from the
same account, for instance from another instance of the same agent or
another worker using a shared service account, or when the service
records no author, as for the mail, rows and stream events of my stack.
Where every other writer uses a different account and the service
exposes authorship, the client's own account is itself an attribution
source in the sense of Definition 1, and an author filter \emph{combined
with waiting out the bound} is as good as a tag.

\textbf{Plan-free intent identity (F6).} S−np binds each model write to
an intent by the tool and the normalised intent fields of the call
itself, without the harness's mapping onto plan steps. For all three
original models it produced exactly the same outcome as Settle in every
hard and every easy world. In hard worlds each model was again at 6.0\%,
with no discordant pair. The plan mapping therefore did not contribute
to the language-model results. Part of this identity is mechanical:
model responses are cached by the request, so where S−np and Settle
present a model with the same history it receives the same response, and
the two conditions can diverge only after their observations differ.
Binding by intent fields does carry the risk that two deliberately
identical writes are merged; no task in my suite contains such a pair.

\textbf{A frontier model (F7).} With \texttt{gpt-5.6-sol} as the agent
and no recovery layer, 52.2\% of hard worlds ended inconsistent (35 of
67), close to \texttt{gpt-5.6-luna}'s 54.5\%. With plan-free Settle
underneath, 6.0\% did (paired difference −46.3 pp, 95\% CI −53.5 pp to
−35.4 pp; exact McNemar \(p\) \textless{} 0.001). The frontier model
behaved like the careful smaller one. After its first ambiguous write it
checked state first in 77\% of episodes and re-issued at once in only
16\%. Its failures fell in the same cells: 74\% of look-alike worlds,
56\% of crash worlds and 27\% of heavy-tailed-delay worlds were
inconsistent, against 0\% after a lost reply with exact reads. A more
capable model does not escape the cells where the observation is
insufficient or the framework re-issues before the model acts. Its cost
per episode was US\$ 0.0197 with Settle and US\$ 0.0223 without. These
figures count paid calls only, and the two conditions share cached
responses for identical request histories, so they understate both.

\begin{longtable}[]{|
  >{\raggedright\arraybackslash}p{(\columnwidth - 2\tabcolsep - 2.5pt - 6\tabcolsep) * \real{0.4182}}|
  >{\raggedright\arraybackslash}p{(\columnwidth - 2\tabcolsep - 2.5pt - 6\tabcolsep) * \real{0.1182}}|
  >{\raggedright\arraybackslash}p{(\columnwidth - 2\tabcolsep - 2.5pt - 6\tabcolsep) * \real{0.3818}}|
  >{\raggedright\arraybackslash}p{(\columnwidth - 2\tabcolsep - 2.5pt - 6\tabcolsep) * \real{0.0818}}|}
\caption{\textbf{Table 12.} The pre-registered follow-up study:
inconsistent final states unless stated otherwise. Sp: tag
wait-then-count; S2: Settle as designed; S−np: Settle with plan-free
intent identity; C5a and Gwa: verify-before-retry and guard + wait with
an author filter.}\tabularnewline
\noalign{\global\arrayrulewidth=1.1pt}\arrayrulecolor{cvdeep}\hline\noalign{\global\arrayrulewidth=0.4pt}\arrayrulecolor{cvgrid}\rowcolor{cvhead}
\begin{minipage}[b]{\linewidth}\raggedright\bfseries
Question
\end{minipage} & \begin{minipage}[b]{\linewidth}\raggedright\bfseries
Worlds
\end{minipage} & \begin{minipage}[b]{\linewidth}\raggedright\bfseries
Result
\end{minipage} & \begin{minipage}[b]{\linewidth}\raggedright\bfseries
Verdict
\end{minipage} \\
\hline
\endfirsthead
\noalign{\global\arrayrulewidth=1.1pt}\arrayrulecolor{cvdeep}\hline\noalign{\global\arrayrulewidth=0.4pt}\arrayrulecolor{cvgrid}\rowcolor{cvhead}
\begin{minipage}[b]{\linewidth}\raggedright\bfseries
Question
\end{minipage} & \begin{minipage}[b]{\linewidth}\raggedright\bfseries
Worlds
\end{minipage} & \begin{minipage}[b]{\linewidth}\raggedright\bfseries
Result
\end{minipage} & \begin{minipage}[b]{\linewidth}\raggedright\bfseries
Verdict
\end{minipage} \\
\hline
\endhead
\endlastfoot
F1(a) Sp vs S, ordinary faults (H1 population) & 519 & identical: 0
discordant pairs & confirmed \\ \hline
F1(b) Sp vs S, transport redelivery & 42 & Sp 100\%, S 0\% &
confirmed \\ \hline
F1(c) Sp vs S, heavy-tailed delay & 42 & Sp 54.8\%, S 45.2\% &
confirmed \\ \hline
F2 S2, content-only writes: overclaim & 39 & S2 0, S 6 & confirmed \\ \hline
F3 S2 vs S (H1 population); fault-free cost & 603 & identical; +19.8 s
per fault-free episode & confirmed \\ \hline
F4(a) delayed redelivery, reversible writes & 51 & S2 0\%; S, Sp, Gw,
C0, O 100\% & confirmed \\ \hline
F4(b) delayed redelivery, e-mail: overclaim & 21 & S2 0\%; all others
100\% & confirmed \\ \hline
F5 look-alikes from another account (Gitea) & 84 & C5, Gw 64.3\%; C5a
25.0\%; Gwa, S, Sp, S2 0\% & confirmed \\ \hline
F6 plan-free intent identity, hard worlds & 134 / 67 / 67 & S−np = S for
all three models & confirmed \\ \hline
F7 \texttt{gpt-5.6-sol}, hard worlds & 67 & C0 52.2\%, S−np 6.0\% &
confirmed \\ \hline
\end{longtable}

\hypertarget{discussion}{%
\section*{10. Discussion}\label{discussion}}
\addcontentsline{toc}{section}{10. Discussion}

\textbf{Where idempotency keys and workflow engines fall short.} Keys
and durable execution answer the question ``has this \emph{request} been
processed?'' The question an agent faces after an ambiguous outcome is a
different one: ``how many \emph{effects} of this operation exist?'' The
two coincide only when the service enforces the key, and the census
shows that this is the exception outside payments. Without a key on the
service, a key the client generates is a string the service ignores,
however durable it is, and the effect it should protect can still be
duplicated by a late commit, a redelivery or a replayed call. A workflow
engine's durable history makes a crashed step resume, but resuming means
executing the activity again: the engine provides at-least-once
\emph{execution} by design {[}12{]}, {[}13{]}, and with key-less tools
that becomes at-least-once \emph{effects}. My measurements locate this
boundary. In-memory keys and the workflow engine left 60.9\% and 61.0\%
of ambiguous key-less episodes inconsistent, no different from plain
retries (61.0\%). On keyed writes the engine matched Settle exactly. And
on key-less writes its durable history did not stop a crashed step from
taking effect twice, in 100.0\% of crash-after-dispatch worlds. Even
where a key exists, its retention bounds the protection: a key that the
service forgets after an hour does not cover an agent that resumes the
next day.

\textbf{What API designers can do.} The most useful thing a key-less API
can offer an agent is not necessarily a key, but a persisted, queryable,
client-controlled field. A metadata map, an external id or a header that
the service stores and lets clients look up exactly is enough to build
attribution and convergence on the client (Proposition 2). The census
shows that many APIs already have such a field and that the gaps are
specific. Some APIs persist the field but cannot look it up, such as
Pub/Sub attributes or SQS message attributes. Others look it up only
through a search documented as lagging by an unstated amount, which
weakens assumption (c). Still others return only a count or a sum, which
removes attribution entirely (Proposition 5). The common design in which
the server returns the new id is no help in the case that matters, since
the id is lost together with the reply. Tool annotations such as
idempotency hints {[}21{]} describe intent but cannot provide either
property. A tool description that says how to \emph{find} an effect
again would provide it, and the tag can usually live in a field that end
users never see: 9 of the taggable key-less APIs in the census document
such a field.

\textbf{When a tag is needed.} Proposition 4 is an indistinguishability
argument, and the follow-up shows where it binds in practice. A
look-alike created by a \emph{different} account was harmless to a
content check that also filtered by author and waited out the bound;
such a check matched Settle. The tag is needed when the look-alike can
come from the same account (another instance of the same agent, or
another worker sharing its service account, which is common), and on
services that record no author, such as the mail, rows and stream events
of my stack. Since one cannot usually rule out the first case, the tag
remains the robust choice, but the claim is narrower than ``content
checks fail under look-alikes''; how often real services lack an author
field is something the census did not code.

\textbf{What agent frameworks can do.} Two framework behaviours caused
failures in my language-model runs. Replaying a pending tool call after
a restart turns every crash during a write into a blind re-issue
(Proposition 6). Leaving clean-up to the model is the second: under
look-alikes, a careful model removed what it took to be its own
duplicate and deleted someone else's record. The model was not careless;
by content, the two records were indistinguishable. Both behaviours are
fixed by two mechanisms that need no better model: a tag that survives
the restart, for instance in a journal written before dispatch, which
makes a restart resolve the outcome instead of repeating the call, and
attribution, which gates every destructive action.

\textbf{Prevention, convergence and what each costs.} The follow-up
study (Section 9) sharpens the choice between the two guarantees. A
tag-based wait-then-count policy, which is prevention without a key, was
exactly as good as Settle on lost replies, requests that never arrived,
held connections, bounded delays, crashes and partial batches, and about
3 s faster per episode of the H1 population. Convergence earned its
place for duplicates created by the transport, which no client-side
prevention can stop; its small advantage under long delays came from
four worlds just past the bound and reflects decision timing rather than
a robust difference. A deployment that knows its transport never
redelivers and its delays are bounded can therefore use the simpler
policy. One that cannot should converge. If late transport duplicates
are possible it should also sweep in every episode, at a cost of about
one in-flight bound of latency per episode, because a sweep that runs
only after visible trouble misses every duplicate that arrives after the
agent has finished; the main study's latency criterion would rule this
out, so the choice is a real trade-off.

\textbf{Convergence is not prevention, and latency is negotiable.}
Settle allows transient duplicates and removes them by the horizon. That
is acceptable for an issue, a row or an event, but not for effects that
others react to before compensation, such as a page that wakes someone
or a webhook that triggers a shipment. For irreversible effects the
theory and my measurements agree that convergence collapses to
prevention. Settle then waits out the in-flight bound before re-issuing,
and reports an unrecoverable duplicate instead of hiding it. Where even
a transient duplicate is harmful, the right tool is a key at the
service, and my results quantify what is lost without one. For
reversible effects the latency cost is not intrinsic. Proposition 7
shows that the residual risk falls only as the last decision moves
later, but nothing requires the agent to wait for that decision: the
journal outlives the episode, and a background reconciler can keep
counting effects by tag and compensate a late surplus after the agent
has reported. My implementation blocks until the horizon and I did not
evaluate this asynchronous variant. In my late-tail worlds it could have
removed at most the 22 delay failures on reversible writes; the other 14
concern irreversible writes, for which only a longer wait before the
re-issue helps.

\hypertarget{limitations-and-threats-to-validity}{%
\subsection*{10.1 Limitations and threats to
validity}\label{limitations-and-threats-to-validity}}
\addcontentsline{toc}{subsection}{10.1 Limitations and threats to
validity}

\textbf{External validity of the stack.} The services are real, but they
run on one machine under synthetic load, and the background actor is the
only other user. Real deployments have more concurrent writers, but I
have no data on how often they produce look-alikes as close as mine. In
the main study the background actor uses the same account and writes at
the moment of the fault, which is the worst case for checks by content;
the follow-up adds look-alikes from a second account. The census
establishes how real contracts are distributed over the classes of the
theory, but it codes documented behaviour, not behaviour under fault.
The two together argue that the classes I test are the ones real
services fall into; they do not show how often each fault occurs in
production. Production traces of ambiguous outcomes would close this
gap, and I know of no public source of them.

\textbf{Faults.} Faults are injected at the network and process level
and validated per service against the service's own records (Appendix
B). Real outages produce mixtures of faults and delays longer than my
twelve-second late commit. The heavy-tailed mode and the replay in
Section 8.6 show how the residual risk grows when the in-flight bound is
violated, as Propositions 2 and 7 predict, but my delay distribution is
a choice, not a measurement. The replay itself is exploratory: it relies
on Settle's decision rule, which it verified on 167 of the 168 recorded
late-commit worlds with exact reads, and not on new runs at each bound.

\textbf{Models.} I test four OpenAI models at their default settings,
all from one provider: three smaller models in the main study and the
frontier model \texttt{gpt-5.6-sol} in the follow-up, which ran one
repeat and two conditions within a fixed budget. The frontier model
behaved like the most careful smaller model, checking state first in
most episodes and failing in the same cells. Models from other
providers, other reasoning settings and agent frameworks that do not
replay pending calls remain untested. A model that compares authors and
creation times, or waits longer, might do better against less
adversarial look-alikes and shorter delays than mine.

\textbf{Baselines.} The guard {[}8{]}, verify-before-retry {[}23{]} and
rollback {[}14{]} were re-implemented from their published descriptions
and may differ from their authors' code in detail. I describe what each
re-implementation does (Section 7.3), use the strongest configuration of
each family I could build without attribution, and release the code.
Settle, the theory and the harness were all designed by me. The frozen
predictions, the matched-latency control, the per-mode breakdown and the
tag-aware ablations limit how much this can favour Settle, but they do
not replace an independent replication. The baselines of the main study
are denied tags by design, so part of the headline comparison confirms
Proposition 4 rather than testing an engineering choice. The follow-up
adds the two fairest competitors I could construct, a tag-based
wait-then-count policy and a waiting guard with an author filter, and
reports where each matches Settle.

\textbf{Pre-registration.} The freezes are internal: they consist of
source hashes recorded in the repository, not of a registration with a
third party, and the freeze, runs and writing took place within a few
days. The follow-up study has its own plan and freeze, written after
review of the main results and before any follow-up episode; its
analysis script needed one correction after the runs (Appendix D). Two
amendments were made to the main study after its freeze, one of them
after the runs. Both are documented with their evidence and re-freezes
(Appendix D). The second corrected a defect that had \emph{inflated} the
baselines' failure rates; the verdicts of H1--H4 and H6 held before and
after it, and H5's cost criterion for \texttt{gpt-5.4-nano} moved from
inconclusive to confirmed. The restriction of H6 to hard worlds was
decided on pilot data before the freeze, but it is a choice a reader may
weigh.

\textbf{Step mapping and grading.} For fault injection and grading,
language-model writes are mapped to plan steps by their key fields; a
rephrased write that escapes the mapping is treated as a new intent, and
its effects are still graded. Grading uses service-side records only.
Mail cannot be deleted, so a duplicated e-mail is always counted. For
Gitea deletions I rely on a compensation log written by the harness,
because deleted issues and comments are no longer visible through
Gitea's API.

\textbf{Tags: visibility and forgery.} In my stack the tag of an issue
or comment is a visible body line. Real deployments can often place it
in a hidden field instead; where only visible fields exist, the tag is a
cost in presentation that some teams may not accept. A visible tag can
also be copied into another record, for instance when someone quotes the
original. Settle would then count the copy as its own, skip a needed
write or compensate another user's record. Requiring the tag together
with the client's own author account would close most of this gap; my
implementation checks the tag only.

\textbf{What reversal leaves behind.} Deleting a surplus issue or row
restores the state my grader checks, but creating it may already have
sent notifications or triggered webhooks, and another user may have
commented on the duplicate before it was removed. The grader does not
count notifications, so reversible in my sense means reversible in the
service's state, not in its side effects. Interactions are a real
hazard: all of rollback's collateral damage arose because it deleted the
agent's own issue together with the other user's comment on it.

\textbf{Implementation gaps.} In the main study, three behaviours of the
implementation differed from the design that the theory supports: Settle
resolved an ambiguous content-only write by content, its final sweep
waited only after visible trouble, and in the language-model runs it
bound writes to intents through the harness's plan mapping. The
follow-up measured the design without these gaps (S2, S−np): it reports
every content-only failure, catches delayed redelivery of reversible
writes at a cost of about 20 s per episode, and gives identical
language-model outcomes without the plan mapping in every ambiguous
world. Two gaps remain untested. Compensations ran on a fault-free path,
so their own ambiguity was not exercised, and a visible tag could be
copied into another record.

\textbf{Design of the fault grid.} The pooled rates depend on the mix of
modes I chose. Two modes, redelivery and partial batches, are failed by
construction by every condition without a final sweep or row-level
checks, and Table 6 therefore also reports the comparison without them.
The crash pair penalises only eager re-issuing, because the driver's
checkpoint re-issues a step that crashed before dispatch in every
condition.

\hypertarget{conclusion}{%
\section*{11. Conclusion}\label{conclusion}}
\addcontentsline{toc}{section}{11. Conclusion}

An agent that loses the reply to a write does not know whether the
action happened, and no amount of reasoning recovers what the
observation does not contain. What recovers it is attribution: the
ability to count the effects of one's own operation. I defined
attributability as identifiability and derived which guarantee each API
contract and read model admits. Prevention needs a durable key unless
the in-flight bound is known and waited out, and a key protects only
within its retention. Convergence needs only attribution, reversibility
and an in-flight bound. A policy that must stop at a horizon has to bet
between two indistinguishable worlds. And abstention is the honest
answer everywhere else. A census of 52 APIs shows that keys are the
exception outside payments, while a taggable field is the rule. Settle
implements the map with a journal written before dispatch, a tag in a
field the API already persists, a reconciliation horizon and
compensation gated on evidence. On real services with real faults and
under a frozen protocol, Settle converged where the theory says
convergence is possible. Its remaining failures were violations of its
assumed in-flight bound, whose risk trades against latency, or lay where
the theory admits no guarantee. It reported every abstention on a shared
aggregate and every unrecoverable duplicate, and in its designed form
also every content-only failure and every late-arriving duplicate of an
e-mail. An internally pre-registered follow-up located the boundaries of
the method more precisely than the main study could. A tag-based
prevention policy is as good as convergence on ordinary faults, so
convergence is needed mainly for transport duplicates. An author filter
combined with waiting is as good as a tag when other writers use other
accounts, so the tag is needed for same-account look-alikes and services
without authorship. And a frontier model fails without the layer in the
same places as smaller ones, while with it all four models converge
alike. Keys and workflow engines remain the right tool when the service
honours a key. For the many services that do not, the question an agent
should ask is not ``did my request arrive?'' but ``how many of my
effects exist?''

\hypertarget{declarations}{%
\section*{Declarations}\label{declarations}}
\addcontentsline{toc}{section}{Declarations}

\textbf{Funding:} This research was self-funded by the author and
received no external funding.

\textbf{Competing interests:} The author declares no competing
interests.

\textbf{Author contributions:} Ankitha Vuyyala: conceptualization,
methodology, software, investigation, validation, formal analysis,
visualization, and writing (original draft, review and editing). The
author reviewed and approved the final manuscript and accepts
responsibility for its content.

\textbf{Data and code availability:} The Settle implementation, the
baselines, the fault-injection harness and service definitions, the API
census, the frozen study protocols and analysis plans, the recorded
episode outcomes and the analysis scripts are publicly available at
\url{https://github.com/ankitha-vuyyala/action-recovery}. The artifacts
corresponding to this manuscript are archived in release v1.0 at
\url{https://github.com/ankitha-vuyyala/action-recovery/releases/tag/v1.0}.
The release includes the raw-records archives. They hold the per-episode
folders of all 29,394 recorded episodes: the 20,500 episodes of the main
study (17,820 scripted and 2,680 language-model episodes) and the 4,272
episodes of the follow-up study, 24,772 in all, together with 4,622
episodes of development, pilot and gate-check runs, which are not part
of the confirmatory or follow-up analyses. Each folder contains the
configuration, dispatch and checkpoint logs, the grader reports and, for
language-model episodes, the full conversation with the model. The
archives also hold the response cache of the 17,171 paid model calls,
whose model, token usage, cost and time are logged in the repository as
well (\texttt{runs/cost\_ledger.jsonl}). The evaluation records in the
repository regenerate every reported number, table and figure, and the
response cache replays every model decision without an API key.
Third-party materials, including the service container images, remain
subject to their respective licenses and attribution requirements.

\textbf{Corresponding author:} Ankitha Vuyyala. Email:
\href{mailto:ankithavuyyala08@gmail.com}{\nolinkurl{ankithavuyyala08@gmail.com}}

\hypertarget{references}{%
\section*{References}\label{references}}
\addcontentsline{toc}{section}{References}

\begingroup\small\linespread{1.08}\selectfont\setlength{\parskip}{3pt}\everypar{\hangindent=2.4em\hangafter=1}

{[}1{]} E. A. Akkoyunlu, K. Ekanadham, and R. V. Huber, ``Some
constraints and tradeoffs in the design of network communications,'' in
\emph{Proc. 5th ACM Symp. Operating Systems Principles (SOSP)}, 1975,
pp.~67--74, doi: 10.1145/800213.806523.

{[}2{]} J. N. Gray, ``Notes on data base operating systems,'' in
\emph{Operating Systems: An Advanced Course, ser. Lecture Notes in
Computer Science, vol.~60, R. Bayer, R. M. Graham, and G. Seegmüller,
Eds. Berlin, Germany: Springer}, 1978, pp.~393--481, doi:
10.1007/3-540-08755-9\_9.

{[}3{]} P. Helland, ``Idempotence is not a medical condition: An
essential property for reliable systems,'' \emph{ACM Queue}, vol.~10,
no. 4, pp.~30--46, 2012, doi: 10.1145/2181796.2187821.

{[}4{]} C. Lee, S. J. Park, A. Kejriwal, S. Matsushita, and J.
Ousterhout, ``Implementing linearizability at large scale and low
latency,'' in \emph{Proc. 25th ACM Symp. Operating Systems Principles
(SOSP)}, 2015, pp.~71--86, doi: 10.1145/2815400.2815416.

{[}5{]} Stripe, ``Idempotent requests,'' \emph{Stripe API Reference}.
{[}Online{]}. Available: \url{https://docs.stripe.com/api/idempotent\_requests} (accessed Oct.~1,
2026).

{[}6{]} S. Sun, ``Did it happen? Counterfactual evaluation of LLM agent
recovery from ambiguous tool outcomes,'' Research Square preprint, 2026,
doi: 10.21203/rs.3.rs-10730245/v1.

{[}7{]} S. Gopnal Swamy, \emph{IdempotencyBench: Do LLM agents act
exactly once? Measuring idempotency violations under retries}, 2026.
{[}Online{]}. Available: \url{https://github.com/gssanjana4/idempotencybench}

{[}8{]} J. Li, ``Where does exactly-once live? Model, harness, and
tool-contract effects on duplicate side effects in LLM agents,''
arXiv:2609.29095, 2026.

{[}9{]} A. Trofimov and B. Novikov, ``When tool calls succeed but
workflows fail: Anomalies at the agent-tool boundary,''
arXiv:2609.15397, 2026.

{[}10{]} B. Mohammadi, N. Potamitis, L. Klein, A. Arora, and L.
Bindschaedler, ``Atomix: Timely, transactional tool use for reliable
agentic workflows,'' arXiv:2602.14849, 2026.

{[}11{]} Z. Chen et al., ``Cordon: Semantic transactions for tool-using
LLM agents,'' arXiv:2606.17573, 2026.

{[}12{]} Temporal Technologies, ``Activity definition,'' \emph{Temporal
Documentation}. {[}Online{]}. Available: \url{https://docs.temporal.io/activity-definition} (accessed Oct.~1, 2026).

{[}13{]} DBOS, ``Workflows,'' \emph{DBOS Docs}. {[}Online{]}. Available: \url{https://docs.dbos.dev/python/tutorials/workflow-tutorial} (accessed
Oct.~1, 2026).

{[}14{]} S. Perera, K. Hapuarachchi, F. Leymann, and R. Khalaf, ``Robust
agent compensation (RAC): Teaching AI agents to compensate,'' in
\emph{Proc. ACM Conf. AI and Agentic Systems (CAIS)}, 2026,
pp.~253--262, doi: 10.1145/3786335.3813141.

{[}15{]} E. Y. Chang and L. Geng, ``SagaLLM: Context management,
validation, and transaction guarantees for multi-agent LLM planning,''
\emph{Proc. VLDB Endow.}, vol.~18, no. 12, pp.~4874--4886, 2025, doi:
10.14778/3750601.3750611.

{[}16{]} G. Wang et al., ``Consistency and completeness: Rethinking
distributed stream processing in Apache Kafka,'' in \emph{Proc. ACM
SIGMOD Int. Conf. Management of Data}, 2021, pp.~2602--2613, doi:
10.1145/3448016.3457556.

{[}17{]} H. Zhang, A. Cardoza, P. B. Chen, S. Angel, and V. Liu,
``Fault-tolerant and transactional stateful serverless workflows,'' in
\emph{Proc. 14th USENIX Symp. Operating Systems Design and
Implementation (OSDI)}, 2020, pp.~1187--1204.

{[}18{]} TigerBeetle, ``Reliable transaction submission,''
\emph{TigerBeetle Docs}. {[}Online{]}. Available: \url{https://docs.tigerbeetle.com/coding/reliable-transaction-submission/}
(accessed Oct.~1, 2026).

{[}19{]} H. Garcia-Molina and K. Salem, ``Sagas,'' in \emph{Proc. ACM
SIGMOD Int. Conf. Management of Data}, 1987, pp.~249--259, doi:
10.1145/38713.38742.

{[}20{]} H. F. Korth, E. Levy, and A. Silberschatz, ``A formal approach
to recovery by compensating transactions,'' in \emph{Proc. 16th Int.
Conf. Very Large Data Bases (VLDB)}, 1990, pp.~95--106.

{[}21{]} Model Context Protocol, ``Specification 2025-06-18: Schema
reference (ToolAnnotations),'' \emph{Model Context Protocol website}.
{[}Online{]}. Available: \url{https://modelcontextprotocol.io/specification/2025-06-18/schema}
(accessed Oct.~1, 2026).

{[}22{]} Z. Wang, ``Callability is not operability: Controlled interface
interventions for LLM agents,'' arXiv:2608.23628, 2026.

{[}23{]} I. K. Mansoor, A. Phadke, and P. Rana, ``Verified tool calls
improve LLM agent reliability under non-atomic failures,''
arXiv:2608.02645, 2026.

{[}24{]} Z. Sun, X. Wang, and G. Li, ``Agentic transaction: Towards
ACID-compliant agent systems,'' arXiv:2608.13900, 2026.

{[}25{]} Y. Zhuang, K. Chen, Y. Duan, S. Zheng, J. Li, and X.-Y. Zhang,
``AgentRewind: Recoverable execution for long-horizon LLM agents,''
arXiv:2608.14380, 2026.

{[}26{]} K. Yang, P. Li, Z. Wu, K. Xu, H. Huang, and X. Huang, ``DART:
Semantic recoverability for structured tool agents,'' arXiv:2605.23311,
2026.

{[}27{]} Y. Zheng, Y. Yang, W. Zhang, and A. Quinn, ``ACRFence:
Preventing semantic rollback attacks in agent checkpoint-restore,''
arXiv:2603.20625, 2026.

{[}28{]} J. Xu, L. Fan, Z. Wang, X. Li, and H. Liu, ``Beyond single-use
tokens: Durable authorization state for replay-resistant LLM agent
actions,'' arXiv:2608.01710, 2026.

{[}29{]} A. Gupta, ``ReliabilityBench: Evaluating LLM agent reliability
under production-like stress conditions,'' arXiv:2601.06112, 2026.

{[}30{]} A. Sigdel and R. Baral, ``ToolMisuseBench: An offline
deterministic benchmark for tool misuse and recovery in agentic
systems,'' arXiv:2604.01508, 2026.

{[}31{]} D. Zhu et al., ``When tools fail: Benchmarking dynamic
replanning and anomaly recovery in LLM agents,'' arXiv:2606.05806, 2026.

{[}32{]} G. Tan et al., ``AgentChaos: Chaos engineering for agent
systems via programmatic fault injection,'' arXiv:2608.06790, 2026, to
appear at ASE 2026.

{[}33{]} Z. Zheng, J. Long, B. Li, and J. Yao, ``Stale does not mean
unsafe: Guard precision for tool-using LLM agents under infrastructure
state races,'' arXiv:2609.29522, 2026.

{[}34{]} Z. Zhai, M. Li, and X. Wang, ``Revisable by design: A theory of
streaming LLM agent execution,'' arXiv:2604.23283, 2026.

{[}35{]} L. Advani, ``From confident closing to silent failure:
Characterizing false success in LLM agents,'' arXiv:2606.09863, 2026.

{[}36{]} The Gitea Authors, \emph{Gitea: Git with a cup of tea}, 2026.
{[}Online{]}. Available: \url{https://github.com/go-gitea/gitea}

{[}37{]} Mailpit contributors, \emph{Mailpit: An email and SMTP testing
tool and API for developers}, 2026. {[}Online{]}. Available: \url{https://github.com/axllent/mailpit}

{[}38{]} Redis, ``Redis Streams,'' \emph{Redis Documentation}.
{[}Online{]}. Available: \url{https://redis.io/docs/latest/develop/data-types/streams/} (accessed
Oct.~1, 2026).

{[}39{]} SeaweedFS contributors, \emph{SeaweedFS: Distributed storage
system for object storage (S3), file systems, and Iceberg tables}, 2026.
{[}Online{]}. Available: \url{https://github.com/seaweedfs/seaweedfs}

{[}40{]} Shopify, \emph{Toxiproxy: A TCP proxy to simulate network and
system conditions for chaos and resiliency testing}, 2026. {[}Online{]}.
Available: \url{https://github.com/Shopify/toxiproxy}

{[}41{]} N. Shinn, F. Cassano, A. Gopinath, K. Narasimhan, and S. Yao,
``Reflexion: Language agents with verbal reinforcement learning,'' in
\emph{Advances in Neural Information Processing Systems (NeurIPS)},
vol.~36, 2023, pp.~8634--8652.

{[}42{]} OpenAI, ``Create a model response,'' \emph{OpenAI API
Reference}. {[}Online{]}. Available: \url{https://developers.openai.com/api/reference/resources/responses/methods/create}
(accessed Oct.~1, 2026).

\endgroup

\hypertarget{appendix-a.-proofs}{%
\section*{Appendix A. Proofs}\label{appendix-a.-proofs}}
\addcontentsline{toc}{section}{Appendix A. Proofs}

Throughout, a policy is a function from the client's history to its next
action, worlds are drawn from a prior, and \(N_\ell(t)\) is the number
of standing effects of operation \(\ell\) at time \(t\).

\textbf{Proposition 1.} \emph{(i) Prevention with a durable key.} Let
\(\kappa\) be a function of durable state, so that every attempt of
\(\ell\), before and after any crash, carries the same \(\kappa\). Let
\(e_1\) be the time of the first execution of any attempt. By the key
contract, the service records \(\kappa\) at \(e_1\), and every execution
with key \(\kappa\) in \([e_1, e_1 + r)\) creates no new effect. If
every execution of every attempt falls in that interval, exactly one
effect is created. Re-issuing until an acknowledgement or a refusal
arrives ends with that effect acknowledged, provided each attempt is
delivered with positive probability. With a known bound \(\Delta\),
every execution falls in \([t_1, t_{\text{last}} + \Delta]\), so
\(t_{\text{last}} - t_1 + \Delta < r\) suffices. Because the loop has no
time bound, prevention holds with the probability that the loop finishes
before this condition fails, and with certainty only if the number of
attempts is capped.

\emph{(i′) Retention.} An execution at \(e > e_1 + r\) is treated as
new. In a world where \(a_1\) executed at \(e_1\) and its
acknowledgement was lost, a late execution after \(e_1 + r\), whether of
a re-issue or of a delayed original, creates a second effect. After that
point the client must resolve the outcome by another attribution source
before re-issuing, which is the situation of Propositions 2 and 3.
State-based expiry is the same argument with the state change in place
of \(e_1 + r\).

\emph{(ii) Failure with an in-memory key.} Let the client crash after
dispatching attempt \(a_1\) with key \(\kappa_1\) and before recording
it. After the restart a fresh key \(\kappa_2 \neq \kappa_1\) is drawn.
In a world in which \(a_1\) is delayed and both \(a_1\) and the
re-issued \(a_2\) execute, the service sees two distinct keys and
creates two effects.

\emph{(iii) Without a key.} This is Li's Proposition 1 {[}8{]}, restated
in my notation. Take two worlds that agree up to time \(t\): in \(w_1\),
attempt \(a_1\) never executes; in \(w_2\), it executes at
\(t + \epsilon\). They are observation-equivalent up to \(t\). A policy
that re-issues at or before \(t\) duplicates in \(w_2\), and a policy
that never re-issues leaves \(\ell\) missing in \(w_1\). If \(\Delta\)
is unknown, such a \(t\) exists for every waiting time. ∎

\textbf{Proposition 2.} Consider one round, and let
\(t^\star = t_{\text{last}} + \Delta + \lambda\).

\begin{enumerate}
\def\labelenumi{\arabic{enumi}.}
\tightlist
\item
  By (d), every attempt dispatched up to \(t_{\text{last}}\) has
  executed by \(t_{\text{last}} + \Delta\) or never will.
\item
  By (b) and (c), a read at \(t^\star\) returns exactly the standing
  effects that carry \(\tau\) and executed before \(t^\star - \lambda\),
  which by step 1 are all that will ever exist from these attempts.
\item
  By (f) and (g), the effects that carry \(\tau\) are exactly the
  standing effects of \(\ell\), and no other actor removes them. The
  count read at \(t^\star\) therefore equals the eventual count of the
  attempts made so far.
\end{enumerate}

If \(N = 0\), no attempt executed and none is pending, so a re-issue
under the same \(\tau\) cannot race an earlier attempt, and a new round
begins. If \(N > 1\), targeted compensation of \(N - 1\) effects chosen
by id leaves exactly one, using (e). By (e) a compensation may be
retried through the faulty channel until a further read, after another
lag \(\lambda\), confirms \(N = 1\). If \(N = 1\), the policy stops. By
(a), \(\tau\) survives crashes, so any history resumes the same
procedure without creating an attempt it cannot attribute. Each round
re-issues at most once, and a round ends without a re-issue with
probability at least the probability that its attempt executes, so the
number of rounds is finite with probability one if that probability is
positive. No effect of another actor is touched, because only effects
that carry \(\tau\) are compensated and, by (g), they all belong to
\(\ell\).

For the bound, drop (d). A round that ends at horizon \(H\) can leave an
inconsistent state only if some attempt executes after \(H - \lambda\):
otherwise steps 2 and 3 hold at \(H\). For attempt \(i\), dispatched at
\(t_i\), this requires that it executes at all and that
\(\delta_i > H - \lambda - t_i\), which has probability
\(p_i(1 - F(H - \lambda - t_i))\) if delays are independent of whether
the attempt executes. A union bound over the attempts gives the stated
bound. ∎

\textbf{Proposition 3.} If effects cannot be neutralised, \(N_\ell = 1\)
at quiescence implies that exactly one execution happened. That is
prevention, so Proposition 1(iii) applies whenever \(\Delta\) is
unknown. If \(\Delta\) is known, waiting until
\(t_{\text{last}} + \Delta + \lambda\) before the count makes \(N = 0\)
equivalent to ``no attempt executed or pending'' (Proposition 2, steps
1--3), so a re-issue in that case cannot create a second execution. A
transport redelivery is an execution that no client action caused and
that no client action can prevent: in a world with redelivery and one
without, the client's history up to the first execution is the same, and
without a key the service accepts both executions. With a tag, the
client can detect the second effect and report it. ∎

\textbf{Proposition 4.} In \(w_1\) the client's attempt executed and
created effect \(e\). In \(w_2\) it did not, and another actor created
\(e'\) with the same content. A content read returns one matching effect
in both worlds, and every other observation is identical, so the two
worlds are observation-equivalent while \(N_\ell = 1\) in \(w_1\) and
\(N_\ell = 0\) in \(w_2\). A content-only policy takes the same action
in both worlds. If it skips, \(\ell\) is missing in \(w_2\). If it
re-issues, \(\ell\) is duplicated in \(w_1\). If it re-issues and then
removes one of the two matching effects, it either removes its new
effect, which leaves \(\ell\) missing in \(w_2\), or removes the earlier
match, which in \(w_2\) is \(e'\) (collateral). Either way it is wrong
in \(w_2\). With a tag, \(e\) carries \(\tau\) and \(e'\) does not, so a
read by tag distinguishes the two worlds. ∎

\textbf{Proposition 5.} \emph{If.} With exclusive writing, the aggregate
changes only through the client's own effects, so
\(N_\ell = (A(t) - A(t_0))/u\), where \(u\) is the unit of the
operation, \(A\) the aggregate and \(A(t_0)\) its value read before the
first attempt; \(\ell\) is attributable. \emph{Worst case.} Suppose
another actor may add one unit of the same size during the window. Then
a world in which the client's attempt executed and the other actor's did
not, and the reverse world, give the same aggregate. They are
observation-equivalent under an aggregate read, while \(N_\ell\)
differs, so any decision about re-issuing or compensating is wrong in
one of them. A fungible compensation removes ``a unit'' and cannot tell
whose. ∎

\textbf{Proposition 6.} Let the client crash after dispatching \(a_1\)
and before recording its outcome. After the restart, compare a world in
which the crash happened just before the dispatch, so that \(\ell\)
never ran, with one in which it happened just after, so that \(\ell\)
may have executed. Without a durable trace of the dispatch and without a
recoverable tag, the persisted state is the same in both. A policy that
issues \(\ell\) duplicates it in the second world if \(a_1\) executed. A
policy that does not issue leaves it missing in the first. A policy that
treats \(\ell\) as possibly dispatched and resolves it by reading needs
a tag it can count by, and an in-memory tag is lost. Conversely, if
\(\tau\) is recoverable, either from a write-ahead record or
deterministically from durable step identity, the policy ``treat the
step as possibly dispatched, wait \(\Delta + \lambda\), count by
\(\tau\), issue if \(N = 0\)'' is correct in both worlds by Proposition
2. ∎

\textbf{Proposition 7.} Fix a history \(h\) up to \(t_d\) that satisfies
the premise. Because \(\ell\) is attributable, the last read is exact
for effects that executed before \(t_d - \lambda\). Among the worlds
consistent with \(h\), the pending attempt either never executes
(\(w_0\)) or executes after \(t_d - \lambda\) (\(w_L\)); the policy
cannot tell them apart. After \(t_d\) it takes no action on \(\ell\), so
the number \(k \in \{0, 1\}\) of its other effects of \(\ell\) that
remain standing is fixed at \(t_d\). Here \(k = 0\) means that it did
not re-issue or that it compensated its re-issue, and \(k = 1\) means
that it re-issued or kept its re-issue. The final count is \(k\) in
\(w_0\) and \(k + 1\) in \(w_L\), so \(k = 0\) is wrong in \(w_0\) and
\(k = 1\) is wrong in \(w_L\). A policy that chooses \(k = 1\) with
probability \(q\) is inconsistent with probability
\[\begin{aligned} &q\Pr[w_L \mid h] + (1 - q)\Pr[w_0 \mid h] \\ &\quad\ge\; \min\{\Pr[w_0 \mid h], \Pr[w_L \mid h]\}. \end{aligned}\]
Settle chooses \(k = 1\), so its error at \(h\) is \(\Pr[w_L \mid h]\),
and it attains the minimum if and only if
\(\Pr[w_L \mid h] \le \Pr[w_0 \mid h]\). If Settle's own re-issue may
fail to execute, its error is larger by \(\Pr[w_0 \mid h]\) times the
probability of that failure. ∎

\hypertarget{appendix-b.-environment-validation}{%
\section*{Appendix B. Environment
validation}\label{appendix-b.-environment-validation}}
\addcontentsline{toc}{section}{Appendix B. Environment validation}

Before the test runs, and on the frozen code, the validation ran 510
trials in total, ten per service and fault mode. The hidden outcome was
read from the services' own records (validation run
\texttt{20261001-164459}). Table 13 reports how many trials produced the
intended outcome; 493 of 510 did.

\begin{longtable}[]{|
  >{\raggedright\arraybackslash}p{(\columnwidth - 2\tabcolsep - 4.5pt - 16\tabcolsep) * \real{0.2025}}|
  >{\raggedright\arraybackslash}p{(\columnwidth - 2\tabcolsep - 4.5pt - 16\tabcolsep) * \real{0.0633}}|
  >{\raggedright\arraybackslash}p{(\columnwidth - 2\tabcolsep - 4.5pt - 16\tabcolsep) * \real{0.1013}}|
  >{\raggedright\arraybackslash}p{(\columnwidth - 2\tabcolsep - 4.5pt - 16\tabcolsep) * \real{0.1519}}|
  >{\raggedright\arraybackslash}p{(\columnwidth - 2\tabcolsep - 4.5pt - 16\tabcolsep) * \real{0.0633}}|
  >{\raggedright\arraybackslash}p{(\columnwidth - 2\tabcolsep - 4.5pt - 16\tabcolsep) * \real{0.1392}}|
  >{\raggedright\arraybackslash}p{(\columnwidth - 2\tabcolsep - 4.5pt - 16\tabcolsep) * \real{0.0633}}|
  >{\raggedright\arraybackslash}p{(\columnwidth - 2\tabcolsep - 4.5pt - 16\tabcolsep) * \real{0.1266}}|
  >{\raggedright\arraybackslash}p{(\columnwidth - 2\tabcolsep - 4.5pt - 16\tabcolsep) * \real{0.0886}}|}
\caption{\textbf{Table 13.} Environment validation: trials whose
realised outcome (number of executions, read from the service) matched
the intended hidden outcome of the mode. The SeaweedFS conditional
create turns a redelivered request into a refusal (one execution, as
intended).}\tabularnewline
\noalign{\global\arrayrulewidth=1.1pt}\arrayrulecolor{cvdeep}\hline\noalign{\global\arrayrulewidth=0.4pt}\arrayrulecolor{cvgrid}\rowcolor{cvhead}
\begin{minipage}[b]{\linewidth}\raggedright\bfseries
Service
\end{minipage} & \begin{minipage}[b]{\linewidth}\raggedright\bfseries
none
\end{minipage} & \begin{minipage}[b]{\linewidth}\raggedright\bfseries
lost ack
\end{minipage} & \begin{minipage}[b]{\linewidth}\raggedright\bfseries
not exec err
\end{minipage} & \begin{minipage}[b]{\linewidth}\raggedright\bfseries
held
\end{minipage} & \begin{minipage}[b]{\linewidth}\raggedright\bfseries
not exec to
\end{minipage} & \begin{minipage}[b]{\linewidth}\raggedright\bfseries
late
\end{minipage} & \begin{minipage}[b]{\linewidth}\raggedright\bfseries
redelivery
\end{minipage} & \begin{minipage}[b]{\linewidth}\raggedright\bfseries
partial
\end{minipage} \\
\hline
\endfirsthead
\noalign{\global\arrayrulewidth=1.1pt}\arrayrulecolor{cvdeep}\hline\noalign{\global\arrayrulewidth=0.4pt}\arrayrulecolor{cvgrid}\rowcolor{cvhead}
\begin{minipage}[b]{\linewidth}\raggedright\bfseries
Service
\end{minipage} & \begin{minipage}[b]{\linewidth}\raggedright\bfseries
none
\end{minipage} & \begin{minipage}[b]{\linewidth}\raggedright\bfseries
lost ack
\end{minipage} & \begin{minipage}[b]{\linewidth}\raggedright\bfseries
not exec err
\end{minipage} & \begin{minipage}[b]{\linewidth}\raggedright\bfseries
held
\end{minipage} & \begin{minipage}[b]{\linewidth}\raggedright\bfseries
not exec to
\end{minipage} & \begin{minipage}[b]{\linewidth}\raggedright\bfseries
late
\end{minipage} & \begin{minipage}[b]{\linewidth}\raggedright\bfseries
redelivery
\end{minipage} & \begin{minipage}[b]{\linewidth}\raggedright\bfseries
partial
\end{minipage} \\
\hline
\endhead
\endlastfoot
Gitea issue & 10/10 & 10/10 & 10/10 & 10/10 & 10/10 & 10/10 & 10/10 & \\ \hline
Mailpit SMTP & 10/10 & 10/10 & 10/10 & 10/10 & 10/10 & 10/10 & 10/10
& \\ \hline
PostgreSQL row & 10/10 & 10/10 & 10/10 & 10/10 & 10/10 & 10/10 & 10/10
& \\ \hline
PostgreSQL batch & 10/10 & & & & & & & 10/10 \\ \hline
Redis stream & 10/10 & 10/10 & 10/10 & 10/10 & 10/10 & 10/10 & 10/10
& \\ \hline
Redis counter & 10/10 & 10/10 & 10/10 & 10/10 & 10/10 & 10/10 & 10/10
& \\ \hline
SeaweedFS S3 & 10/10 & 10/10 & 10/10 & 10/10 & 10/10 & 10/10 & 10/10
& \\ \hline
TigerBeetle & 10/10 & 10/10 & 0/10 & 10/10 & 3/10 & 10/10 & 10/10 & \\ \hline
\end{longtable}

Three further checks passed:

\begin{itemize}
\tightlist
\item
  \textbf{Late commits.} In every late-commit trial, the write landed
  after the caller had given up (minimum share 100\%).
\item
  \textbf{Tag fidelity.} The operation tag was read back byte-exact on
  every taggable service (minimum 100\%).
\item
  \textbf{Observation equivalence.} For every service, the paired modes
  (lost\_ack vs not\_exec\_err; held vs not\_exec\_to vs late) produced
  identical tool-level observations: the total variation distance
  between the distributions of (outcome kind, exception type) was 0.
  Elapsed times differ by tens of milliseconds at most, well below the
  variation between services.
\end{itemize}

For one service the raw client exception \emph{did} differ between the
paired modes: the SMTP client raised different exception types for a
reply cut after commit and a connection cut before delivery. The tool
layer therefore normalises all transport failures to one exception and
message, while the raw exception is kept in the logs for audit. A
deployed agent whose tool layer passes raw exceptions through may be
able to tell some of these worlds apart. My results are thus
conservative for such agents only in that respect.

\textbf{The TigerBeetle exception.} The TigerBeetle client library
retries internally and deduplicates by session. A request whose first
transmission was cut is therefore often delivered by the library's own
retry: the not\_exec modes realised the intended \emph{absent} outcome
in 0 and 3 of 10 trials. This is the native-key contract working as
designed. I keep these cells as measured, and they enter no equivalence
claim.

Crash modes are validated per episode by the child process's exit code
(75 at the crash, 0 after the restart). The late-tail mode uses the same
mechanism as the late mode, with a per-world delay.

\hypertarget{appendix-c.-agent-interface}{%
\section*{Appendix C. Agent
interface}\label{appendix-c.-agent-interface}}
\addcontentsline{toc}{section}{Appendix C. Agent interface}

\textbf{System instruction} (all language-model conditions):

\begin{quote}
You are an autonomous operations agent. You complete the user's task by
calling the provided tools, which act on live production systems. The
user is not available while you work, so do not ask them questions.
Follow the task instructions exactly. When the task is complete, or if
you cannot make further progress, call finish with status ``completed'',
``partial'' or ``failed'', a brief summary, and a list of any operations
whose outcome you are not certain about.
\end{quote}

\textbf{User message.} The task goal, followed by the sentence ``Each
action must take effect exactly once.'' An example for T7:

\begin{quote}
Triage ticket TCK-5476 for Woodgrove: open the issue `Ticket TCK-5476:
Woodgrove cannot export reports' (body: `Reported by Woodgrove. Priority
3.'), add exactly these two comments to it, in order: `Triage:
reproduced for Woodgrove; assigning to reports team.' and `ETA for
Woodgrove: fix in next release.', and record triage:TCK-5476 with amount
3. Each action must take effect exactly once.
\end{quote}

\textbf{Reflection prompt} (condition C7, sent after any turn with a
tool error):

\begin{quote}
One or more of your tool calls returned an error. Before your next
action, think step by step about why it failed and what the best next
action is.
\end{quote}

\textbf{Nudge} (sent if a turn contains no tool call; at most three
times):

\begin{quote}
Continue the task using the tools. When you are done, call finish.
\end{quote}

\textbf{Tools.} Each task exposes the following tools. Their
descriptions state documented service behaviour only; they say nothing
about faults or recovery methods.

\begin{itemize}
\tightlist
\item
  \emph{Writes:} create\_issue, add\_comment, send\_email,
  send\_bulk\_email, insert\_record, insert\_records, append\_event,
  adjust\_stock, increment\_counter, create\_transfer (with an optional
  transfer\_id, described as an idempotency key), put\_object.
\item
  \emph{Reads:} list\_issues, get\_issue, search\_issues (described as
  updated asynchronously), search\_sent\_email, find\_records,
  list\_events, get\_stock, get\_counter, find\_transfers,
  object\_exists.
\item
  \emph{The agent's own corrections:} delete\_issue, delete\_comment,
  delete\_record, delete\_event, reverse\_transfer.
\item
  \emph{Control:} wait (up to 60 s) and finish.
\end{itemize}

A task sees only the writes it needs, their reads and corrections, wait
and finish.

\textbf{Limits.} At most 30 model turns and 40 tool calls per episode,
with one model call per turn. Parallel tool calls were disabled, and
\texttt{max\_output\_tokens} was 6000, including reasoning.

\textbf{Endpoint and settings.}

\begin{itemize}
\tightlist
\item
  \emph{Endpoint and reasoning:} the OpenAI Responses API, stateless
  (\texttt{store=false}), with each model at its default reasoning
  setting.
\item
  \emph{Carried reasoning:} for reasoning models, the encrypted
  reasoning items are passed back with the transcript.
\item
  \emph{Model versions:} the resolved versions are recorded per call in
  the cost ledger.
\item
  \emph{Caching:} responses are cached by the SHA-256 of the request,
  which covers the model, instructions, input items, tools and
  parameters, together with the repeat index.
\end{itemize}

\textbf{Crash and resume.} The conversation is checkpointed to disk
after every model turn and tool result. After a crash, the restarted
agent re-executes every function call that has no recorded output. This
is the replay semantics of common agent frameworks, applied identically
in all conditions.

\textbf{Step mapping.} Faults target a plan step. A model's write is
mapped to a plan step when its key fields match the step after
normalising case, white space and trailing periods: the title for
issues, the recipient and subject for mail, the item for records, the
memo for transfers. For comments, the mapping uses the issue and the
text, and falls back to the next unused comment step of that issue. A
write that matches no step is an extra intent, and the grader still sees
its effects.

\hypertarget{appendix-d.-pre-registration-and-deviations}{%
\section*{Appendix D. Pre-registration and
deviations}\label{appendix-d.-pre-registration-and-deviations}}
\addcontentsline{toc}{section}{Appendix D. Pre-registration and
deviations}

\textbf{Freeze.} The protocol was frozen on 2026-10-01 22:37:09 UTC. The
freeze consisted of SHA-256 hashes of 43 files, with digest prefix
\texttt{51df7cbdbebd}. The frozen files were:

\begin{itemize}
\tightlist
\item
  all harness, recovery-layer, driver, task, grader and analysis code;
\item
  the container definitions;
\item
  the study design;
\item
  the theory;
\item
  the analysis plan;
\item
  this protocol.
\end{itemize}

The test presets check the freeze before every start or resume, and each
start is logged with the digest.

\textbf{Order of work.}

\begin{enumerate}
\def\labelenumi{\arabic{enumi}.}
\tightlist
\item
  Environment validation and a scripted development grid on the
  development split.
\item
  A paid pilot of 84 recorded episodes on the development split: four
  models (gpt-5.6-luna, gpt-5.4-mini, gpt-5.4-nano and gpt-4.1-mini),
  five conditions, easy worlds only. gpt-5.4-mini was not carried into
  the main study: in the pilot it cost about four times as much per
  episode as gpt-5.6-luna or gpt-5.4-nano (US\$ 0.0043 against 0.0011)
  and, like gpt-5.4-nano and gpt-4.1-mini, re-issued writes without
  checking, so it would have added cost without adding a distinct
  behaviour.
\item
  A check of the hard worlds with one model.
\item
  Full re-validation on the final code.
\item
  A harness check of the three test-only templates on a non-test
  instance; only errors and fault-free outcomes were inspected.
\item
  The freeze.
\item
  The test runs.
\item
  The analysis.
\end{enumerate}

\textbf{Hypotheses} (exact criteria; all evaluated by
\texttt{analysis/confirmatory.py}):

\begin{itemize}
\tightlist
\item
  \textbf{H1.} In class-A, targeted-reversible, ambiguous worlds,
  Settle's inconsistency is at most 5\%. Settle is also at least 15
  points below the best non-oracle baseline, with a 95\%
  cluster-bootstrap upper bound below 0 and a Holm-adjusted exact
  McNemar \(p < 0.05\). The best baseline is chosen on test data, which
  is conservative.
\item
  \textbf{H2.} Under look-alikes, each content-based method (C5, C6,
  S−t) is inconsistent in at least 10\% of ambiguous worlds, and
  Settle's rate is at most half of each, with a CI upper bound below 0
  and Holm \(p < 0.05\).
\item
  \textbf{H3.} On class-K writes, Settle is equivalent to the workflow
  engine with durable keys (C4): the 90\% CI of the difference lies
  within ±5 points.
\item
  \textbf{H4.}

  \begin{enumerate}
  \def\labelenumi{(\alph{enumi})}
  \tightlist
  \item
    With a shared aggregate, every non-oracle condition is inconsistent
    in at least 25\% of ambiguous worlds.
  \item
    Settle overclaims in at most 5\% of shared-aggregate or irreversible
    ambiguous worlds.
  \item
    In irreversible worlds, every condition shows duplicates or missing
    effects.
  \end{enumerate}
\item
  \textbf{H5.}

  \begin{itemize}
  \tightlist
  \item
    Settle adds at most 2 s in fault-free worlds (95\% upper bound) and
    at most \(\Delta\) + 10 s on average in ambiguous worlds, both
    relative to the guard.
  \item
    Per model, Settle's US\$ per episode relative to the guard is at
    most 1.25 (95\% upper bound).
  \end{itemize}
\item
  \textbf{H6.} For every model, in hard worlds (late\_tail,
  crash\_after, look-alike or lagged reads), Settle is below each
  language-model baseline, with a CI upper bound below 0 and Holm
  \(p < 0.05\).
\item
  \textbf{Ablations.} Each ablation is at least 10 points worse than
  Settle in the cell where its component is needed (lower 95\% bound
  above 0).
\end{itemize}

\textbf{Deviations from the original study design.} Each of the
following was decided on development data before the freeze, and is
logged with its evidence in the released change log.

\begin{enumerate}
\def\labelenumi{\arabic{enumi}.}
\tightlist
\item
  \emph{H1, statistical model.} The mixed-effects logistic model was
  replaced by cluster bootstrap plus exact McNemar, because of
  quasi-separation at near-zero rates.
\item
  \emph{H3, comparator and margin.} The comparator changed from
  in-memory keys (C3) to durable keys (C4), and the margin widened from
  ±3 to ±5 points. The reason is that in-memory keys fail in crash
  worlds exactly as Proposition 1 predicts, and that contrast is
  reported separately.
\item
  \emph{Hard worlds in the language-model grid.} The pilot showed that a
  reasoning model checks state before re-issuing. The language-model
  grid therefore gained the hard worlds, and H6 was restricted to them.
\item
  \emph{Settle guards the agent's own deletions.} This was added after a
  development episode in which the agent deleted another user's
  look-alike issue.
\item
  \emph{The OpenAI endpoint changed} from chat completions to the
  Responses API. Chat completions refused function tools for
  gpt-5.6-luna, the main model of the language-model study (the only one
  that reasons by default, and the one run with two repeats), at its
  default reasoning setting, before any charge.
\end{enumerate}

\textbf{Amendments after the freeze.} Both amendments are recorded in
the protocol's amendment log with their evidence. Each was followed by a
re-freeze, and every freeze is kept.

\begin{itemize}
\tightlist
\item
  \textbf{Amendment 1 (during the runs; re-freeze 2026-10-02 00:10:16
  UTC, digest prefix \texttt{a674ca5aedfc}).}

  \begin{itemize}
  \tightlist
  \item
    \emph{Cause.} The spend guard stopped the first model run on a
    Windows race in the lock file of the cost ledger.
  \item
    \emph{Fix.} The lock treats \texttt{PermissionError} like
    \texttt{FileExistsError}.
  \item
    \emph{Impact.} The lock guards only the spending ledger, so no
    outcome could change. The stopped episode was re-run.
  \end{itemize}
\item
  \textbf{Amendment 2 (after the runs; re-freeze 2026-10-03 05:06:59
  UTC, digest prefix \texttt{dd8b4adba131}).}

  \begin{itemize}
  \tightlist
  \item
    \emph{Cause.} The audit before writing found that models choosing
    their own transfer keys collided with the same key chosen in other
    episodes that replay the same task. TigerBeetle ids are global, and
    the tool layer reported the ledger's refusal as success.
  \item
    \emph{Fix.} Keys are now scoped to the episode, and refusals are
    surfaced as errors.
  \item
    \emph{Re-run.} Every T2 and T8 language-model episode, in all
    conditions, was re-run (540 episodes). The 12 episodes whose agent
    process had crashed without a report (9 language-model, 3 scripted)
    were also re-run; one of them was already among the T2/T8 episodes.
    A crashed agent process now counts as unfinished.
  \item
    \emph{Scope.} The scripted grid is unaffected by the key collision.
    Its keys are unique per run, and a status scan of its ledger calls
    found no \texttt{EXISTS\_WITH\_DIFFERENT\_*} refusal; the only
    repeat answers were the expected \texttt{EXISTS} for re-issued keys.
    Its only re-runs were the three crashed episodes.
  \item
    \emph{Reporting.} Section 8.9 gives the numbers before and after.
    The verdicts of H1--H4 and H6 did not change. H5's cost criterion
    for \texttt{gpt-5.4-nano} moved from inconclusive to confirmed.
  \end{itemize}
\end{itemize}

\textbf{Exploratory analyses added after the runs.} Two analyses were
added while writing, after all test runs and the confirmatory analysis
had finished: the API census (Section 5) and the horizon replay (Section
8.6). Neither enters any hypothesis, neither changes a confirmatory
number, and both are labelled as exploratory where they appear. The
census was coded from public documentation and does not use the test
data. The replay reads the recorded test episodes, but computes no rate
that a hypothesis depends on. Their code lives outside the frozen source
tree, so the freeze still verifies.

\textbf{Follow-up study.} The follow-up study (Section 9) has its own
analysis plan (\texttt{docs/07\_followup\_plan.md}) and freeze (digest
prefix \texttt{87ddadb5}), written after the main results had been
reviewed and before any follow-up episode ran; the main study's freeze
was archived first. Its spending was capped at US\$ 9.90 and came to
US\$ 6.09. A free smoke check on the development instance confirmed that
the new fault mode and the new read variant realised their intended
hidden outcomes before the grids ran. One correction was made after the
runs, to the analysis script only: in two F4 criteria the frozen code
read a missing rate as \texttt{rate\ or\ 1}, which also turns a measured
rate of exactly 0 into 1, so the criteria ``S2 ≤ 5\%'' and ``S2
overclaims ≤ 5\%'' failed although S2's measured rates were 0\%. With
the correction both criteria are confirmed; every other result of the
script is unchanged. The frozen script is released unchanged at
\texttt{analysis/followup.py}, so the freeze check still passes, and the
corrected copy, which also treats a missing rate as failing in every
criterion, is released at \texttt{paper/followup/followup\_corrected.py}
with both outputs. Under the frozen script, F4(a) and F4(b) read ``not
confirmed''.

\textbf{Exclusions.} Unfinished episodes (runner errors, grading errors,
model errors or budget stops) are archived and re-run. No scripted or
language-model record remained unfinished, so none is excluded.

\hypertarget{appendix-e.-the-api-census}{%
\section*{Appendix E. The API census}\label{appendix-e.-the-api-census}}
\addcontentsline{toc}{section}{Appendix E. The API census}

Table 14 lists the 52 coded APIs. ``Key'' records whether the examined
create operation has a server-side key, and its documented retention.
``partial'' marks the near-keys discussed in Section 5, which I do not
count as keys. ``Lookup by tag'' is the best documented way to find an
object by the value of the tag field: an exact filter or query
(\emph{exact}), a search documented as lagged or fuzzy (\emph{search}),
only scanning a list of recent objects (\emph{list scan}), or
\emph{none}. ``Undo'' records whether the effect can be undone through
the API. The released table also holds whether the tag field is hidden
from end users, the coder's notes, and a link to every documentation
page relied on. Documentation was read in October 2026, and vendors
change their APIs, so individual entries may have changed since.

\begingroup\def\tablefont{\fontsize{7.5}{9}\selectfont}\renewcommand{\arraystretch}{0.95}

\begin{longtable}[]{|
  >{\raggedright\arraybackslash}p{(\columnwidth - 2\tabcolsep - 3.7pt - 12\tabcolsep) * \real{0.2791}}|
  >{\raggedright\arraybackslash}p{(\columnwidth - 2\tabcolsep - 3.7pt - 12\tabcolsep) * \real{0.0291}}|
  >{\raggedright\arraybackslash}p{(\columnwidth - 2\tabcolsep - 3.7pt - 12\tabcolsep) * \real{0.2674}}|
  >{\raggedright\arraybackslash}p{(\columnwidth - 2\tabcolsep - 3.7pt - 12\tabcolsep) * \real{0.2267}}|
  >{\raggedright\arraybackslash}p{(\columnwidth - 2\tabcolsep - 3.7pt - 12\tabcolsep) * \real{0.0407}}|
  >{\raggedright\arraybackslash}p{(\columnwidth - 2\tabcolsep - 3.7pt - 12\tabcolsep) * \real{0.1163}}|
  >{\raggedright\arraybackslash}p{(\columnwidth - 2\tabcolsep - 3.7pt - 12\tabcolsep) * \real{0.0407}}|}
\caption{\textbf{Table 14.} The census: write contracts of 52 APIs,
coded from vendor documentation. Class K: server-side key; A: tag field
with exact lookup; A−: tag field with only lagged search or list
scanning; C: no usable tag field or no lookup.}\tabularnewline
\noalign{\global\arrayrulewidth=1.1pt}\arrayrulecolor{cvdeep}\hline\noalign{\global\arrayrulewidth=0.4pt}\arrayrulecolor{cvgrid}\rowcolor{cvhead}
\begin{minipage}[b]{\linewidth}\raggedright\bfseries
API (operation)
\end{minipage} & \begin{minipage}[b]{\linewidth}\raggedright\bfseries
Class
\end{minipage} & \begin{minipage}[b]{\linewidth}\raggedright\bfseries
Server-side key (retention)
\end{minipage} & \begin{minipage}[b]{\linewidth}\raggedright\bfseries
Tag field
\end{minipage} & \begin{minipage}[b]{\linewidth}\raggedright\bfseries
Hidden
\end{minipage} & \begin{minipage}[b]{\linewidth}\raggedright\bfseries
Lookup by tag
\end{minipage} & \begin{minipage}[b]{\linewidth}\raggedright\bfseries
Undo
\end{minipage} \\
\hline
\endfirsthead
\noalign{\global\arrayrulewidth=1.1pt}\arrayrulecolor{cvdeep}\hline\noalign{\global\arrayrulewidth=0.4pt}\arrayrulecolor{cvgrid}\rowcolor{cvhead}
\begin{minipage}[b]{\linewidth}\raggedright\bfseries
API (operation)
\end{minipage} & \begin{minipage}[b]{\linewidth}\raggedright\bfseries
Class
\end{minipage} & \begin{minipage}[b]{\linewidth}\raggedright\bfseries
Server-side key (retention)
\end{minipage} & \begin{minipage}[b]{\linewidth}\raggedright\bfseries
Tag field
\end{minipage} & \begin{minipage}[b]{\linewidth}\raggedright\bfseries
Hidden
\end{minipage} & \begin{minipage}[b]{\linewidth}\raggedright\bfseries
Lookup by tag
\end{minipage} & \begin{minipage}[b]{\linewidth}\raggedright\bfseries
Undo
\end{minipage} \\
\hline
\endhead
\endlastfoot
Stripe (create PaymentIntent) & K & native (24 h) & metadata & yes &
search (\textless1 min lag) & yes \\ \hline
PayPal Orders v2 (create order) & K & native (6 h) &
purchase\_units.custom\_id & yes & list scan & partial \\ \hline
Adyen Checkout (/payments) & K & native (7--14 d) & reference / metadata
& unclear & none & yes \\ \hline
Square (CreatePayment) & K & native (undocumented) & reference\_id &
unclear & list scan & yes \\ \hline
Braintree (transaction sale) & A & partial (30 s duplicate check) &
order\_id & unclear & exact & yes \\ \hline
Shopify Admin (orderCreate) & A & none & metafields & yes & exact &
yes \\ \hline
WooCommerce (create order) & A− & none & meta\_data & unclear & list
scan & yes \\ \hline
Razorpay (create order) & A & partial (unique receipt rejected) &
receipt & unclear & exact & no \\ \hline
Wise (create transfer) & K & native (undocumented) &
customerTransactionId & no & list scan & partial \\ \hline
Checkout.com (request payment) & K & native (72 h) & reference /
metadata & unclear & exact & yes \\ \hline
Mollie (create payment) & K & native (1 h) & metadata & unclear & list
scan & yes \\ \hline
GoCardless (create payment) & K & native (≥ 30 d) & metadata & unclear &
list scan & partial \\ \hline
GitHub (create issue) & A & none & label (exact) / body (search only) &
no & exact & partial \\ \hline
GitLab (create issue) & A & none & labels & no & exact & yes \\ \hline
Jira Cloud (create issue) & A− & none & issue property / labels & yes &
search (index lag) & yes \\ \hline
Linear (issueCreate) & A & partial (client-chosen id) & id (UUID) &
unclear & exact & yes \\ \hline
Asana (create task) & A & none & external.gid & yes & exact & yes \\ \hline
Trello (create card) & A− & none & desc & no & search & yes \\ \hline
Notion (create page in database) & A & none & text property & no & exact
& partial \\ \hline
Confluence Cloud (create page) & A & partial (unique title per space) &
title & no & exact & yes \\ \hline
Bitbucket Cloud (create PR comment) & A & none & content.raw & no &
exact & yes \\ \hline
Azure DevOps (create work item) & A & none & hidden custom field & yes &
exact & yes \\ \hline
monday.com (create\_item) & A & none & text column & no & exact & yes \\ \hline
ClickUp (create task) & A & none & text custom field & no & exact &
yes \\ \hline
Slack (chat.postMessage) & A− & none & metadata & unclear & list scan &
yes \\ \hline
Discord (create message) & A− & partial (nonce for minutes) & nonce &
unclear & list scan & yes \\ \hline
Microsoft Graph mail (sendMail) & A & none & extended property & unclear
& exact & no \\ \hline
Microsoft Teams (Graph) (channel message) & C & none & none & n/a & list
scan & partial \\ \hline
Gmail (messages.send) & A− & none & custom header & yes & search & no \\ \hline
SendGrid (mail send) & A & none & custom\_args & yes & exact (paid
add-on) & no \\ \hline
Mailgun (send message) & A & none & o:tag & unclear & exact (events) &
no \\ \hline
Twilio SMS (create message) & C & none & none & n/a & list scan & no \\ \hline
Postmark (send email) & A & none & Metadata & yes & exact & no \\ \hline
Amazon SES v2 (SendEmail) & A− & none & EmailTags & yes & list scan
(sampled) & no \\ \hline
Telegram Bot (sendMessage) & C & none & none & n/a & none & partial \\ \hline
WhatsApp Cloud (send message) & C & none & biz\_opaque\_callback\_data
(webhook only) & yes & none & no \\ \hline
Intercom (create message) & C & none & none & n/a & search (content) &
partial \\ \hline
Salesforce (create record) & A & partial (upsert by external id) &
External ID field & unclear & exact & yes \\ \hline
HubSpot (create contact) & A & partial (unique property) & unique
property & unclear & exact & yes \\ \hline
Zendesk (create ticket) & K & native (2 h) & external\_id & unclear &
exact & yes \\ \hline
Freshdesk (create ticket) & A− & none & custom field / tag & unclear &
search (minutes lag) & yes \\ \hline
Google Calendar (events.insert) & A & partial (client id; collisions not
guaranteed) & extendedProperties.private / id & unclear & exact & yes \\ \hline
Microsoft Graph calendar (create event) & K & native (undocumented) &
extended property & unclear & exact & partial \\ \hline
Airtable (create records) & A & partial (upsert on merge fields) & text
field & no & exact & yes \\ \hline
Google Sheets (values.append) & A− & none & cell value & no & list scan
& partial \\ \hline
Amazon S3 (PutObject (If-None-Match)) & K & conditional create (no
expiry) & object key & yes & exact & yes \\ \hline
Amazon DynamoDB (PutItem (attribute\_not\_exists)) & K & conditional
create (no expiry) & primary key & yes & exact & yes \\ \hline
Google Cloud Pub/Sub (publish) & C & none & attributes (no lookup) & yes
& none & no \\ \hline
Apache Kafka (producer send) & A− & partial (producer session only) &
record header & yes & list scan & no \\ \hline
Elasticsearch (index (op\_type=create)) & K & conditional create (no
expiry) & \_id & yes & exact & yes \\ \hline
PagerDuty Events v2 (trigger) & K & native (while alert open) &
dedup\_key & unclear & exact & partial \\ \hline
Amazon SQS (SendMessage) & C & partial (FIFO 5 min) & attributes (no
lookup) & yes & none & partial \\ \hline
\end{longtable}

\endgroup

\end{document}
